%% On Hecke algebras of the p-adic metaplectic cover SL2bar
%% Draft 14 (November 2026). Corrected and expanded version of
%%   D. Joyner, "A correspondence for the generalized Hecke algebra of the
%%   metaplectic cover SL(2,F)-bar, F p-adic", New York J. Math. 4 (1998).
%% Companion file: hecke-algebras-SL2bar-revision-log.tex
%% Compiles with pdflatex. Times via mathptmx (newtx is not installed here).
\documentclass[12pt,letterpaper]{amsart}
\usepackage[T1]{fontenc}
\usepackage{mathptmx}
\renewcommand{\ttdefault}{pcr}
\DeclareMathAlphabet{\mathcal}{OMS}{cmsy}{m}{n}
\usepackage[letterpaper,margin=1in]{geometry}
\usepackage{amsmath,amssymb,amsthm}
\usepackage{microtype}
\usepackage{enumitem}
\usepackage{booktabs}
\usepackage[hidelinks]{hyperref}

\numberwithin{equation}{section}
\setcounter{tocdepth}{2}

\newtheorem{theorem}{Theorem}[section]
\newtheorem{proposition}[theorem]{Proposition}
\newtheorem{lemma}[theorem]{Lemma}
\newtheorem{corollary}[theorem]{Corollary}
\theoremstyle{definition}
\newtheorem{definition}[theorem]{Definition}
\newtheorem{example}[theorem]{Example}
\newtheorem{question}[theorem]{Question}
\theoremstyle{remark}
\newtheorem{remark}[theorem]{Remark}
\newtheorem{vremark}{Remark}
\renewcommand{\thevremark}{2026.\arabic{vremark}}

\DeclareMathOperator{\vol}{vol}
\DeclareMathOperator{\ord}{ord}
\DeclareMathOperator{\pr}{pr}
\DeclareMathOperator{\Ad}{Ad}
\DeclareMathOperator{\tr}{tr}
\DeclareMathOperator{\Hom}{Hom}

% Table of contents: sections flush left, subsections indented under the
% section titles, numbers in fixed-width boxes so the titles line up.
\makeatletter
\def\l@section{\@tocline{1}{3pt plus 1pt}{0pt}{1.8em}{}}
\def\l@subsection{\@tocline{2}{0pt}{1.8em}{2.4em}{}}
\renewcommand{\tocsection}[3]{%
  \makebox[\@tempdima][l]{\@ifnotempty{#2}{#2.}}#3}
\let\tocsubsection\tocsection
\makeatother

\newcommand{\cO}{\mathcal{O}}
\newcommand{\cH}{\mathcal{H}}
\newcommand{\cP}{\mathcal{P}}
\newcommand{\fB}{\mathfrak{B}}
\newcommand{\fb}{\mathfrak{b}}
\newcommand{\C}{\mathbb{C}}
\newcommand{\Z}{\mathbb{Z}}
\newcommand{\Fq}{\mathbb{F}_q}
\newcommand{\Gbar}{\overline{G}}
\newcommand{\gen}{\mathrm{gen}}
\newcommand{\fbar}{\overline{f}}
\newcommand{\mat}[4]{\begin{pmatrix} #1 & #2\\ #3 & #4\end{pmatrix}}
\newcommand{\Hgen}{\cH(\Gbar/\!/B_k^*)_{\gen}}
\newcommand{\Hlin}{\cH(G/\!/B_k)}

\DeclareRobustCommand{\lowp}{\ensuremath{p}}
\begin{document}

\title[On Hecke algebras of \texorpdfstring{$\overline{SL_2}$}{SL2bar}]{On
Hecke algebras of the \texorpdfstring{\lowp}{p}-adic metaplectic cover
\texorpdfstring{$\overline{SL_2}$}{SL2bar}\\[6pt]
{\normalsize Draft 14. Corrected and expanded version}}

\author{David Joyner}
\email{wdjoyner@gmail.com}
\thanks{This paper is a significant rewrite of D.~Joyner, \emph{A
correspondence for the generalized Hecke algebra of the metaplectic
cover $\overline{SL(2,F)}$, $F$ $p$-adic}, New York J. Math.\
\textbf{4} (1998), 223--235 \cite{Jo98}, correcting, changing and
expanding on various topics. The revision log contains the list of
corrections from the original paper.}

\author{Claude}
\address{Anthropic, San Francisco, California, USA}

\date{Original draft: 1998. Corrected and expanded draft: 2026.}

\keywords{metaplectic groups, generalized Hecke algebras,
representation theory, $p$-adic groups, $SL(2)$}
\subjclass[2020]{Primary 22E50. Secondary 11F70, 20C08}

\begin{abstract}
This corrected and expanded version was prepared with substantial help
from AI systems, as described in the Acknowledgments.
It is known that the genuine representation theory of the metaplectic
cover of $SL(2)$ resembles that of $PGL(2)$, or more generally of a
quotient of $SL(2)$ by part of its center, rather than that of $SL(2)$
itself. This paper asks why a comparison with $SL(2)$ fails at the level
of Hecke algebras, and how far it is from being correct.
Let $G=SL(2,F)$, where $F$ is a $p$-adic field containing the $n$-th
roots of unity. Let $\Gbar$ be the $n$-fold metaplectic cover of $G$,
and let $B_k$ be the $k$-th filtration subgroup of an Iwahori subgroup.
For $k$ sufficiently large (for every $k\geq 1$ when $p\nmid n$), $B_k$
has a canonical lift $B_k^*$ to $\Gbar$, and every double coset
$B_kgB_k$ carries a genuine bi-$B_k^*$-invariant function
$\fbar_{\tilde g}$. Part~I follows the approach of the original paper.
The relations among the characteristic functions $f_g$ in $\Hlin$
transfer to the $\fbar_{\tilde g}$ with suitably chosen lifts, and when
$p\nmid n$ the part of the algebra supported on $SL(2,\cO)$ transfers
exactly. For the classical case $n=2$ we compute the remaining signs
explicitly. We then show, for every $n\geq 2$, that no algebra
homomorphism sends each $f_g$ to a multiple of $\fbar_{\tilde g}$ by a
constant of absolute value one. The proof rests on the fact that the
cover does not split. Part~II treats $n\geq 3$, where the obstruction
is local. Lifts of $t_0=h(\varpi)$ and of $h(u)$ commute only up to
$(u,\varpi)_n^{2}$, so no support-preserving homomorphism exists for
any choice of constants. The Hecke operators attached to the
translations $t_0^m$ with $n\nmid 2m$ have trace zero in every
admissible genuine representation. The same translations carry no
genuine functions at Iwahori level. There we give a presentation of the
genuine Iwahori--Hecke algebra and a direct proof of Savin's theorem
for $SL(2)$. The genuine algebra is isomorphic to the Iwahori--Hecke
algebra of $SL(2)$ for odd $n$ and of $PGL(2)$ for even $n$, by an
explicit map that rescales the translations and multiplies each basis
function by the square root of a ratio of volumes. We compute the
orbital integrals of all basis functions at split elements at Iwahori
level, and show that Savin's isomorphism transfers them, with an
explicit transfer factor. We also compute the elliptic orbital integrals
on the linear side, reduce the genuine ones to three values, and show
that the values found numerically imply an elliptic transfer, of
endoscopic shape for even $n$.
\end{abstract}

\maketitle

\tableofcontents

%======================================================================
\section{Introduction}
\label{section:intro}

This section states the problem, explains why it is natural, and
summarizes the results.

Let $F$ be a $p$-adic field of characteristic $0$ with uniformizer
$\varpi$, ring of integers $\cO$, and residue field $\Fq$ of
cardinality $q$. Assume that $F$ contains the group $\mu_n$ of $n$-th
roots of unity and that $|\mu_n|=n$. Let $G=SL(2,F)$ and let $\Gbar$
be its $n$-fold metaplectic cover, the central extension of $G$ by
$\mu_n$ constructed by Kubota \cite{Ku67}. Representations of $\Gbar$
on which $\mu_n$ acts by a faithful character are called
\emph{genuine}. They are the local objects behind modular forms of
fractional weight.

\subsection{Motivation}
For $n=2$ it has long been known that the genuine representation theory
of $\Gbar$ resembles that of $PGL(2)$ rather than that of $SL(2)$.
Weil's representation and Waldspurger's Shimura correspondence
\cite{Wa80} relate genuine representations of the double cover of
$SL(2)$ to representations of $PGL(2)$. Flicker and Kazhdan
\cite{Fl80,FK86} compared covers of $GL(r)$ with $GL(r)$ by the trace
formula. At the level of Iwahori--Hecke algebras, Savin \cite{Sa88}
proved that the genuine Iwahori--Hecke algebra of the $n$-fold cover is
isomorphic to the Iwahori--Hecke algebra of $G/Z_n$, where $Z_n$ is the
$n$-torsion of the center (Theorem~\ref{thm:savin}). This is $PGL(2)$
when $n$ is even. For $SL(2)$ we give a direct proof, with an explicit
isomorphism (Theorem~\ref{thm:savinsl2}), and
Section~\ref{subsec:savinbg} explains the background. So the natural linear partner of the cover is a
quotient of $SL(2)$ by part of its center, and its Hecke algebra sees
the translations of $SL(2)$ only after a rescaling.

The original version of this paper \cite{Jo98} asserted that at
positive level the comparison works with $SL(2)$ itself, that is, that
$f_g\mapsto\fbar_g$ is an algebra isomorphism from $\cH(G/\!/B_k)$ onto
the genuine Hecke algebra. That assertion is false. The aim of this
version is to repair it honestly. We explain why the comparison with
$SL(2)$ fails at the level of Hecke algebras, and we measure how far
it is from being correct. The answer is summarized in
Table~\ref{tab:howfar} below.

\subsection{The problem}
For an open compact subgroup $K_0$ of $G$, the representations of $G$
with non-zero $K_0$-fixed vectors are controlled by the
\emph{generalized Hecke algebra} $\cH(G/\!/K_0)$ of bi-$K_0$-invariant
compactly supported functions. Howe and Moy \cite{HM85} described these
algebras for $GL(n)$ and the filtration subgroups of an Iwahori
subgroup by generators and relations. If $K_0$ lifts to a subgroup
$K_0^*$ of $\Gbar$, there is an analogous algebra
$\cH(\Gbar/\!/K_0^*)_{\gen}$ of genuine functions, which controls the
genuine representations with $K_0^*$-fixed vectors. A natural question
is the following.

\begin{quote}
\emph{Are $\cH(G/\!/K_0)$ and $\cH(\Gbar/\!/K_0^*)_{\gen}$ isomorphic
by a map that sends the characteristic function of each double coset
$K_0gK_0$ to a function supported over the same double coset?}
\end{quote}

\noindent
Such an isomorphism would match the representations of $G$ and of
$\Gbar$ with fixed vectors, together with their traces on these
algebras. We study the question for $K_0=B_k$, the $k$-th filtration
subgroup of the standard Iwahori subgroup $B$ of $G$.

\subsection{Main results}
The level $k$ is called \emph{deep} if $k\geq 1$ and $p\nmid n$, or if
$k$ exceeds an explicit bound $j_1$ in general
(Definition~\ref{def:deep}). Throughout, $t_0=h(\varpi)$ is the
diagonal matrix with entries $\varpi,\varpi^{-1}$, and $h(u)$ is the
diagonal matrix with entries $u,u^{-1}$.

Part~I (Sections~\ref{section:sl2}--\ref{section:n2}) follows the
approach of the original paper. For deep $k$ we prove the following.
\begin{enumerate}[leftmargin=2em]
\item $B_k$ has a canonical lift $B_k^*$ to $\Gbar$, compatible with
conjugation by all of $\Gbar$ (Lemma~\ref{lemma:canonical}).
\item Every double coset $B_kgB_k$ supports a genuine
bi-$B_k^*$-invariant function $\fbar_{\tilde g}$, normalized by a
choice of lift $\tilde g$ of $g$. Hence $f_g\mapsto\fbar_{\tilde g}$ is
a linear isomorphism $\Phi$ from $\Hlin$ onto $\Hgen$
(Proposition~\ref{prop:admissible}). It is support-preserving in the
sense of Definition~\ref{def:supportpreserving}: it maps a function
supported on a double coset to a function supported over the same
double coset.
\item If the volumes of $B_kg_1B_k$ and $B_kg_2B_k$ multiply to the
volume of $B_kg_1g_2B_k$, then
$\fbar_{\tilde g_1}*\fbar_{\tilde g_2}=\fbar_{\tilde g_1\tilde g_2}$
(Lemma~\ref{lemma:genuine}).
\item $\Hlin$ is generated by the $f_g$ with $g$ in
$\{s_1,s_2\}\cup B$, which satisfy the relations of
Proposition~\ref{prop:main}. Each of these relations holds in $\Hgen$
for the $\fbar_{\tilde g}$, with lifts chosen to suit that relation
(Proposition~\ref{prop:main2}).
\item If $p\nmid n$, the part of $\Hlin$ supported on $K=SL(2,\cO)$
transfers to $\Hgen$ exactly (Proposition~\ref{prop:Ksub}).
\item For $n=2$ and $p$ odd, lifts of $t_0$ commute with lifts of the
diagonal units, every lift of $-1$ has order~$2$, and the relation
between $t_0$ and $s_1$ acquires the explicit sign $(-1,\varpi)_2$
(Lemma~\ref{lemma:signs}).
\item For every $n\geq 2$, no algebra homomorphism $\Hlin\to\Hgen$
sends each $f_g$ to $c_g\fbar_{\tilde g}$ with $|c_g|=1$
(Theorem~\ref{thm:unimodular}). Suppose $\Psi$ were such a
homomorphism. What fails is the multiplication table. In $\Hlin$ the coefficient of the double coset of $xz$ in
$f_x*f_z$ is a positive integer $N$, the number of cosets $x'B_k$ in
$B_kxB_k$ with $x'^{-1}xz\in B_kzB_k$. Evaluate both sides of
$\Psi(f_x*f_z)=\Psi(f_x)*\Psi(f_z)$ at a lift of $xz$. The left side is
$N$ times a number of absolute value one, and the right side is a sum of
$N$ numbers of absolute value one, one for each coset. They agree only
if all $N$ terms are equal. Requiring this for every pair $x,z$ makes
the function $\psi(\tilde y)=\Psi(f_y)(\tilde y)$, where $f_y$ is the
basis function of the double coset of $y$, a homomorphism from $\Gbar$
to $\C^\times$ that is the identity on $\mu_n$. Its kernel would then
map isomorphically onto $G$, which would give a splitting of the
cover. No splitting exists, because a
splitting $s$ would satisfy $s(h(a))\,s(h(b))=(b,a)_n\,s(h(ab))$
(Lemma~\ref{lemma:nonsplit}). Products with $N=1$, those whose volumes
multiply, cause no trouble (Lemma~\ref{lemma:genuine}). The failure is
in products with $N\geq 2$ (Remark~\ref{rem:n2where}). These constants $N$ do not shrink as $k$
grows. For $z=x^{-1}$ we have $N=q^{\ell(w)}$ at every level, where
$x\in BwB$ (Remark~\ref{rem:Nk}). On a monoid of products with $N=1$,
$\Phi$ is multiplicative exactly when the lifts can be chosen
multiplicatively (Proposition~\ref{prop:monoid}). Choosing other
lifts does not help, since it only changes the constants by roots of
unity. Constants of other absolute values are ruled out for $n\geq 3$
(Corollary~\ref{corollary:main}), and for $n=2$ the question is open
(Question~\ref{q:n2}).
\end{enumerate}

\noindent
Part~II (Sections~\ref{section:torus}--\ref{section:savin}) treats
$n\geq 3$. Its last section, on Iwahori level, treats every $n\geq 2$.
\begin{enumerate}[leftmargin=2em,resume]
\item For every unit $u$ and every choice of lifts,
$\fbar_{\tilde t_0}*\fbar_{\tilde h(u)}=(u,\varpi)_n^{-2}\,
\fbar_{\tilde h(u)}*\fbar_{\tilde t_0}$, while $f_{t_0}$ and
$f_{h(u)}$ commute (Theorem~\ref{thrm:main}). So no support-preserving
homomorphism exists, whatever the constants
(Corollary~\ref{corollary:main}), and the expected trace identities
fail (Example~\ref{ex:trace}).
\item At Iwahori level ($k=0$, $p\nmid n$) the double cosets $Bt_0^mB$
and $Bt_0^ms_1B$ support non-zero genuine functions if and only if $n$
divides $2m$ (Proposition~\ref{prop:level0} and
Lemma~\ref{lemma:level0s1}). This holds for every $n\geq 2$.
\item The genuine Iwahori--Hecke algebra is generated by two elements
$X$ and $Y$ with the defining relations $X^2=(q-1)X+q$ and
$Y^2=c_1Y+c_0$, where $c_0$ and $c_1$ are computed explicitly
(Theorem~\ref{thm:presentation}). This gives a direct proof of Savin's
theorem \cite{Sa88} for $SL(2)$. The isomorphism multiplies each basis
function by $q^{e}$, where $q^{2e}$ is the ratio of the volumes of the
two double cosets involved (Theorem~\ref{thm:savinsl2}). For even $n$
no isomorphism with the Iwahori--Hecke algebra of $SL(2)$ maps $Y$ to a
multiple of $f_{s_1}$ or $f_{s_2}$ (Corollary~\ref{cor:wronggroup}).
\item The orbital integrals at split elements of all basis functions of
the Iwahori--Hecke algebras of $SL(2)$, $PGL(2)$ and $\Gbar$ are
computed explicitly (Propositions~\ref{prop:orbsl2} and
\ref{prop:orbpgl}, Corollary~\ref{cor:orbgen}). Savin's isomorphism
transfers them, with an explicit transfer factor of absolute value one
(Theorem~\ref{thm:orbtransfer}). At deep level, every orbital integral
of $\fbar_{\tilde t_0^m}$ vanishes when $n\nmid 2m$
(Proposition~\ref{prop:orbdeep}).
\item At Iwahori level the elliptic orbital integrals of all basis
functions of the Iwahori--Hecke algebras of $SL(2)$ and $PGL(2)$ are
counts of edges of the Bruhat--Tits tree, computed explicitly
(Propositions~\ref{prop:ellsl2} and \ref{prop:ellpgl}). The genuine ones
are determined by three values (Proposition~\ref{prop:ellgen}). The
values computed numerically imply that Savin's isomorphism transfers
elliptic orbital integrals too, class by class for odd $n$, and in
endoscopic form for even $n$ (Theorem~\ref{thm:ellcond}).
\end{enumerate}

\noindent
\textbf{The cases $n=2$ and $n\geq 3$.} For $n=2$ the failure is
global. Each relation that we compute can be matched on its own. The
contradiction appears only when all structure constants are matched at
once, which would produce a splitting of the cover.
The argument needs constants of absolute value one, as in the map
$f_g\mapsto\fbar_{\tilde g}$. Whether a homomorphism with other
constants exists for $n=2$ is open (Question~\ref{q:n2}). For
$n\geq 3$ the failure is local. It is already visible in a single
commutation relation, and it holds for arbitrary constants.

\subsection{How far off is the comparison with \texorpdfstring{$SL(2)$}{SL(2)}?}
The comparison can be tested at several levels, from the coarsest to
the finest. Table~\ref{tab:howfar} summarizes what we prove.

\begin{table}[ht]
\caption{How the genuine Hecke algebra compares with $\Hlin$, at deep
level $k$ unless stated otherwise.}
\label{tab:howfar}
\begin{tabular}{@{}p{0.42\textwidth}p{0.25\textwidth}p{0.25\textwidth}@{}}
\toprule
Test & $n=2$, $p$ odd & $n\geq 3$, $p\nmid n$ \\
\midrule
Supports and volumes of double cosets & agree & agree \\
Supports at Iwahori level ($k=0$) & all $t_0^m$ & exactly $t_0^m$ with $n\mid 2m$ \\
Each relation of Prop.~\ref{prop:main}, on its own & holds up to a sign &
holds up to a root of unity \\
Subalgebra supported on $SL(2,\cO)$ & identical & identical \\
Commutation of $f_{t_0}$ and $f_{h(u)}$ & holds & fails \\
Homomorphism $f_g\mapsto c_g\fbar_{\tilde g}$ with $|c_g|=1$ & none &
none \\
Same, with arbitrary $c_g\neq 0$ & open & none \\
Trace identities & open & fail \\
Spherical part over $t_0^m$ & no vanishing forced & zero when $n\nmid 2m$ \\
Iwahori--Hecke algebra ($k=0$) & of $PGL(2)$ type & of $SL(2)$ type if
$n$ is odd, of $PGL(2)$ type if $n$ is even \\
Split orbital integrals ($k=0$) & transferred to $PGL(2)$ & transferred
to $G/Z_n$ \\
Elliptic orbital integrals ($k=0$), if Remark~\ref{vrem:ellcomputed}
holds & stable part to $PGL(2)$ & to $SL(2)$ ($n$ odd), as for $n=2$
($n$ even) \\
\bottomrule
\end{tabular}
\end{table}

At deep level the two algebras have the same supports and the same
volumes, and each relation of $\Hlin$ holds in the genuine algebra up to
a root of unity. The discrepancy is a twist by Kubota's cocycle, seen
through the double cosets of $B_k$ (Remark~\ref{rem:twist}). A choice of
lifts removes the twist from any one relation, but not from all of them
at once.

For $n=2$ the twist is invisible on the torus and on $SL(2,\cO)$. It is
detected only by products of double cosets whose volumes do not
multiply, and its non-triviality is the same phenomenon as the
non-splitting of the cover (Theorem~\ref{thm:unimodular} and
Remark~\ref{rem:n2where}).

For $n\geq 3$ the twist is already visible on the torus. Lifts of
$t_0^m$ and $h(u)$ commute only up to $(u,\varpi)_n^{2m}$, and in every
admissible genuine representation the Hecke operators attached to
$t_0^m$ with $n\nmid 2m$ have trace zero
(Proposition~\ref{prop:traceless}). These are the translations that
carry no genuine functions at Iwahori level
(Proposition~\ref{prop:level0}), where Savin's theorem replaces $SL(2)$
by $G/Z_n$. So the translations $t_0^m$ with $n\nmid 2m$ are where
$SL(2)$ is ``off''. At Iwahori level they disappear from the genuine
algebra, and at deep level they remain but are invisible to traces
(Remark~\ref{rem:samelattice}).

Savin's isomorphism at Iwahori level is of a different kind from the
maps considered above. It moves functions to other double cosets and
rescales them by powers of $\sqrt q$ (Theorem~\ref{thm:savinsl2}).
These powers are the square roots of the ratios of the volumes of the
double cosets, so the isomorphism carries an orthonormal basis to an
orthonormal basis. Savin's isomorphism transfers the orbital integrals
at split elements, and these powers are exactly what the Weyl
discriminant factors require (Theorem~\ref{thm:orbtransfer}). For odd $n$ the target is the
Iwahori--Hecke algebra of $SL(2)$ itself, but the translation $t_0^{nk}$
goes to $t_0^k$. For even $n$ the target is that of $PGL(2)$, and the
affine generator has the eigenvalues $\pm\sqrt{c_0}$ instead of $q$
and $-1$ (Section~\ref{subsec:savinbg}). Whether the
same holds at elliptic elements, and whether there is an analog at
deep level, is Question~\ref{q:transfer}.

\subsection{Status of results}
\label{subsec:status}
A statement with a proof in this paper is proved here. A statement
quoted from the literature carries a citation in its heading.
Statements that we could not verify against their sources are flagged
in remarks numbered ``Remark~2026.$m$''. Questions are open.

\subsection{Running example}
Our running example is $F=\mathbb{Q}_7$, with $n=2$ in Part~I and $n=3$
in Part~II. Since $6$ divides $7-1$, the field $\mathbb{Q}_7$ contains
$\mu_6$, so both covers are defined. Since $p=7$ divides neither $n$,
every level $k\geq 1$ is deep.

\subsection{Literature review}
\label{section:lit}

This subsection summarizes the work most closely related to ours. The
first part covers work up to 1998, the year of the original version
\cite{Jo98}. The second part covers later work.

\subsubsection{Work up to 1998}
Kubota \cite{Ku67} constructed the topological $n$-fold covering of
$SL(2)$ over a local field by an explicit cocycle built from the
$n$-th power Hilbert symbol. We use his cocycle throughout.

Gelbart's monograph \cite{Ge76} developed the representation theory
of the two-fold cover of $SL(2)$ through Weil's representation, and it
is the starting point for the study of genuine representations of
$p$-adic covering groups. Rallis and Schiffmann \cite{RS77} studied
the supercuspidal genuine representations of the two-fold cover.

Waldspurger \cite{Wa80} established the Shimura correspondence between
genuine representations of the two-fold cover of $SL(2)$ and
representations of $PGL(2)$, locally and globally, by means of the
theta correspondence. Flicker \cite{Fl80} compared Hecke algebras of
smooth functions on $GL(2,F)$ and on its $n$-fold covers by matching
orbital integrals, and combined this with the trace formula to obtain
a global Shimura correspondence. Kazhdan and Patterson \cite{KP84}
developed the theory of metaplectic forms on covers of $GL(r)$, and
Flicker and Kazhdan \cite{FK86} extended the trace formula comparison
to $GL(r)$. None of these works uses the Hecke algebras of open compact
subgroups that we study.

Howe and Moy \cite{HM85} described the Hecke algebras of the
filtration subgroups of an Iwahori subgroup of $GL(n)$ by generators
and relations. Our Section~\ref{section:sl2} adapts their method to
$SL(2)$.

As described in \cite{Ka21}, Savin \cite{Sa88} proved that for the
$n$-fold cover of a
simply-laced, simply connected split group with $\gcd(n,p)=1$, the
genuine Iwahori--Hecke algebra is isomorphic to the Iwahori--Hecke
algebra of the linear group $G/Z_n$, where $Z_n$ is the $n$-torsion of
the center of $G$. This is the Iwahori-level counterpart of our
question. For $SL(2)$ we reprove it in Theorem~\ref{thm:savinsl2},
using the presentation of Iwahori--Hecke algebras by Iwahori and
Matsumoto \cite{IM65}.

Bushnell and Kutzko \cite{BK98} introduced types. A type is a pair
consisting of an open compact subgroup and an irreducible
representation of it, whose Hecke algebra governs one block of the
category of smooth representations. Schultz \cite{Sc98} studied a
lifting of characters between $\overline{SL_2(F)}$ and $SO_{2,1}(F)$.

\subsubsection{Work after 1998}
Savin \cite{Sa04} studied the unramified representations of covering
groups of split groups more generally. Loke and Savin \cite{LS10}
treated the two-fold cover of $SL_2(\mathbb{Q}_2)$, where the residual
characteristic divides $n$.

W.-W.~Li \cite{Li11} developed endoscopy for the metaplectic groups
$Mp(2n)$. He proved the transfer of orbital integrals, and the
fundamental lemma for the unit element of the spherical Hecke algebra
when the residual characteristic is large. Luo \cite{Lu18} proved the
fundamental lemma for the whole spherical Hecke algebra of $Mp(2n)$.
For $n=1$ these results concern our cover with $n=2$, whose elliptic
endoscopic groups include $PGL(2)=SO(3)$.

Gan and Savin \cite{GS12} used types for the even and odd Weil
representations to define two Hecke algebras governing the
corresponding Bernstein components of the metaplectic group $Mp(2n)$
in odd residual characteristic, and identified them with Iwahori--Hecke
algebras of odd orthogonal groups. Takeda and Wood \cite{TW18}
extended this to even residual characteristic.

Karasiewicz \cite{Ka21} proved that a genuine Hecke algebra on the
double cover of a simple, simply-laced, simply connected Chevalley
group over $\mathbb{Q}_2$ admits a Bernstein presentation, and that it
is isomorphic to the Iwahori--Hecke algebra of $G/Z_2$.

Wang \cite{Wa24} studied the Hecke algebras of types for tame genuine
principal series. He showed that each shares its affine Hecke algebra
part with a Hecke algebra of the principal endoscopic group, and he
gave a counterexample showing that the two algebras are not isomorphic
in general. The principal endoscopic group has the dual group of the
cover in the sense of Weissman \cite{We18} as its dual group.

Witthaus \cite{Wi22} classified the irreducible genuine representations
modulo $p$ of the double cover of $GL(2,\mathbb{Q}_p)$, $p$ odd. He
works with the pro-$p$ Iwahori--Hecke algebra, which he decomposes by
idempotents attached to characters of the finite torus. The metaplectic
twist merges the orbits of these characters. In \cite[\S6.1]{Wi22} he
writes Savin's isomorphism explicitly for $GL(2)$, in the form that we
prove for $SL(2)$ in Theorem~\ref{thm:savinsl2}. He also refers to work
of Peskin \cite{Pe14} on the double cover of $SL(2,\mathbb{Q}_p)$
modulo $p$.

\subsubsection{How the present results fit}
The cited works on Hecke algebra correspondences match a genuine
Hecke algebra with the Hecke algebra of a \emph{different} linear
group, such as $G/Z_n$, an odd orthogonal group, or an endoscopic
group. Our results show that at positive level the comparison with $G$
itself, by a support-preserving map, fails. For every $n\geq 2$ it
fails for maps with constants of absolute value one, because the cover
does not split. For $n\geq 3$ it fails for all constants, because the
inverse image of the diagonal torus in $\Gbar$ is not commutative. This is consistent with the rescaling of translations in
Savin's theorem and with Wang's counterexample.

\subsection{Notation}
\label{section:notation}

This subsection lists the notation used from here on. Each symbol is also
defined at its first use.

\medskip
\begin{center}
\begin{tabular}{@{}p{0.27\textwidth}p{0.67\textwidth}@{}}
\toprule
Symbol & Meaning \\
\midrule
$F$, $\cO$, $\varpi$, $q$ & $p$-adic field, its integers, a uniformizer,
the residue field size \\
$\mu_n$, $(\ ,\ )_n$ & $n$-th roots of unity in $F$, $n$-th power
Hilbert symbol \\
$G$, $\Gbar$, $\pr$ & $SL(2,F)$, its $n$-fold metaplectic cover, the
projection $\Gbar\to G$ \\
$\beta$, $x(g)$ & Kubota's cocycle and the function used to define it \\
$\overline{H}$ & the inverse image $\pr^{-1}(H)$ of $H\subset G$ \\
$\tilde g$ & a lift of $g\in G$ to $\Gbar$ \\
$u_\pm(x)$, $h(a)$, $T$ & unipotent and diagonal elements, the diagonal
torus \\
$s_1$, $s_2$, $t$, $t_0$ & the elements of
\S\ref{subsec:subgroups}, with $t_0=h(\varpi)$ \\
$K$, $B$, $B_k$ & $SL(2,\cO)$, the Iwahori subgroup, its $k$-th
filtration subgroup \\
$\Gamma_k$, $U^\pm_k$, $T_k$ & principal congruence subgroup and its
unipotent and diagonal parts \\
$W_a$, $\ell$ & representatives of the affine Weyl group, the length \\
$f_g$ & characteristic function of $B_kgB_k$ \\
$\cH(G/\!/B_k)$ & generalized Hecke algebra of $G$ with respect to
$B_k$ \\
$\cP$, $\sigma$ & small subgroups with canonical lifts, the canonical
lift (Lemma~\ref{lemma:canonical}) \\
$s$ & a splitting of an open compact subgroup, for example of $K$ when
$p\nmid n$ \\
$w(x)$ & the element $u_+(x)u_-(-x^{-1})u_+(x)$ of $G$ \\
$B_k^*$, $K_0^*$ & lifts of $B_k$, $K_0$ to subgroups of $\Gbar$ \\
$\cH(\Gbar/\!/K_0^*)_{\gen}$ & genuine generalized Hecke algebra \\
$\fbar_{\tilde g}$ & genuine bi-$K_0^*$-invariant function with value
$1$ at $\tilde g$ \\
\bottomrule
\end{tabular}
\end{center}

\begin{example}
In the running example $F=\mathbb{Q}_7$, $n=3$, we may take
$\varpi=7$, so $t_0=h(7)$, $q=7$, $\cO=\mathbb{Z}_7$, and
$\mu_3=\{1,\omega,\omega^2\}$ with $\omega\in\mathbb{Z}_7$ a primitive
cube root of unity congruent to $2$ modulo $7$.
\end{example}

\subsection{Organization}
Section~\ref{section:intro} also contains a review of the literature
(Section~\ref{section:lit}) and a table of notation
(Section~\ref{section:notation}).
Section~\ref{section:background} collects background on the Hilbert
symbol, the cover, the relevant subgroups and Hecke algebras. Its last
subsection, Section~\ref{subsec:savinbg}, is a guide to Savin's
isomorphism for $SL(2)$. Part~I
consists of Sections~\ref{section:sl2}--\ref{section:n2}.
Section~\ref{section:sl2} treats $\Hlin$. Section~\ref{section:genuine}
constructs the genuine basis, and Section~\ref{section:metaplectic}
transfers the relations. These three sections hold for every $n$.
Section~\ref{section:n2} specializes to $n=2$ and proves the
non-existence result for constants of absolute value one, which holds
for every $n\geq 2$. Part~II consists of
Sections~\ref{section:torus}--\ref{section:savin}.
Section~\ref{section:torus} gives the torus obstruction for $n\geq 3$.
Section~\ref{section:reps} relates the algebras to representations.
Section~\ref{section:savin} treats Iwahori level for every $n$. It
gives a presentation of the genuine Iwahori--Hecke algebra, a direct
proof of Savin's theorem for $SL(2)$, and explicit orbital integrals at
split and elliptic elements.
Section~\ref{section:summary} summarizes the results and lists open
questions. A glossary follows.

%======================================================================
\section{Background}
\label{section:background}

This section explains the prerequisites: the Hilbert symbol, Kubota's
cover, the subgroups $B_k$ and their volumes, generalized Hecke
algebras, and Iwahori--Hecke algebras together with Savin's
isomorphism. Each subsection ends with an example from the running
example $F=\mathbb{Q}_7$, usually with $n=3$.

\subsection{The Hilbert symbol}
This subsection recalls the Hilbert symbol, the arithmetic input to the
cover. The fact we use most is that its values on pairs consisting of a
unit and the uniformizer fill out $\mu_n$ (Lemma~\ref{lemma:hilbert}).

The $n$-th power Hilbert symbol is a pairing
$(\ ,\ )_n:F^\times\times F^\times\to\mu_n$ from local class field
theory. We use the following standard properties.
\begin{enumerate}[leftmargin=3.2em]
\item[(H1)] It is bimultiplicative.
\item[(H2)] It is skew-symmetric: $(a,b)_n(b,a)_n=1$.
\item[(H3)] $(a,-a)_n=1$ for all $a$, and $(a,1-a)_n=1$ for $a\neq 1$.
\item[(H4)] $(a,b)_n=1$ for all $b\in F^\times$ if and only if
$a\in (F^\times)^n$.
\item[(H5)] If $p\nmid n$ then $(u,v)_n=1$ for all $u,v\in\cO^\times$.
\end{enumerate}
From (H1)--(H3) we get $(a,a)_n=(-1,a)_n$. Properties (H3) are the
Steinberg relations. Property (H4) says the pairing is non-degenerate
modulo $n$-th powers.

\begin{lemma}
\label{lemma:hilbert}
The map $u\mapsto (u,\varpi)_n$ sends $\cO^\times$ onto $\mu_n$. In
particular, if $n\geq 3$, there is a unit $u$ with
$(u,\varpi)_n^2\neq 1$.
\end{lemma}

\begin{proof}
The element $\varpi$ has order $n$ in $F^\times/(F^\times)^n$. By (H1)
and (H4), the character $a\mapsto (a,\varpi)_n$ of
$F^\times/(F^\times)^n$ also has order $n$, so its image is $\mu_n$.
Since $(\varpi,\varpi)_n=(-1,\varpi)_n$, we have
$(\varpi^mu,\varpi)_n=((-1)^mu,\varpi)_n$ for every unit $u$. So the
units already give the whole image. When $n\geq 3$, the squaring map
on $\mu_n$ is not trivial.
\end{proof}

\begin{example}
\label{ex:hilbert}
When $p\nmid n$, the symbol $(u,\varpi)_n$ of a unit $u$ is the image
in $\mu_n$ of $\bar u^{(q-1)/n}$, up to inversion depending on the
normalization, where $\bar u$ is the residue of $u$. For $n=2$ and $p$
odd this is the Legendre symbol of $\bar u$. In the running example
$F=\mathbb{Q}_7$, $n=3$, the cube roots of unity modulo $7$ are
$1,2,4$, and $(u,7)_3$ corresponds to $\bar u^{\pm 2}$. For $u=3$ we get
$3^2=9\equiv 2$, so $(3,7)_3$ and $(3,7)_3^2$ are primitive cube roots
of unity. For $n=2$ the squares modulo $7$ are $1,2,4$. So $3$ and $-1$
are not squares modulo $7$, and $(3,7)_2=(-1,7)_2=-1$.
\end{example}

\subsection{The metaplectic cover}
\label{subsec:cover}
This subsection defines Kubota's cover and computes it on the diagonal
torus and on the unipotent subgroups. The main facts are that the
inverse image of the torus is not commutative when $n\geq 3$, and that
the cover does not split over $G$ for any $n\geq 2$.

For $g=\mat abcd\in G$ put $x(g)=c$ if $c\neq 0$ and $x(g)=d$
otherwise. Kubota's cocycle is
\[
\beta(g_1,g_2)=\left(\frac{x(g_1g_2)}{x(g_1)},
\frac{x(g_1g_2)}{x(g_2)}\right)_n,\qquad g_1,g_2\in G.
\]
The group $\Gbar$ is the set of pairs $(g,\zeta)$ with $g\in G$ and
$\zeta\in\mu_n$, with multiplication
\[
(g_1,\zeta_1)(g_2,\zeta_2)=(g_1g_2,\beta(g_1,g_2)\zeta_1\zeta_2).
\]
The projection is $\pr(g,\zeta)=g$, and we identify $\zeta\in\mu_n$
with the central element $(1,\zeta)$. For $H\subset G$ we write
$\overline{H}=\pr^{-1}(H)$. A \emph{lift} of $g\in G$ is an element
$\tilde g\in\Gbar$ with $\pr(\tilde g)=g$. A \emph{splitting} of a
subgroup $H\subset G$ is a homomorphism $s:H\to\Gbar$ with
$\pr\circ s=\mathrm{id}_H$.

\begin{vremark}
\label{vrem:kubota}
\emph{Cited, not re-checked against the source.} By \cite[\S1]{Ku67},
$\beta$ is a $2$-cocycle, and its restriction to some open subgroup of
$G$ is a coboundary. So $\Gbar$ is a locally compact totally
disconnected group, and $\pr$ is a covering that splits over some open
subgroup of $G$. We rely on this throughout. Kubota also shows that when
$p\nmid n$ the cover splits over $K=SL(2,\cO)$. We use this only in
Proposition~\ref{prop:Ksub}. The formula for $\beta$
can be rewritten with (H1)--(H3) as
$(x(g_1),x(g_2))_n\,(-x(g_1)^{-1}x(g_2),x(g_1g_2))_n$.
\end{vremark}

Let
\[
u_+(x)=\mat 1x01,\qquad u_-(x)=\mat 10x1,\qquad h(a)=\mat a00{a^{-1}},
\]
and let $U^\pm=\{u_\pm(x)\,|\,x\in F\}$ and
$T=\{h(a)\,|\,a\in F^\times\}$.

\begin{lemma}
\label{lemma:cover}
\begin{enumerate}[leftmargin=2em,label=(\alph*)]
\item $\beta(h(a),h(b))=(b,a)_n$. Consequently, for any lifts,
\[
\tilde h(a)\,\tilde h(b)=(b,a)_n^2\ \tilde h(b)\,\tilde h(a).
\]
\item $\beta$ is identically $1$ on $U^+\times U^+$ and on
$U^-\times U^-$. So $u_\pm(x)\mapsto (u_\pm(x),1)$ are splittings of
$U^\pm$, and they are the only splittings of $U^\pm$.
\item The subgroup $\overline{\{\pm 1\}}$ is central in $\Gbar$.
\end{enumerate}
\end{lemma}

\begin{proof}
(a) We have $x(h(a))=a^{-1}$, so
$\beta(h(a),h(b))=(b^{-1},a^{-1})_n=(b,a)_n$. Hence
\[
(h(a),1)(h(b),1)=(b,a)_n(a,b)_n^{-1}\,(h(b),1)(h(a),1),
\]
and $(a,b)_n^{-1}=(b,a)_n$ by (H2). The relation does not depend on
the lifts, because $\mu_n$ is central.

(b) On $U^+$ the function $x$ is identically $1$, so $\beta=1$ there.
On $U^-$, let $g_1=u_-(x)$ and $g_2=u_-(y)$. If $x=0$ or $y=0$, then
$\beta(g_1,g_2)=1$ directly. If $x+y=0$, then $x(g_1g_2)=1$ and
$\beta(g_1,g_2)=(x^{-1},-x^{-1})_n=1$ by (H3). Otherwise put
$a=(x+y)/x$. Then $(x+y)/y=a/(a-1)$, and
\[
\beta(g_1,g_2)=(a,a)_n\,(a,a-1)_n^{-1}.
\]
Here $(a,a)_n=(-1,a)_n=(a,-1)_n^{-1}$. Also
$(a,a-1)_n=(a,-1)_n(a,1-a)_n=(a,-1)_n$. Hence
\[
\beta(g_1,g_2)=(a,-1)_n^{-2}=(a,(-1)^{-2})_n=(a,1)_n=1.
\] Two splittings of
$U^\pm\cong F$ differ by a homomorphism $F\to\mu_n$. Such a
homomorphism is trivial because $F$ is divisible.

(c) Let $z=\pm 1$ with lift $\tilde z$. The map
$g\mapsto\tilde z\tilde g\tilde z^{-1}\tilde g^{-1}$ does not depend on
the lift $\tilde g$, and it is a homomorphism $G\to\mu_n$ because $z$
is central in $G$. It is trivial because $G$ is perfect.
\end{proof}

\begin{example}
\label{ex:cover}
For $n=2$, Lemma~\ref{lemma:cover}(a) shows that $\overline{T}$ is
abelian. In the running example $n=3$, $F=\mathbb{Q}_7$, take $a=7$ and
$b=3$. Then $\tilde h(7)\tilde h(3)=(3,7)_3^2\,\tilde h(3)\tilde
h(7)$, and $(3,7)_3^2\neq 1$ by Example~\ref{ex:hilbert}. So
$\overline{T}$ is not abelian. This elementary fact is the source of
Theorem~\ref{thrm:main}.
\end{example}

For $n=2$ the torus gives no obstruction. The next lemma shows that the
cover is nevertheless not split. For $x\in F^\times$ put
\[
w(x)=u_+(x)\,u_-(-x^{-1})\,u_+(x)=\mat 0x{-x^{-1}}0 .
\]
Then $w(-1)=s_1$ and $w(a)w(-1)=h(a)$ for $a\in F^\times$.

\begin{lemma}
\label{lemma:nonsplit}
Let $n\geq 2$. The cover $\Gbar\to G$ does not split, even as a
homomorphism of abstract groups. More precisely, if $s:G\to\Gbar$ were
a splitting, then $s(h(a))\,s(h(b))=(b,a)_n\,s(h(ab))$ for all
$a,b\in F^\times$.
\end{lemma}

\begin{proof}
Suppose that $s$ is a splitting. By Lemma~\ref{lemma:cover}(b), $s$
agrees with $u\mapsto(u,1)$ on $U^+$ and on $U^-$. We compute two
values of $\beta$. First, $u_+(x)u_-(y)=\mat{1+xy}xy1$, so
$\beta(u_+(x),u_-(y))=(y,1)_n=1$ for $y\neq 0$. Second, put
$g=u_+(x)u_-(-x^{-1})=\mat 0x{-x^{-1}}1$. Then $x(g)=-x^{-1}$,
$x(u_+(x))=1$ and $x(gu_+(x))=x(w(x))=-x^{-1}$, so
$\beta(g,u_+(x))=(1,-x^{-1})_n=1$. Hence $s(w(x))=(w(x),1)$. Next,
$x(w(a))=-a^{-1}$, $x(w(-1))=1$ and $x(h(a))=a^{-1}$, so
\[
s(h(a))=(w(a),1)(w(-1),1)=\bigl(h(a),(-1,a^{-1})_n\bigr).
\]
By Lemma~\ref{lemma:cover}(a), $\beta(h(a),h(b))=(b,a)_n$, so
\[
s(h(a))\,s(h(b))=\bigl(h(ab),(b,a)_n(-1,a^{-1})_n(-1,b^{-1})_n\bigr)
=(b,a)_n\,s(h(ab)).
\]
Since $s$ is a homomorphism, $(b,a)_n=1$ for all $a,b$. This
contradicts Lemma~\ref{lemma:hilbert}, which gives a unit $u$ with
$(u,\varpi)_n\neq 1$.
\end{proof}

\begin{example}
\label{ex:nonsplit}
In the running example with $n=2$, take $a=7$ and $b=3$. A splitting
would satisfy $s(h(7))s(h(3))=(3,7)_2\,s(h(21))=-s(h(21))$, by
Example~\ref{ex:hilbert}. So the cover of $SL(2,\mathbb{Q}_7)$ of
degree~$2$ does not split, although its torus $\overline{T}$ is
abelian.
\end{example}

\subsection{Subgroups}
\label{subsec:subgroups}
This subsection introduces the Iwahori subgroup $B$, its filtration
subgroups $B_k$ and the principal congruence subgroups $\Gamma_k$, and
it computes the volumes of the double cosets $B_kwB_k$. These volumes
decide when Hecke functions multiply without extra terms.

Let
\[
s_1=\mat 0{-1}10,\quad
s_2=\mat 0{-\varpi^{-1}}{\varpi}0,\quad
t=\mat 0{-1}{\varpi}0,\quad
t_0=h(\varpi)=\mat{\varpi}00{\varpi^{-1}}.
\]
Then $s_1^2=s_2^2=-1$, $t_0=s_1^{-1}s_2$, $t\notin G$, and
$ts_1t^{-1}=s_2$. Moreover $s_1t_0^{-1}s_1^{-1}=t_0$ and
$t_0^{-1}s_1=s_2$. Note that $s_1$ has order $4$, whereas the permutation
matrix used in \cite{HM85} for $GL(n)$ has order $2$.

Let $\fB=\mat{\cO}{\cO}{\varpi\cO}{\cO}$. The Iwahori subgroup is
$B=\fB^\times\cap G$, the group of matrices in $SL(2,\cO)$ that are
upper triangular modulo $\varpi$. For $k\geq 1$ the $k$-th filtration
subgroup and its Lie algebra are
\[
B_k=(1+\varpi^k\fB)\cap G,\qquad
\fb_k=\left\{\mat ab{\varpi c}{-a}\ \Big|\ a,b,c\in\varpi^k\cO\right\}.
\]
We also put $B_0=B$. Since $t^2=-\varpi$, we have
$B_k=(1+t^{2k}\fB)\cap G$. Let $K=SL(2,\cO)$, and for $j\geq 1$ let
\[
\Gamma_j=G\cap\bigl(1+\varpi^jM_2(\cO)\bigr)
\]
be the principal congruence subgroup of level $j$. Let
$U^\pm_j=\{u_\pm(x)\,|\,x\in\varpi^j\cO\}$ and
$T_j=\{h(a)\,|\,a\in 1+\varpi^j\cO\}$.

\begin{lemma}
\label{lemma:subgroups}
Let $k\geq 1$.
\begin{enumerate}[leftmargin=2em,label=(\alph*)]
\item $B_k$ is a normal subgroup of $B$, and $t$ normalizes $B$ and
$B_k$.
\item The map sending $\mat{1+\varpi^ka}{\varpi^kb}{\varpi^{k+1}c}{d}$
to $(a,b,c)$ modulo $\varpi$ is a surjective homomorphism
$B_k\to\Fq^3$ with kernel $B_{k+1}$. Likewise, the map sending
$\mat ab{\varpi c}{-a}$ to $\varpi^{-k}(a,b,c)$ modulo $\varpi$
induces an isomorphism $\fb_k/\fb_{k+1}\cong\Fq^3$. In particular
$B_k$ is a pro-$p$ group.
\item (Iwahori factorization.) Every element of $B_k$ can be written
uniquely as $u_-(x)h(a)u_+(y)$ with $x\in\varpi^{k+1}\cO$,
$a\in 1+\varpi^k\cO$ and $y\in\varpi^k\cO$. That is,
$B_k=U^-_{k+1}T_kU^+_k$. Similarly $\Gamma_k=U^-_kT_kU^+_k$.
\item $B_k\cdot(s_1^{-1}B_ks_1)=(s_1^{-1}B_ks_1)\cdot B_k
=U_k^-B_k=\Gamma_k$. The group $B_k$ is normal in $\Gamma_k$ of index
$q$. Similarly $B_k\cdot(s_2^{-1}B_ks_2)=t\Gamma_kt^{-1}=U^+_{k-1}B_k$.
\item $\Gamma_k$ is normal in $K$. For $x\in\cO^\times$ we have
$B_ku_-(x)B_k=u_-(x)\Gamma_k$, which has volume $q\cdot\vol(B_k)$.
\end{enumerate}
\end{lemma}

\begin{proof}
(a) For $b\in\fB^\times$ we have $b(1+\varpi^k\fB)b^{-1}=1+\varpi^k\fB$.
Also $t\fB t^{-1}=\fB$.

(b) For $g=1+X$ and $g'=1+X'$ in $B_k$ we have $gg'=1+X+X'+XX'$ with
$XX'\in\varpi^{2k}\fB\subset\varpi^{k+1}\fB$. This gives the
homomorphism property. Surjectivity holds because, given $a$, $b$ and
$c$, the entry $d$ can be solved for from $\det=1$. The statement for
$\fb_k$ is clear.

(c) If $g=\mat{a}{b}{c}{d}\in B_k$, then $a$ is a unit and
$g=u_-(c/a)h(a)u_+(b/a)$. Uniqueness is clear. The same computation
applies to $\Gamma_k$.

(d) Conjugation by $s_1$ gives
$s_1^{-1}\mat abcd s_1=\mat d{-c}{-b}a$. So
$s_1^{-1}B_ks_1=U^-_kT_kU^+_{k+1}$. The product $(s_1^{-1}B_ks_1)B_k$
is contained in $\Gamma_k$ and contains $U^-_kT_kU^+_k=\Gamma_k$ by
(c). Taking inverses gives the same for $B_k(s_1^{-1}B_ks_1)$. The map
$1+\varpi^kX\mapsto X$ modulo $\varpi$ identifies
$\Gamma_k/\Gamma_{k+1}$ with the abelian group $\mathfrak{sl}_2(\Fq)$,
and $B_k$ is the preimage of the upper triangular matrices. Hence
$B_k$ is normal in $\Gamma_k$ of index $q$. The statement for $s_2$
follows by conjugating with $t$. Here we use $ts_1t^{-1}=s_2$, part
(a), and $tu_-(x)t^{-1}=u_+(-x/\varpi)$.

(e) $\Gamma_k$ is the kernel of reduction modulo $\varpi^k$ on $K$.
Under the identification in (d), $K$ acts on $\mathfrak{sl}_2(\Fq)$ by
the adjoint action. For a unit $x$, the subalgebra of upper triangular
matrices is not stable under $u=u_-(x)$. Two distinct planes in the
three-dimensional space $\mathfrak{sl}_2(\Fq)$ meet in a line. Both
$B_k$ and $uB_ku^{-1}$ contain $\Gamma_{k+1}$, because $u\in K$
normalizes $\Gamma_{k+1}$. So $[B_k:B_k\cap uB_ku^{-1}]=q$, and
$\vol(B_kuB_k)=q\,\vol(B_k)=\vol(u\Gamma_k)$. Since
$B_kuB_k\subset u\Gamma_k$, the two sets are equal.
\end{proof}

By the Iwahori--Bruhat decomposition, $G$ is the disjoint union of the
double cosets $BwB$, where $w$ runs over
\[
W_a=\{t_0^m,\ t_0^ms_1\ |\ m\in\Z\},
\]
a set of representatives for the affine Weyl group. Define the length
by $\ell(t_0^m)=2|m|$ and $\ell(t_0^ms_1)=|2m+1|$, so that
$\ell(s_1)=\ell(s_2)=1$. Put $\ell(-w)=\ell(w)$.

\begin{lemma}
\label{lemma:length}
For $k\geq 0$ and $w\in \pm W_a$, we have
$\vol(B_kwB_k)=q^{\ell(w)}\vol(B_k)$.
\end{lemma}

\begin{proof}
Since $-1$ is central, we may assume $w\in W_a$. We have
$\vol(B_kwB_k)/\vol(B_k)=[B_k:B_k\cap wB_kw^{-1}]$. Conjugation by
$t_0^m$ multiplies the upper right entry by $\varpi^{2m}$ and the lower
left entry by $\varpi^{-2m}$. Conjugation by $s_1$ was computed in the
proof of Lemma~\ref{lemma:subgroups}(d). Hence $wB_kw^{-1}$ is the set
of elements of $G$ whose diagonal entries lie in $1+\varpi^k\cO$, whose
upper right entry lies in $\varpi^i\cO$, and whose lower left entry
lies in $\varpi^j\cO$. Here $(i,j)=(k+2m,\,k+1-2m)$ for $w=t_0^m$, and
$(i,j)=(k+1+2m,\,k-2m)$ for $w=t_0^ms_1$. For $k=0$, read $\cO^\times$ for $1+\varpi^0\cO$. The intersection
with $B_k$ is described by the exponents $\max(i,k)$ and
$\max(j,k+1)$. By the factorization of Lemma~\ref{lemma:subgroups}(c),
applied to both groups, or for $k=0$ by the Iwahori factorization
$B=U^-_1T(\cO)U^+_0$, which is proved in the same way, the index is $q^{e}$ with
$e=(\max(i,k)-k)+(\max(j,k+1)-k-1)$. This exponent is $2|m|$ in the
first case and $|2m+1|$ in the second.
\end{proof}

\begin{example}
\label{ex:subgroups}
In the running example $q=7$. For $k=1$, the quotient $B_1/B_2$ is
isomorphic to $\mathbb{F}_7^3$. The double cosets $B_1s_1B_1$ and
$B_1t_0B_1$ have volumes $7$ and $49$ times $\vol(B_1)$, and
$B_1t_0s_1B_1$ has volume $343$ times $\vol(B_1)$.
\end{example}

\subsection{Generalized Hecke algebras}
This subsection recalls generalized Hecke algebras and the basic rule
for multiplying characteristic functions of double cosets. Everything
here holds for any locally profinite unimodular group.

\begin{definition}
\label{def:1}
Let $H$ be a locally compact, totally disconnected, unimodular group,
and let $K_0$ be an open compact subgroup of $H$. Normalize Haar
measure on $H$ so that $K_0$ has measure one. Let $\cH(H/\!/K_0)$
denote the space of functions on $H$ that are locally constant,
compactly supported, and bi-$K_0$-invariant. With the convolution
\[
(f*g)(x)=\int_H f(y)g(y^{-1}x)\,dy
\]
it is an algebra, the \textbf{generalized Hecke algebra}. For $h\in H$,
$f_h$ denotes the characteristic function of $K_0hK_0$.
\end{definition}

\begin{example}
\label{ex:hecke}
If $H$ is finite and $K_0=\{1\}$, then $\cH(H/\!/K_0)$ is the group
algebra $\C[H]$, and $f_h*f_{h'}=f_{hh'}$. If $K_0$ is normal in $H$,
then $\cH(H/\!/K_0)$ is the group algebra of the discrete group
$H/K_0$. For $H=G$ and $K_0=B_k$, the group $B$ normalizes $B_k$, so
the $f_x$ with $x\in B$ span a copy of $\C[B/B_k]$.
\end{example}

\begin{lemma}
\label{lemma:general}
In the notation of Definition~\ref{def:1}, the following hold.
\begin{enumerate}[leftmargin=2em,label=(\alph*)]
\item $f_1$ is the identity of $\cH(H/\!/K_0)$.
\item $(f_{h_1}*f_{h_2})(x)=\vol(K_0h_1K_0\cap xK_0h_2^{-1}K_0)$. For any
$\phi\in\cH(H/\!/K_0)$,
\[
(f_h*\phi)(x)=\sum_{v\in K_0/(K_0\cap hK_0h^{-1})}\phi(h^{-1}v^{-1}x).
\]
\item $\int_H f*g=\int_H f\cdot\int_H g$.
\item If $\vol(K_0h_1K_0)\vol(K_0h_2K_0)=\vol(K_0h_1h_2K_0)$, then
$f_{h_1}*f_{h_2}=f_{h_1h_2}$ and
\begin{equation}
\label{eq:unique}
\{x\in K_0h_1K_0\ |\ x^{-1}h_1h_2\in K_0h_2K_0\}=h_1K_0 .
\end{equation}
\end{enumerate}
\end{lemma}

\begin{proof}
Parts (a), (b) and (c) follow from the definitions. For (b) we use
$K_0hK_0=\bigsqcup_v vhK_0$. For (d), part (b) shows that
$f_{h_1}*f_{h_2}$ is bi-$K_0$-invariant, takes non-negative integer
values, and is at least $1$ on $K_0h_1h_2K_0$. By (c) its integral is
$\vol(K_0h_1h_2K_0)$. Hence it equals $1$ on $K_0h_1h_2K_0$ and
vanishes elsewhere. Evaluating (b) at $x=h_1h_2$ gives
$\vol(K_0h_1K_0\cap h_1h_2K_0h_2^{-1}K_0)=1$. This set contains
$h_1K_0$ and is a union of left cosets of $K_0$, which gives
(\ref{eq:unique}). Part (d) is \cite[Chapter~3,
Proposition~2.2]{HM85}.
\end{proof}

\subsection{Iwahori--Hecke algebras and Savin's isomorphism}
\label{subsec:savinbg}
This subsection gives background on the Iwahori--Hecke algebras of
$SL(2)$ and $PGL(2)$, on the lattice of translations that the genuine
theory sees, and on Savin's isomorphism. It explains what Savin's
theorem says for $SL(2)$, why its target is $SL(2)$ for odd $n$ and
$PGL(2)$ for even $n$, how the isomorphism acts on double cosets, and
what it means for representations. The statements that concern the
cover are proved in Section~\ref{section:savin}, which gives a direct
proof of Savin's theorem for $SL(2)$ (Theorem~\ref{thm:savinsl2}). In
this subsection $p\nmid n$, and we put
\begin{equation}
\label{eq:nprime}
n'=\frac{n}{\gcd(n,2)},\qquad r=t_0^{-n'}s_1 .
\end{equation}
Thus $n'=n$ for odd $n$ and $n'=n/2$ for even $n$. Since
$t_0^{-1}s_1=s_2$, we have $r=s_2$ when $n=2$. In the notation of
Section~\ref{subsec:cover}, $r=w(-\varpi^{-n'})$ and $r^2=-1$.

\subsubsection{The Iwahori--Hecke algebra of \texorpdfstring{$SL(2)$}{SL(2)}}
Normalize Haar measure on $G$ by $\vol(B)=1$, and write $f_w$ for the
characteristic function of $BwB$. By Lemma~\ref{lemma:length}, which
also holds for $k=0$, we have $\vol(BwB)=q^{\ell(w)}$. Since $-1\in B$,
the functions $f_w$ with $w\in W_a$ form a basis of $\cH(G/\!/B)$.
Modulo $\pm 1$, the set $W_a$ is the infinite dihedral group generated
by the images of $s_1$ and $s_2$. Every $w\in W_a$ is, up to sign, a
unique alternating product of $s_1$ and $s_2$, and $\ell(w)$ is the
number of factors. For example, $t_0=s_1^{-1}s_2$ has length $2$.

\begin{proposition}[Iwahori--Matsumoto {\cite[\S3]{IM65}}]
\label{prop:IM}
\begin{enumerate}[leftmargin=2em,label=(\alph*)]
\item If $\ell(xy)=\ell(x)+\ell(y)$, then $f_x*f_y=f_{xy}$.
\item $f_{s_i}*f_{s_i}=(q-1)f_{s_i}+qf_1$ for $i=1,2$.
\item The algebra $\cH(G/\!/B)$ is generated by $f_{s_1}$ and
$f_{s_2}$, and the two relations in (b) are defining relations.
\end{enumerate}
\end{proposition}

\begin{proof}
Part (a) is Lemma~\ref{lemma:general}(d). Parts (b) and (c) are a
special case of \cite[\S3]{IM65}. They are also the case of the trivial
cover in Proposition~\ref{prop:quadratic} and
Theorem~\ref{thm:presentation} below. The proofs given there apply
verbatim when every Hilbert symbol is replaced by $1$ and $n'$ by $1$.
\end{proof}

There is no braid relation, because $s_1s_2$ has infinite order. Each
$f_{s_i}$ is invertible, with inverse $q^{-1}(f_{s_i}-(q-1)f_1)$.

\begin{remark}[Bernstein's form]
\label{rem:bernstein}
Put $\theta=q^{-1}f_{t_0^{-1}}$. Since $t_0^{-1}=-s_2s_1$ and the
lengths add, $\theta=q^{-1}f_{s_2}*f_{s_1}$, so $\theta$ is
invertible, and $\theta^m=q^{-m}f_{t_0^{-m}}$ for $m\geq 0$ by
Proposition~\ref{prop:IM}(a). Expanding $\theta^{-1}=qf_{s_1}^{-1}*
f_{s_2}^{-1}$ with the formula for the inverses gives
\[
f_{s_1}*\theta-\theta^{-1}*f_{s_1}=(q-1)(\theta+f_1).
\]
This is the relation between the finite part and the translations in
Bernstein's presentation (see the Glossary). The translations are
scaled by $q^{-\ell/2}$, since $\ell(t_0^{-m})=2m$.
\end{remark}

\subsubsection{The Iwahori--Hecke algebra of \texorpdfstring{$PGL(2)$}{PGL(2)}}
Let $G'=PGL(2,F)$, and write $[g]$ for the image of $g\in GL(2,F)$. Let
$I'$ be the image of the Iwahori subgroup of $GL(2,\cO)$, the group of
matrices in $GL(2,\cO)$ whose lower left entry lies in $\varpi\cO$.
Normalize Haar measure by $\vol(I')=1$, and write $f'_w$ for the
characteristic function of $I'wI'$. The element $[t]$ normalizes $I'$,
and $[t]^2=[-\varpi]=1$. Put
\[
\lambda=\left[\mat{\varpi}001\right]=[s_1][t].
\]
The image of $t_0$ in $G'$ is $\lambda^2$. The elements $\lambda^k$ and
$\lambda^k[s_1]$, for $k\in\Z$, represent the double cosets
$I'\backslash G'/I'$. They form the extended affine Weyl group, which is
the infinite dihedral group generated by the two involutions $[s_1]$
and $[t]$. Every element is a unique alternating product of $[s_1]$ and
$[t]$, and its length $\ell'$ is the number of factors $[s_1]$. Thus
$\ell'(\lambda^k)=|k|$, while $\ell'(\lambda^k[s_1])$ is $k+1$ for
$k\geq 0$ and $-k-1$ for $k<0$. In particular $\ell'(\lambda)=1$ and
$\ell'([t])=0$. As for $SL(2)$, $\vol(I'wI')=q^{\ell'(w)}$
\cite[\S3]{IM65}.

\begin{proposition}[Iwahori--Matsumoto {\cite[\S3]{IM65}}]
\label{prop:IMpgl}
\begin{enumerate}[leftmargin=2em,label=(\alph*)]
\item If $\ell'(xy)=\ell'(x)+\ell'(y)$, then $f'_x*f'_y=f'_{xy}$.
\item $f'_{[s_1]}*f'_{[s_1]}=(q-1)f'_{[s_1]}+qf'_1$ and
$f'_{[t]}*f'_{[t]}=f'_1$.
\item The algebra $\cH(G'/\!/I')$ is generated by $f'_{[s_1]}$ and
$f'_{[t]}$, and the relations in (b) are defining relations.
\end{enumerate}
\end{proposition}

\begin{proof}
Part (a) is Lemma~\ref{lemma:general}(d), and the second relation in
(b) is a case of it. The rest is a special case of \cite[\S3]{IM65}. It
also follows from the argument of Theorem~\ref{thm:presentation}, with
alternating products of $[s_1]$ and $[t]$ in place of those of $s_1$
and $r$.
\end{proof}

The generator $f'_{[t]}$ has eigenvalues $\pm1$, whereas $f_{s_2}$ has
eigenvalues $q$ and $-1$. This is the whole difference between the two
presentations. In the language of Hecke algebras with unequal
parameters, the two generators of $\cH(G/\!/B)$ both have parameter
$q$, while those of $\cH(G'/\!/I')$ have parameters $q$ and $1$. The
analog of Remark~\ref{rem:bernstein} is as follows. Put
$\theta'=q^{-1/2}f'_{\lambda^{-1}}=q^{-1/2}f'_{[t]}*f'_{[s_1]}$. Then
$\theta'^2=q^{-1}f'_{\lambda^{-2}}$ corresponds to $\theta$, and
\[
f'_{[s_1]}*\theta'-\theta'^{-1}*f'_{[s_1]}=(q-1)\,\theta' .
\]

\subsubsection{The lattice seen by the genuine algebra}
Let $s$ be the splitting of $B$ constructed at the beginning of
Section~\ref{section:savin}, and $B^*=s(B)$. By
Proposition~\ref{prop:level0} and Lemma~\ref{lemma:level0s1}, the
double cosets $BwB$, $w\in W_a$, that carry non-zero genuine
bi-$B^*$-invariant functions are those of $t_0^m$ and $t_0^ms_1$ with
$n\mid 2m$, that is, with $n'\mid m$. So the genuine algebra sees only
the translations $t_0^m$ with $m\in n'\Z$, and only the reflections
$t_0^ms_1$ with $m\in n'\Z$. The element $r$ is the genuine analog of
$s_2$. Modulo $\pm1$, the admissible elements form the infinite
dihedral group generated by $s_1$ and $r$ (Lemma~\ref{lemma:words}),
and $\ell(r)=2n'-1$.

In the language of root data, the translations $t_0^m=h(\varpi^m)$
correspond to the multiples $m\alpha^\vee$ of the coroot of $SL(2)$. By
Lemma~\ref{lemma:cover}(a), the commutator of lifts of $h(a)$ and
$h(b)$ is $(b,a)_n^2$. In the notation of \cite[\S2.1]{Wa24}, which
follows \cite{We18}, this means that Kubota's cover has
$Q(\alpha^\vee)=1$ and $B_Q(\alpha^\vee,\alpha^\vee)=2$, up to sign
conventions. The lattice
\[
Y_{Q,n}=\{y\in\Z\alpha^\vee\ |\ B_Q(y,\alpha^\vee)\in n\Z\}
=n'\Z\alpha^\vee
\]
is the lattice found above. The modified coroot is
$n_\alpha\alpha^\vee=n\alpha^\vee$, where
$n_\alpha=n/\gcd(n,Q(\alpha^\vee))=n$. The split group $G_{Q,n}$ with
cocharacter lattice $Y_{Q,n}$ and coroot $n\alpha^\vee$ is simply
connected when $n\alpha^\vee$ generates $Y_{Q,n}$, that is, when $n$ is
odd. It is adjoint when $n\alpha^\vee=2\cdot n'\alpha^\vee$, that is,
when $n$ is even. So $G_{Q,n}$ is $SL(2)$ for odd $n$ and $PGL(2)$ for
even $n$, in agreement with $G/Z_n$ in Theorem~\ref{thm:savin} below.
The dual group of the cover in the sense of \cite{We18} is the dual
group of $G_{Q,n}$, namely $PGL(2,\C)$ for odd $n$ and $SL(2,\C)$ for
even $n$. For $n=2$ this is the familiar fact that the double cover of
$SL(2)$ behaves like a group with dual group $Sp(2,\C)=SL(2,\C)$, and
its genuine representations are matched with those of $PGL(2)=SO(3)$
\cite{Wa80}. We use this dictionary only for orientation. The results
of Section~\ref{section:savin} are proved directly.
Table~\ref{tab:lattice} collects the dictionary.

\begin{table}[ht]
\caption{The dictionary of Savin's isomorphism for $SL(2)$, with
$n'=n/\gcd(n,2)$ and $r=t_0^{-n'}s_1$. The last row but one refers to
$w=t_0^{n'k}$ and $w=t_0^{n'k}s_1$.}
\label{tab:lattice}
\begin{tabular}{@{}lll@{}}
\toprule
 & $n$ odd & $n$ even \\
\midrule
genuine translations & $t_0^m$ with $n\mid m$ & $t_0^m$ with $\frac
n2\mid m$ \\
partner group $G'$ & $SL(2,F)$ & $PGL(2,F)$ \\
image of $t_0^{n'}$ & $t_0$, of length $2$ & $\lambda$, of length $1$ \\
genuine reflection $r$ & length $2n-1$ & length $n-1$ \\
image of $r$ & $s_2$, of length $1$ & $[t]$, of length $0$ \\
relation of $Y=\fbar_{(r,1)}$ & $Y^2=q^{n-1}(q-1)Y+q^{2n-1}$ &
$Y^2=(-1)^{\frac{q-1}{2}}q^{n-1}$ \\
scaling of $\fbar_{(w,1)}$ & $q^{(n-1)|k|}$ & $q^{(n-1)|k|/2}$ \\
dual group & $PGL(2,\C)$ & $SL(2,\C)$ \\
\bottomrule
\end{tabular}
\end{table}

\subsubsection{Savin's theorem}
Savin proved the following for the covers of a large class of split
groups. We state it for $SL(2)$.

\begin{theorem}[Savin {\cite{Sa88}}, as described in {\cite{Ka21}}]
\label{thm:savin}
Assume $\gcd(p,n)=1$. The genuine Iwahori--Hecke algebra of the
$n$-fold cover $\Gbar$ of $G=SL(2,F)$ is isomorphic to the
Iwahori--Hecke algebra of the split group $G/Z_n$, where $Z_n$ is the
$n$-torsion of the center of $SL_2$.
\end{theorem}

\begin{remark}
\label{rem:savin}
The quotient $G/Z_n$ is taken as algebraic groups. Its group of
$F$-points is $SL(2,F)$ if $n$ is odd and $PGL(2,F)$ if $n$ is even.
\end{remark}

\begin{vremark}
\label{vrem:savin}
\emph{Cited from secondary sources.} We have not been able to consult
\cite{Sa88}. The general statement of Theorem~\ref{thm:savin} follows
\cite{Ka21}. An AI-written summary of \cite{Sa88} was also available to
us, and we used it only for orientation. For $SL(2)$ the theorem is
proved below (Theorem~\ref{thm:savinsl2}), with an explicit
isomorphism and independently of these sources. Whether our
normalization agrees with Savin's in every detail should still be
checked in \cite{Sa88}.
\end{vremark}

For $SL(2)$ the isomorphism can be written down explicitly. Let
$X=\fbar_{(s_1,1)}$ and $Y=\fbar_{(r,1)}$ be the genuine functions on
$\overline{Bs_1B}$ and $\overline{BrB}$ with value $1$ at $(s_1,1)$
and $(r,1)$, and write $1$ for the identity element. Then $X$ and $Y$
generate the genuine Iwahori--Hecke algebra, with the defining relations
\[
X^2=(q-1)X+q,\qquad Y^2=c_1Y+c_0,
\]
where $c_1=q^{n-1}(q-1)$ and $c_0=q^{2n-1}$ for odd $n$, and $c_1=0$
and $c_0=(-1)^{(q-1)/2}q^{n-1}$ for even $n$
(Theorem~\ref{thm:presentation}). Compare this with
Propositions~\ref{prop:IM} and \ref{prop:IMpgl}. For odd $n$, the
element $q^{1-n}Y$ satisfies the relation of $f_{s_2}$. For even $n$,
the element $a^{-1}Y$ satisfies the relation of $f'_{[t]}$ when
$a^2=c_0$. Savin's isomorphism sends $X$ to $f_{s_1}$, respectively
$f'_{[s_1]}$, and $Y$ to $q^{n-1}f_{s_2}$, respectively to
$af'_{[t]}$ (Theorem~\ref{thm:savinsl2}).

On the basis it acts as follows. Let $\Psi$ send $t_0^{n'k}$ to
$\mu^k$ and $t_0^{n'k}s_1$ to $\mu^ks_1$, where $\mu=t_0$ for odd $n$
and $\mu=\lambda$ for even $n$, and where $s_1$ stands for $[s_1]$ in
$PGL(2)$. Write $f'=f$ and $\ell'=\ell$ when $n$ is odd. Then
\begin{equation}
\label{eq:savinbasis}
\fbar_{(w,1)}\longmapsto c_w\,q^{e(w)}f'_{\Psi(w)},\qquad
e(w)=\frac{\ell(w)-\ell'(\Psi(w))}{2}=\frac{(n-1)|k|}{\gcd(n,2)},
\end{equation}
with $|c_w|=1$. Since $\vol(BwB)=q^{\ell(w)}$ and
$\vol(I'\Psi(w)I')=q^{\ell'(\Psi(w))}$, the factor $q^{e(w)}$ is the
square root of the ratio of the two volumes. So the isomorphism sends
$\fbar_{(w,1)}/\vol(BwB)^{1/2}$ to
$c_wf'_{\Psi(w)}/\vol(I'\Psi(w)I')^{1/2}$. These normalized functions
are orthonormal for the inner product
$\langle\phi,\psi\rangle=\int\phi\,\overline\psi$. Thus the
isomorphism carries an orthonormal basis to an orthonormal basis, up to
constants of absolute value one. This is the precise sense in which
Savin's isomorphism rescales the translations by powers of $\sqrt q$.
For example, for $n=2$ the genuine function on $\overline{Bt_0^kB}$
goes to $q^{|k|/2}$ times a constant times $f'_{\lambda^k}$. For $n=3$
the genuine function on $\overline{Bt_0^{3k}B}$ goes to $q^{2|k|}$
times a constant times $f_{t_0^k}$.

Witthaus \cite[Lemma~6.1]{Wi22} writes down Savin's isomorphism in the
same form for the double cover of $GL(2,\mathbb{Q}_p)$, $p$ odd, whose
linear partner is $GL(2,\mathbb{Q}_p)$. The genuine function on the
double coset of $\mathrm{diag}(p^{2a},p^{2b})$ with $a\geq b$ goes to
$p^{(a-b)/2}$ times the characteristic function of the double coset of
$\mathrm{diag}(p^a,p^b)$. The genuine function on the double coset of
$s$ goes to the characteristic function of the double coset of $s$. The
exponent $(a-b)/2$ is again half the difference of the two lengths,
$2(a-b)$ and $a-b$.

\subsubsection{What the isomorphism means for representations}
By Borel's theorem \cite{Bo76}, the irreducible representations of $G'$
with non-zero $I'$-fixed vectors correspond to the simple modules of
$\cH(G'/\!/I')$. By Proposition~\ref{prop:borel}(d), the irreducible
genuine representations of $\Gbar$ with non-zero $B^*$-fixed vectors
correspond to the simple modules of the genuine Iwahori--Hecke
algebra. Savin's isomorphism therefore gives a bijection between these
two sets of representations, which Savin calls a local Shimura
correspondence. We do not compare it with Waldspurger's correspondence
\cite{Wa80} for $n=2$, which is defined through the theta
correspondence and depends on an additive character. Three features of
the isomorphism of Theorem~\ref{thm:savinsl2} are worth recording.
\begin{enumerate}[leftmargin=2em,label=(\roman*)]
\item It maps the idempotent $(1+X)/(1+q)$ to $(1+f'_{[s_1]})/(1+q)$,
the identity element of the spherical Hecke algebra of $G'$. So it
restricts to an isomorphism from the genuine spherical Hecke algebra of
$\Gbar$ onto the spherical Hecke algebra of $G'$
(Corollary~\ref{cor:spherical0}). The genuine representations with a
vector fixed by the lift of $K=SL(2,\cO)$ correspond to the unramified
representations of $G'$.
\item It preserves the trace $\phi\mapsto\phi(1)$ and the involution
$\phi^*(x)=\overline{\phi(x^{-1})}$ (Remark~\ref{vrem:plancherel}).
\item For $GL(2)$ and $n=2$, Witthaus \cite[Lemma~6.3]{Wi22} shows that
the corresponding bijection matches genuine unramified principal series
with the unramified principal series whose Satake parameters are
obtained by squaring. We have not written out the analog for $SL(2)$.
\end{enumerate}

\begin{vremark}
\label{vrem:plancherel}
\emph{Not checked in detail.} On the Iwahori-spherical part of the
tempered dual, the trace and the involution determine the Plancherel
measure. So one expects the bijection of Savin's isomorphism to
preserve Plancherel measures and formal degrees, with Haar measures
normalized by $\vol(B)=\vol(I')=1$. We have not checked the details.
\end{vremark}

\subsubsection{Reduction modulo \texorpdfstring{$p$}{p}}
When $n\geq 2$ the constants $c_0$ and $c_1$ are divisible by $q$. Over
a field of characteristic $p$ that contains the $n$-th roots of unity,
the genuine relations therefore become $X^2=-X$ and $Y^2=0$, whereas
Proposition~\ref{prop:IM} gives $f_{s_2}^2=-f_{s_2}$. So modulo $p$ no
isomorphism sends $Y$ to a multiple of $f_{s_2}$, and Savin's
isomorphism, which multiplies by powers of $q^{1/2}$, has no direct
analog. Witthaus \cite[Proposition~3.33]{Wi22} finds the same
phenomenon for the pro-$p$ Iwahori--Hecke algebra of the double cover
of $GL(2,\mathbb{Q}_p)$. There an element supported on the double
coset of an antidiagonal element squares to zero, while the element
attached to $s$ satisfies $S^2=-S$. He uses that algebra to classify the
irreducible genuine representations of the cover modulo $p$. Peskin
\cite{Pe14} studied the double cover of $SL(2,\mathbb{Q}_p)$ modulo $p$
through spherical Hecke algebras.

\subsubsection{What this says about the comparison with \texorpdfstring{$SL(2)$}{SL(2)}}
At Iwahori level the picture is complete. For even $n$ the genuine
algebra is of $PGL(2)$ type. Its generator $Y$ has eigenvalues
$\pm\sqrt{c_0}$, whose ratio is $-1$, whereas $f_{s_1}$ and $f_{s_2}$
have eigenvalues $q$ and $-1$. So no isomorphism onto $\cH(G/\!/B)$
sends $Y$ to a multiple of $f_{s_1}$ or of $f_{s_2}$
(Corollary~\ref{cor:wronggroup}). For odd $n$ the partner is $SL(2)$
itself, but the isomorphism sends the genuine function on
$\overline{Bt_0^{nk}B}$ to a multiple of $f_{t_0^k}$, and the double
cosets $Bt_0^mB$ with $n\nmid m$ carry no genuine functions at all. So
for odd $n$ the group $SL(2)$ is right and the map
$f_g\mapsto\fbar_{\tilde g}$ is wrong, while for even $n$ the group
itself is wrong. As abstract algebras all of these are isomorphic to the
group algebra of the infinite dihedral group
(Remark~\ref{rem:abstract}). The difference lies in the extra structure
that Savin's isomorphism preserves: the trace, the involution, the
finite part spanned by $1$ and $X$, and the double cosets up to the
rescaling $\Psi$. At deep level we know of no analog of Savin's theorem
(Question~\ref{q:savin}).

\begin{example}
\label{ex:savinbg}
In the running example $F=\mathbb{Q}_7$, so $q=7$ and
$(-1)^{(q-1)/2}=-1$. For $n=2$ we have $r=s_2$, $X^2=6X+7$ and
$Y^2=-7$. Savin's isomorphism sends $Y$ to $i\sqrt7\,f'_{[t]}$. For
$n=3$ we have $r=t_0^{-3}s_1$, of length $5$, and $Y^2=294\,Y+16807$.
Savin's isomorphism sends $Y$ to $49f_{s_2}$. The genuine function on
$\overline{Bt_0^{-3}B}$, of volume $7^6$, goes to $49$ times a constant
of absolute value one times $f_{t_0^{-1}}$, of volume $7^2$, and
$49=(7^6/7^2)^{1/2}$.
\end{example}

%======================================================================
\part{General structure and the case \texorpdfstring{$n=2$}{n=2}}

Part~I follows the approach of the original paper. Its first three
sections hold for every $n$, and Part~II uses them.
Section~\ref{section:n2} then specializes to $n=2$.

\section{The structure of \texorpdfstring{$\cH(G/\!/B_k)$}{H(G//Bk)}}
\label{section:sl2}

This section describes the linear algebra $\Hlin$ by generators and
relations, following the method of Howe and Moy \cite{HM85}. The
relations found here are the ones we try to match in the genuine
algebra. Throughout this section $k\geq 1$, Haar measure on $G$ is
normalized by $\vol(B_k)=1$, and $\chi_S$ is the characteristic
function of a subset $S\subset G$.

\begin{proposition}
\label{prop:main}
Let $k\geq 1$. The algebra $\Hlin$ is generated by the functions $f_g$
with $g\in\{s_1,s_2\}\cup B$. These are finitely many functions, since
$f_g$ depends only on $B_kgB_k$. Moreover
$f_{t_0}=f_{-1}*f_{s_1}*f_{s_2}$ and
$f_{t_0^{-1}}=f_{-1}*f_{s_2}*f_{s_1}$. The following relations hold.

\smallskip\noindent
\textup{(A)} Relations for the $s_i$.
\begin{itemize}[leftmargin=3em]
\item[(i)] We have
\begin{align*}
f_{s_1}*f_{s_1}&=q\,\chi_{-\Gamma_k}
=q\sum_{x\in\varpi^k\cO/\varpi^{k+1}\cO}f_{-u_-(x)},\\
f_{s_2}*f_{s_2}&=q\,\chi_{-t\Gamma_kt^{-1}}
=q\sum_{x\in\varpi^{k-1}\cO/\varpi^{k}\cO}f_{-u_+(x)}.
\end{align*}
\item[(ii)] $f_{t_0}*f_{s_1}=f_{s_1}*f_{t_0^{-1}}=f_{t_0s_1}$.
\end{itemize}

\noindent
\textup{(B)} Relations for elements of $B$.
\begin{itemize}[leftmargin=3em]
\item[(i)] $f_1$ is the identity of $\Hlin$.
\item[(ii)] $f_x*f_y=f_{xy}$ for $x,y\in B$. Thus the $f_x$ with $x\in
B$ span a copy of the group algebra $\C[B/B_k]$.
\end{itemize}

\noindent
\textup{(C)} Mixed relations.
\begin{itemize}[leftmargin=3em]
\item[(i)] $f_g*f_x=f_{gx}$ and $f_x*f_g=f_{xg}$ for $x\in B$ and
$g\in G$. In particular $f_{g}*f_x=f_{gxg^{-1}}*f_{g}$ for
$x\in B\cap g^{-1}Bg$.
\item[(ii)] $f_{s_1}*f_{u_+(x)}*f_{s_1}=q\,f_{-u_-(-x)}$ for
$x\in\cO^\times$.
\item[(iii)] $f_{s_2}*f_{u_-(\varpi x)}*f_{s_2}
=q\,f_{-u_+(-x/\varpi)}$ for $x\in\cO^\times$.
\item[(iv)] $f_{t_0}*f_{t_0}=f_{t_0^2}$. More generally
$f_{w}*f_{w'}=f_{ww'}$ whenever $w,w'\in \pm W_a$ and
$\ell(w)+\ell(w')=\ell(ww')$.
\end{itemize}

\noindent
The right-hand sides of \textup{C(ii)} and \textup{C(iii)} are
expressed in terms of the generators by means of the identity
\begin{equation}
\label{eq:bruhat}
s_1u_+(x)s_1=-u_-(-x)=-\,u_+(-x^{-1})\,s_1\,h(-x)\,u_+(-x^{-1}),
\qquad x\in F^\times.
\end{equation}
It gives $f_{s_1}*f_{u_+(x)}*f_{s_1}
=q\,f_{-u_+(-x^{-1})}*f_{s_1}*f_{h(-x)u_+(-x^{-1})}$. Conjugation by
$t$ gives the analog for C(iii).
\end{proposition}

\begin{proof}
\emph{Generation.} If $g=b_1wb_2$ with $b_1,b_2\in B$ and $w\in W_a$,
then $f_g=f_{b_1}*f_w*f_{b_2}$ by C(i). By C(iv) and
Lemma~\ref{lemma:length}, for $m\geq 0$ we have
$f_{t_0^m}=f_{t_0}^{*m}$ and $f_{t_0^ms_1}=f_{t_0}^{*m}*f_{s_1}$. For
$m\geq 1$ we have $f_{t_0^{-m}}=f_{t_0^{-1}}^{*m}$ and
$f_{t_0^{-m}s_1}=f_{t_0^{-1}}^{*(m-1)}*f_{s_2}$, because
$t_0^{-1}s_1=s_2$. Finally $t_0=s_1^{-1}s_2=(-1)s_1s_2$ and
$t_0^{-1}=(-1)s_2s_1$, with
$\ell(s_1s_2)=\ell(s_2s_1)=2=\ell(s_1)+\ell(s_2)$. So the formulas for
$f_{t_0^{\pm1}}$ follow from C(iv) and C(i).

\emph{\textup{B(i), B(ii), C(i)}.} B(i) is
Lemma~\ref{lemma:general}(a). Since $B$ normalizes $B_k$, for $x\in B$
we have $B_kxB_k=xB_k$ and $B_kgxB_k=B_kgB_kx$. Hence
$\vol(B_kxB_k)=1$ and $\vol(B_kgxB_k)=\vol(B_kgB_k)$, and
Lemma~\ref{lemma:general}(d) applies.

\emph{\textup{C(iv), A(ii)}.} These follow from
Lemmas~\ref{lemma:length} and \ref{lemma:general}(d). For A(ii) note
that $t_0s_1=s_1t_0^{-1}$ and $\ell(t_0s_1)=3$.

\emph{\textup{A(i)}.} Let $s=s_1$ and $Q=\Gamma_k$. By
Lemma~\ref{lemma:subgroups}(d), $f_s*f_s$ is supported on
\[
B_ksB_ksB_k=B_ks^2(s^{-1}B_ks)B_k=-B_kQ=-Q.
\]
Let $y=s^2\gamma$ with $\gamma\in Q$, and let $v$ run over the $q$
cosets of $B_k/(B_k\cap sB_ks^{-1})$. By
Lemma~\ref{lemma:general}(b),
\[
(f_s*f_s)(y)=\sum_v f_s(s^{-1}v^{-1}s^2\gamma)
=\sum_v f_s(s\,v^{-1}\gamma),
\]
because $s^2$ is central. As $v^{-1}\gamma\in Q=(s^{-1}B_ks)B_k$, we
have $sv^{-1}\gamma\in B_ksB_k$. So every term is $1$ and
$f_s*f_s=q\chi_{-Q}$. Since $B_k$ is normal in $Q$ and
$Q=\bigsqcup_x u_-(x)B_k$, this is the stated sum. The case of $s_2$
follows by applying the automorphism $f\mapsto f\circ\Ad(t)^{-1}$ of
$\Hlin$, which sends $f_g$ to $f_{tgt^{-1}}$
(Lemma~\ref{lemma:subgroups}(a)). Here we use
$tu_-(x)t^{-1}=u_+(-x/\varpi)$.

\emph{\textup{C(ii)}.} Let $s=s_1$ and $u=u_+(x)$ with $x\in\cO$. By
C(i), $f_s*f_u*f_s=f_{su}*f_s$, and $B_ksuB_k=\bigsqcup_v vsuB_k$ with
$v$ as before. The support lies in
$suB_ksB_k=sus\,(s^{-1}B_ks)B_k=sus\,\Gamma_k$. For
$y=sus\gamma$ with $\gamma\in\Gamma_k$ we get
\[
(f_{su}*f_s)(y)=\sum_vf_s\bigl((su)^{-1}v^{-1}(su)\,s\gamma\bigr).
\]
Since $su\in K$ normalizes $\Gamma_k$, the element
$\gamma'=(su)^{-1}v^{-1}(su)$ lies in $\Gamma_k$. Then
$\gamma's\gamma=s(s^{-1}\gamma's)\gamma\in s\Gamma_k\subset B_ksB_k$.
Hence $f_{su}*f_s=q\,\chi_{sus\Gamma_k}$. By (\ref{eq:bruhat}),
$sus=-u_-(-x)$. If $x$ is a unit, then
$-u_-(-x)\Gamma_k=B_k(-u_-(-x))B_k$ by
Lemma~\ref{lemma:subgroups}(e).

\emph{\textup{C(iii)}.} Apply $\Ad(t)$ to C(ii) with $x$ replaced by
$-x$. Here we use $tu_+(-x)t^{-1}=u_-(\varpi x)$ and
$tu_-(x)t^{-1}=u_+(-x/\varpi)$.

The identity (\ref{eq:bruhat}) is a direct matrix computation. The
displayed formula after it follows from (\ref{eq:bruhat}) and C(i).
\end{proof}

The next example shows that $f_{t_0}$ is not invertible up to a scalar
in $\Hlin$. This is a basic difference between positive level and the
affine Weyl group, where $t_0$ is invertible.

\begin{example}
\label{ex:t0}
Let
\[
R=(t_0B_kt_0^{-1})B_k
=G\cap\mat{1+\varpi^k\cO}{\varpi^k\cO}{\varpi^{k-1}\cO}{1+\varpi^k\cO},
\]
a group containing $B_k$ with index $q^2$. Then
$B_kt_0^{-1}B_k=t_0^{-1}R$. By Lemma~\ref{lemma:general}(b),
\[
(f_{t_0}*f_{t_0^{-1}})(y)=\sum_{v\in B_k/(B_k\cap t_0B_kt_0^{-1})}
\chi_{t_0^{-1}R}(t_0^{-1}v^{-1}y)=q^2\chi_R(y).
\]
Thus $f_{t_0}*f_{t_0^{-1}}=q^2\chi_R$, which is not a multiple of
$f_1$. Here $\ell(t_0)+\ell(t_0^{-1})=4\neq 0=\ell(1)$, so C(iv) does
not apply. Similarly $f_{s_1}*f_{t_0}\neq f_{s_1t_0}=f_{s_2}$, since the
two sides have integrals $q^3$ and $q$. In the running example with
$k=1$, $f_{t_0}*f_{t_0^{-1}}=49\,\chi_R$, and $R/B_1$ has $49$
elements.
\end{example}

\begin{remark}
The analog of Proposition~\ref{prop:main} for $GL(n)$ is
\cite[Chapter~3, Theorem~2.1]{HM85}. There it is shown in addition that
the relations form a defining set. We have not carried out the
corresponding argument for the list (A)--(C). The results of this
paper do not use it.
\end{remark}

\begin{question}
\label{q:defining}
Do the relations \textup{(A)--(C)} of Proposition~\ref{prop:main},
together with the expressions of the $f_g$ on their right-hand sides
in terms of the generators, form a defining set of relations for
$\Hlin$?
\end{question}

%======================================================================
\section{Splittings and genuine functions}
\label{section:genuine}

This section constructs the genuine counterpart of the basis
$\{f_g\}$. The point is that $B_k$ must be lifted to $\Gbar$ in a way
compatible with conjugation, and that a double coset carries a genuine
function only if a commutation condition holds. At deep level both
issues are resolved.

\subsection{Canonical splittings}
This subsection shows that sufficiently small open subgroups of $G$
have a distinguished lift to $\Gbar$, compatible with conjugation by
every element of $\Gbar$. These lifts make all later constructions
independent of the choice of cocycle.

Let $e$ be the ramification index of $F$ over $\mathbb{Q}_p$, and
$\ord$ the normalized valuation of $F$. For an integer $j>e/(p-1)$ the
exponential series converges on $\varpi^jM_2(\cO)$ and defines a
bijection onto $1+\varpi^jM_2(\cO)$ whose inverse is the logarithm
series. Let $g\in GL(2,F)$. Conjugation by $g$ is continuous on
$M_2(F)$ and commutes with the partial sums of both series. Hence the
exponential series commutes with conjugation by $g$ whenever $X$ and
$gXg^{-1}$ lie in $\varpi^jM_2(\cO)$, and the logarithm series commutes
with conjugation by $g$ whenever $b$ and $gbg^{-1}$ lie in
$1+\varpi^jM_2(\cO)$. This holds even when $g$ does not lie in
$GL(2,\cO)$. Moreover $\det\exp(X)=\exp(\tr X)$.

\begin{definition}
\label{def:deep}
Fix an integer $j_0>e/(p-1)$ such that $\pr$ splits over
$\Gamma_{j_0}$. Such $j_0$ exists by Remark~\ref{vrem:kubota}. Let
$j_1=j_0+\ord(n)$. Let $\cP$ be the following family of open subgroups
of $G$.
\begin{itemize}[leftmargin=2em]
\item If $p\nmid n$, $\cP$ consists of all open pro-$p$ subgroups of
$G$.
\item If $p\mid n$, $\cP$ consists of all open subgroups of
$\Gamma_{j_1}$.
\end{itemize}
The level $k\geq 1$ is called \textbf{deep} if $B_k$, $\Gamma_k$ and
$t\Gamma_kt^{-1}$ belong to $\cP$. Equivalently, $k\geq 1$ when
$p\nmid n$, and $k\geq j_1+1$ when $p\mid n$.
\end{definition}

The family $\cP$ is closed under intersections and under passage to
open subgroups. The equivalence in the definition holds because
$t\Gamma_kt^{-1}=U^+_{k-1}B_k$ is contained in $\Gamma_{k-1}$ and is
pro-$p$ for $k\geq 1$.

\begin{example}
\label{ex:deep}
In the running example $p=7$ does not divide $n=3$, so every $k\geq 1$
is deep. For the classical double cover over $F=\mathbb{Q}_2$ we have
$p\mid n$, $e=1$ and $\ord(2)=1$. Then $j_0\geq 2$, and the deep levels
are $k\geq j_0+2$.
\end{example}

\begin{lemma}[Canonical splitting]
\label{lemma:canonical}
For each $P\in\cP$ there is a splitting $s_P:P\to\Gbar$ with the
following properties.
\begin{enumerate}[leftmargin=2em,label=(\alph*)]
\item If $P'\subset P$ are in $\cP$, then $s_P|_{P'}=s_{P'}$. Hence
$\sigma(x)=s_P(x)$ is a well-defined lift of each
$x\in\bigcup_{P\in\cP}P$.
\item If $x\in P$, $g\in G$ and $gxg^{-1}\in P'$ with $P,P'\in\cP$,
then $\tilde g\,\sigma(x)\,\tilde g^{-1}=\sigma(gxg^{-1})$ for every
lift $\tilde g$ of $g$.
\item $\sigma(u_\pm(x))=(u_\pm(x),1)$ whenever $u_\pm(x)$ lies in a
member of $\cP$.
\item If $p\nmid n$, then $s_P$ is the unique splitting of $P$.
\end{enumerate}
\end{lemma}

\begin{proof}
\emph{Case $p\nmid n$.} Let $P$ be an open pro-$p$ subgroup. Then
$\overline{P}$ is a profinite group and a central extension of the
pro-$p$ group $P$ by the group $\mu_n$, whose order is prime to $p$.
Hence $\overline{P}$ has a unique Sylow pro-$p$ subgroup $S_P$, we have
$\overline{P}=\mu_n\times S_P$, and $\pr:S_P\to P$ is an isomorphism.
Let $s_P$ be its inverse. An element of $\overline{P}$ lies in $S_P$ if
and only if it generates a pro-$p$ subgroup. This gives (a). For (b),
the element $\tilde g\sigma(x)\tilde g^{-1}$ lies in the pro-$p$ group
$\tilde gS_P\tilde g^{-1}$ and in $\overline{P'}$. Hence it lies in
$S_{P'}$, and it projects to $gxg^{-1}$. For (d), two splittings of $P$
differ by a homomorphism $P\to\mu_n$. Such a homomorphism is trivial,
because every element of a pro-$p$ group is an $n$-th power when
$p\nmid n$. For (c), the restrictions of $\sigma$ and of
$u\mapsto(u,1)$ to the pro-$p$ group $U^\pm\cap P$ are both splittings
by Lemma~\ref{lemma:cover}(b). They differ by a homomorphism from a
pro-$p$ group to $\mu_n$, which is trivial.

\emph{Case $p\mid n$.} Let $s'$ be a splitting of $\Gamma_{j_0}$. For
$b\in\Gamma_{j_1}$ put $X=\log b\in\varpi^{j_1}M_2(\cO)$. Then
$\tr X=0$ and $X/n\in\varpi^{j_0}M_2(\cO)$, so
$c=\exp(X/n)\in\Gamma_{j_0}$ satisfies $c^n=b$. Therefore
\[
s'(b)=s'(c)^n=\tilde c^{\,n}\qquad\mbox{for every lift $\tilde c$ of
$c$,}
\]
because $\mu_n$ is central of exponent $n$. Define $s_P=s'|_P$ for
$P\subset\Gamma_{j_1}$. Part (a) is clear, and the formula shows that
$s_P$ does not depend on $s'$. For (b), let $b=x$ and $b'=gxg^{-1}$,
both in $\Gamma_{j_1}$, with $X=\log b$. Then $\log b'=gXg^{-1}$, so
$c'=\exp(gXg^{-1}/n)=gcg^{-1}$. The element
$\tilde g\,s'(c)\,\tilde g^{-1}$ is a lift of $c'$. Hence
\[
\tilde g\,\sigma(b)\,\tilde g^{-1}
=\bigl(\tilde g\,s'(c)\,\tilde g^{-1}\bigr)^n=\sigma(b').
\]
For (c), if $b=u_\pm(x)$, then $c=u_\pm(x/n)$ and
$(c,1)^n=(u_\pm(x),1)$ by Lemma~\ref{lemma:cover}(b).
\end{proof}

\begin{example}
\label{ex:canonical}
Let $k$ be deep and $v\in 1+\varpi^k\cO$. Apply
Lemma~\ref{lemma:canonical}(b) with $x=h(v)\in B_k$ and $g=t_0$, which
commutes with $x$. It gives
$\tilde t_0\sigma(h(v))\tilde t_0^{-1}=\sigma(h(v))$, so
$(v,\varpi)_n^2=1$ by Lemma~\ref{lemma:cover}(a). This is consistent
with the fact that, at deep level, every element of $1+\varpi^k\cO$ is
an $n$-th power. For a unit $u$ that is not an $n$-th power no such
statement holds. In the running example $u=3$ is such a unit.
\end{example}

\subsection{Genuine Hecke algebras}
This subsection defines the genuine Hecke algebra of a lifted open
compact subgroup. It identifies the double cosets that carry genuine
functions and proves the genuine analog of the multiplication rule of
Lemma~\ref{lemma:general}(d).

A function $f$ on $\Gbar$ is called \textbf{genuine} if
$f(\zeta\tilde x)=\zeta f(\tilde x)$ for all $\tilde x\in\Gbar$ and
$\zeta\in\mu_n$. Here we identify $\mu_n$ with the complex $n$-th roots
of unity by a fixed isomorphism. Let $K_0\subset G$ be an open compact
subgroup with a splitting $s$, and put $K_0^*=s(K_0)$. Haar measure on
$G$ is normalized by $\vol(K_0)=1$, the measure on $\mu_n$ has total
mass $1$, and $\Gbar$ carries the product measure. Then
$\vol(\overline{K_0})=1$.

\begin{definition}
\label{def:genuineHecke}
Let $\cH(\Gbar/\!/K_0^*)_{\gen}$ be the space of functions on $\Gbar$
that are genuine, locally constant, compactly supported, and
bi-$K_0^*$-invariant. With the convolution
\[
(\phi_1*\phi_2)(\tilde y)=\int_{\Gbar}\phi_1(\tilde x)
\phi_2(\tilde x^{-1}\tilde y)\,d\tilde x
=\int_G\phi_1(\tilde x)\phi_2(\tilde x^{-1}\tilde y)\,dx
\]
it is the \textbf{genuine generalized Hecke algebra} of $\Gbar$ with
respect to $K_0^*$. In the second integral $\tilde x$ is any lift of
$x$, and the integrand does not depend on the choice. The algebra
depends on the splitting $s$.
\end{definition}

\begin{example}
\label{ex:identity}
The identity element of $\cH(\Gbar/\!/K_0^*)_{\gen}$ is the function
$e_{K_0^*}$ supported on $\overline{K_0}$ with
$e_{K_0^*}(\zeta s(x))=\zeta$. In the notation introduced below,
$e_{K_0^*}=\fbar_{(1,1)}$.
\end{example}

\begin{definition}
\label{def:admissible}
An element $g\in G$ is called \textbf{$K_0^*$-admissible} if, for some
(equivalently, any) lift $\tilde g$,
\begin{equation}
\label{eq:admissible}
s(gxg^{-1})=\tilde g\,s(x)\,\tilde g^{-1}\qquad
\mbox{for all }x\in K_0\cap g^{-1}K_0g.
\end{equation}
\end{definition}

\begin{example}
Every element of $K_0$ is $K_0^*$-admissible, because $s$ is a
homomorphism. At Iwahori level with $n=3$ the element $t_0$ is not
admissible (Proposition~\ref{prop:level0}). At deep level every element
is admissible (Proposition~\ref{prop:admissible}).
\end{example}

\begin{lemma}
\label{lemma:support}
Let $g\in G$ with lift $\tilde g$.
\begin{enumerate}[leftmargin=2em,label=(\alph*)]
\item The space of genuine bi-$K_0^*$-invariant functions supported on
$\overline{K_0gK_0}$ has dimension $1$ if $g$ is $K_0^*$-admissible,
and dimension $0$ otherwise. Admissibility depends only on the double
coset $K_0gK_0$.
\item If $g$ is admissible, let $\fbar_{\tilde g}$ be the unique such
function with $\fbar_{\tilde g}(\tilde g)=1$. Thus
$\fbar_{\tilde g}(\zeta\,s(x_1)\,\tilde g\,s(x_2))=\zeta$. Then
$\fbar_{\zeta\tilde g}=\zeta^{-1}\fbar_{\tilde g}$, and
$\fbar_{s(x_1)\tilde gs(x_2)}=\fbar_{\tilde g}$ for $x_1,x_2\in K_0$.
\item Choose one lift $\tilde g$ of a representative of each
admissible double coset. The resulting functions $\fbar_{\tilde g}$
form a basis of $\cH(\Gbar/\!/K_0^*)_{\gen}$.
\item For admissible $g$ and $\phi\in\cH(\Gbar/\!/K_0^*)_{\gen}$,
\[
(\fbar_{\tilde g}*\phi)(\tilde y)
=\sum_{v\in K_0/(K_0\cap gK_0g^{-1})}
\phi\bigl(\tilde g^{-1}s(v)^{-1}\tilde y\bigr).
\]
\end{enumerate}
\end{lemma}

\begin{proof}
(a) A genuine bi-$K_0^*$-invariant function $f$ supported on
$\overline{K_0gK_0}$ is determined by $f(\tilde g)$. A function with
$f(\tilde g)=1$ exists if and only if $s(x_1)\tilde gs(x_2)=\zeta\tilde
g$ implies $\zeta=1$. The relation $x_1gx_2=g$ means $x_1=gxg^{-1}$
with $x=x_2^{-1}\in K_0\cap g^{-1}K_0g$. In that case
\[
s(x_1)\tilde gs(x_2)=s(gxg^{-1})\,\bigl(\tilde gs(x)\tilde
g^{-1}\bigr)^{-1}\,\tilde g .
\]
The rest of (a), and parts (b) and (c), follow at once. For (d), write
$K_0gK_0=\bigsqcup_v vgK_0$. On $vgK_0$ we may take
$\tilde x=s(v)\tilde gs(x_2)$. Then $\fbar_{\tilde g}(\tilde x)=1$ and
$\phi(\tilde x^{-1}\tilde y)=\phi(\tilde g^{-1}s(v)^{-1}\tilde y)$, and
each coset $vgK_0$ has volume $1$.
\end{proof}

The following lemma is the genuine analog of
Lemma~\ref{lemma:general}(d).

\begin{lemma}
\label{lemma:genuine}
Let $g_1,g_2\in G$ be $K_0^*$-admissible, and suppose that
\[
\vol(K_0g_1K_0)\vol(K_0g_2K_0)=\vol(K_0g_1g_2K_0).
\]
Then $g_1g_2$ is $K_0^*$-admissible and
$\fbar_{\tilde g_1}*\fbar_{\tilde g_2}=\fbar_{\tilde g_1\tilde g_2}$
for all lifts $\tilde g_1$ and $\tilde g_2$. In terms of the cocycle,
\[
\fbar_{(g_1,1)}*\fbar_{(g_2,1)}=\beta(g_1,g_2)^{-1}\fbar_{(g_1g_2,1)}.
\]
\end{lemma}

\begin{proof}
The absolute value of $\fbar_{\tilde g_i}$ is $f_{g_i}\circ\pr$. So
$\fbar_{\tilde g_1}*\fbar_{\tilde g_2}$ is supported on the inverse
image of the support of $f_{g_1}*f_{g_2}$, which is $K_0g_1g_2K_0$ by
Lemma~\ref{lemma:general}(d). The convolution is genuine and
bi-$K_0^*$-invariant. We compute its value at
$\tilde g_1\tilde g_2$:
\[
(\fbar_{\tilde g_1}*\fbar_{\tilde g_2})(\tilde g_1\tilde g_2)
=\int_G\fbar_{\tilde g_1}(\tilde x)\,
\fbar_{\tilde g_2}(\tilde x^{-1}\tilde g_1\tilde g_2)\,dx .
\]
By (\ref{eq:unique}) the integrand vanishes unless $x\in g_1K_0$. For
$\tilde x=\tilde g_1s(x_2)$ the integrand is
\[
\fbar_{\tilde g_1}(\tilde g_1s(x_2))\,
\fbar_{\tilde g_2}(s(x_2)^{-1}\tilde g_2)=1 .
\]
Since $\vol(g_1K_0)=1$, the value is $1$. So the convolution is a
non-zero genuine bi-$K_0^*$-invariant function supported over
$K_0g_1g_2K_0$. By Lemma~\ref{lemma:support}(a), $g_1g_2$ is
admissible, and by Lemma~\ref{lemma:support}(b) the convolution equals
$\fbar_{\tilde g_1\tilde g_2}$. The last formula follows from the
identity $(g_1,1)(g_2,1)=\beta(g_1,g_2)(g_1g_2,1)$ and
Lemma~\ref{lemma:support}(b).
\end{proof}

\begin{remark}
\label{rem:cocycle}
These constructions depend only on the group $\Gbar$. Replace $\beta$
by a cohomologous cocycle
$\beta'(g,h)=\beta(g,h)\,s_0(g)^{-1}s_0(h)^{-1}s_0(gh)$, where
$s_0:G\to\mu_n$ is any function with $s_0(1)=1$. Then
$(g,\zeta)\mapsto(g,s_0(g)^{-1}\zeta)$ is an isomorphism from the group
defined by $\beta'$ to the group defined by $\beta$, and it carries
genuine Hecke algebras to genuine Hecke algebras. In contrast, the
function $(x,\zeta)\mapsto\zeta\,\chi_{K_0gK_0}(x)$ depends on the
cocycle and need not be bi-$K_0^*$-invariant. This is why we work with
lifts $\tilde g$ and not with the elements $(g,1)$.
\end{remark}

\subsection{Deep level}
This subsection specializes to $K_0=B_k$ at deep level. There the
canonical lift makes every element of $G$ admissible, so every double
coset of $B_k$ carries a genuine function, unique up to a root of
unity. The result is a basis of the genuine Hecke algebra indexed in
the same way as the basis of $\Hlin$, and a linear isomorphism between
the two algebras that preserves supports.
Sections~\ref{section:n2} and \ref{section:torus} ask whether this map
can be made multiplicative.

\begin{proposition}
\label{prop:admissible}
Let $k$ be deep, and let $B_k^*=\sigma(B_k)$ be the canonical lift of
$B_k$. Then every $g\in G$ is $B_k^*$-admissible. Consequently, for any
choice of lifts $\tilde g$ of representatives of the double cosets
$B_kgB_k$, the functions $\fbar_{\tilde g}$ form a basis of $\Hgen$,
and
\[
\Phi:\Hlin\to\Hgen,\qquad f_g\mapsto\fbar_{\tilde g},
\]
is an isomorphism of vector spaces with $|\Phi(f)|=|f|\circ\pr$.
\end{proposition}

\begin{proof}
Condition (\ref{eq:admissible}) for $K_0=B_k$ and $s=\sigma|_{B_k}$ is
Lemma~\ref{lemma:canonical}(b) with $P=P'=B_k$. The rest follows from
Lemma~\ref{lemma:support}.
\end{proof}

The isomorphism $\Phi$ is natural up to roots of unity. Changing the
lift $\tilde g$ multiplies $\fbar_{\tilde g}$ by an element of
$\mu_n$, so $\Phi$ depends on the choice of lifts only through one root
of unity for each double coset. It is not an algebra homomorphism in
general. The rest of the paper asks whether some modification of
$\Phi$ is one, and the following notion describes the modifications we
consider.

\begin{definition}
\label{def:supportpreserving}
Let $K_0$ be an open compact subgroup of $G$ with a splitting $s$, and
put $K_0^*=s(K_0)$. A linear map
$\Psi:\cH(G/\!/K_0)\to\cH(\Gbar/\!/K_0^*)_{\gen}$ is
\textbf{support-preserving} if, for every double coset $K_0gK_0$, the
function $\Psi(f_g)$ vanishes outside $\pr^{-1}(K_0gK_0)$.
\end{definition}

By Lemma~\ref{lemma:support}(a), $\Psi$ is support-preserving if and
only if $\Psi(f_g)=c_g\fbar_{\tilde g}$ for every $K_0^*$-admissible
$g$, with complex constants $c_g$ that may be zero, and $\Psi(f_g)=0$
for the other $g$. At deep level every $g$ is admissible, and these
are exactly the maps $\Phi\circ D_c$, where $D_c(f_g)=c_gf_g$ rescales the basis, and
such a map is bijective if and only if every $c_g$ is non-zero. So
$\Phi$ is the support-preserving isomorphism with all constants equal to
$1$, and every support-preserving map is $\Phi$ up to rescaling the
basis. Theorem~\ref{thm:unimodular} and Corollary~\ref{corollary:main}
show that no support-preserving map with all constants non-zero is an
algebra homomorphism, for $n\geq 2$ when the constants have absolute
value one, and for $n\geq 3$ in general. Savin's isomorphism
(Theorem~\ref{thm:savinsl2}) is of a different kind: for odd $n$ its
inverse sends $f_{t_0^k}$ to a multiple of $\fbar_{(t_0^{nk},1)}$, so
it moves supports.

\begin{remark}[No other choice of genuine functions]
\label{rem:nochoice}
By Lemma~\ref{lemma:support}(a), every genuine bi-$B_k^*$-invariant
function supported over a single double coset $B_kgB_k$ is a multiple
of $\fbar_{\tilde g}$. A different lift of $g$ only changes the multiple
by a root of unity. So the only freedom in a map that keeps each basis
function on its own double coset is a choice of constants $c_g$.
Theorem~\ref{thm:unimodular} and Corollary~\ref{corollary:main} treat
all such choices, except constants of varying absolute value when
$n=2$. Maps that move $f_g$ to other double cosets are not of
this form. Savin's isomorphism at Iwahori level is of that kind
(Theorem~\ref{thm:savinsl2}).
\end{remark}

\begin{example}
In the running example, take $k=1$. The double cosets $B_1gB_1$ are
indexed by $B/B_1\times(\pm W_a)\times B/B_1$ modulo the obvious
identifications. For each of them, Proposition~\ref{prop:admissible}
gives a function $\fbar_{\tilde g}$ on the cover of $SL(2,\mathbb{Q}_7)$
with values in $\{0\}\cup\mu_3$.
\end{example}

%======================================================================
\section{The structure of
\texorpdfstring{$\cH(\Gbar/\!/B_k^*)_{\gen}$}{the genuine Hecke algebra}}
\label{section:metaplectic}

This section transfers the relations of Proposition~\ref{prop:main} to
the genuine algebra, for every $n$. Each relation holds after the
elements of $G$ are replaced by suitable lifts.
Sections~\ref{section:n2} and \ref{section:torus} ask whether one
choice of lifts works for all relations at once. Throughout, $k$ is deep,
$B_k^*=\sigma(B_k)$ and $\vol(B_k)=1$. For $x$ in $\Gamma_k$ or in
$t\Gamma_kt^{-1}$ we use the canonical lift $\sigma(x)$. Other elements
receive arbitrary lifts, denoted by a tilde. Recall that $\tilde
s_i^{\,2}$ is central by Lemma~\ref{lemma:cover}(c).

\begin{proposition}
\label{prop:main2}
Let $k$ be deep. The algebra $\Hgen$ is generated by the functions
$\fbar_{\tilde g}$ with $g\in\{s_1,s_2\}\cup B$. The following
relations hold, for arbitrary lifts unless otherwise stated.

\smallskip\noindent
\textup{($\overline{\mbox{A}}$)}
\begin{itemize}[leftmargin=3em]
\item[(i)] We have
\begin{align*}
\fbar_{\tilde s_1}*\fbar_{\tilde s_1}
&=q\sum_{x\in\varpi^k\cO/\varpi^{k+1}\cO}
\fbar_{\tilde s_1^{\,2}\,(u_-(x),1)},\\
\fbar_{\tilde s_2}*\fbar_{\tilde s_2}
&=q\sum_{x\in\varpi^{k-1}\cO/\varpi^{k}\cO}
\fbar_{\tilde s_2^{\,2}\,(u_+(x),1)}.
\end{align*}
\item[(ii)] $\fbar_{\tilde t_0}*\fbar_{\tilde s_1}
=\fbar_{\tilde t_0\tilde s_1}
=\varepsilon\,\fbar_{\tilde s_1}*\fbar_{\tilde t_0^{-1}}$, where
$\varepsilon\in\mu_n$ is defined by
$\tilde s_1\tilde t_0^{-1}\tilde s_1^{-1}=\varepsilon\,\tilde t_0$.
\end{itemize}

\noindent
\textup{($\overline{\mbox{B}}$)}
\begin{itemize}[leftmargin=3em]
\item[(i)] $\fbar_{(1,1)}$ is the identity of $\Hgen$.
\item[(ii)] $\fbar_{\tilde x}*\fbar_{\tilde y}=\fbar_{\tilde x\tilde y}$
for $x,y\in B$. Thus the $\fbar_{\tilde x}$ with $x\in B$ span a copy
of the genuine part of the group algebra of the finite group
$\overline{B}/B_k^*$. This is a twisted group algebra of $B/B_k$, and
$\fbar_{(x,1)}*\fbar_{(y,1)}=\beta(x,y)^{-1}\fbar_{(xy,1)}$.
\end{itemize}

\noindent
\textup{($\overline{\mbox{C}}$)}
\begin{itemize}[leftmargin=3em]
\item[(i)] $\fbar_{\tilde g}*\fbar_{\tilde x}=\fbar_{\tilde g\tilde x}$
and $\fbar_{\tilde x}*\fbar_{\tilde g}=\fbar_{\tilde x\tilde g}$ for
$x\in B$ and $g\in G$. In particular
$\fbar_{\tilde g}*\fbar_{\tilde x}
=\fbar_{\tilde g\tilde x\tilde g^{-1}}*\fbar_{\tilde g}$ for
$x\in B\cap g^{-1}Bg$.
\item[(ii)] $\fbar_{\tilde s_1}*\fbar_{\tilde u}*\fbar_{\tilde s_1}
=q\,\fbar_{\tilde s_1\tilde u\tilde s_1}$ for $u=u_+(x)$ with
$x\in\cO^\times$.
\item[(iii)] $\fbar_{\tilde s_2}*\fbar_{\tilde u}*\fbar_{\tilde s_2}
=q\,\fbar_{\tilde s_2\tilde u\tilde s_2}$ for $u=u_-(\varpi x)$ with
$x\in\cO^\times$.
\item[(iv)] $\fbar_{\tilde w}*\fbar_{\tilde w'}
=\fbar_{\tilde w\tilde w'}$ whenever $w,w'\in \pm W_a$ and
$\ell(w)+\ell(w')=\ell(ww')$.
\end{itemize}
\end{proposition}

\begin{proof}
Relations ($\overline{\mbox{B}}$), ($\overline{\mbox{C}}$)(i) and
($\overline{\mbox{C}}$)(iv), and the first equality in
($\overline{\mbox{A}}$)(ii), follow from Lemma~\ref{lemma:genuine}
with the volume computations in the proof of
Proposition~\ref{prop:main}. The statement about generators follows by
the argument given there, because every basis element
$\fbar_{\tilde g}$ is a root of unity times a product of generators.
For the second equality in ($\overline{\mbox{A}}$)(ii),
Lemma~\ref{lemma:genuine} gives $\fbar_{\tilde s_1}*\fbar_{\tilde
t_0^{-1}}=\fbar_{\tilde s_1\tilde t_0^{-1}}$. Since $\tilde s_1\tilde
t_0^{-1}=\varepsilon\,\tilde t_0\tilde s_1$, we get
$\fbar_{\tilde s_1\tilde t_0^{-1}}
=\varepsilon^{-1}\fbar_{\tilde t_0\tilde s_1}$ by
Lemma~\ref{lemma:support}(b).

($\overline{\mbox{A}}$)(i). Let $s=s_1$ and $Q=\Gamma_k$. As in the
linear case the convolution is supported over $-Q$. Let $\gamma\in Q$,
and let $v$ run over the $q$ cosets of $B_k/(B_k\cap sB_ks^{-1})$.
Lemma~\ref{lemma:support}(d) and the centrality of $\tilde s^{\,2}$
give
\[
(\fbar_{\tilde s}*\fbar_{\tilde s})(\tilde s^{\,2}\sigma(\gamma))
=\sum_v\fbar_{\tilde s}\bigl(\tilde s\,\sigma(v^{-1}\gamma)\bigr).
\]
Write $v^{-1}\gamma=(s^{-1}bs)b'$ with $b,b'\in B_k$. By
Lemma~\ref{lemma:canonical}(b),
$\sigma(v^{-1}\gamma)=\tilde s^{-1}\sigma(b)\tilde s\,\sigma(b')$. So
$\tilde s\sigma(v^{-1}\gamma)=\sigma(b)\tilde s\sigma(b')$, and each
term equals $1$. Hence $\fbar_{\tilde s}*\fbar_{\tilde s}$ takes the
value $q$ at $\tilde s^{\,2}\sigma(\gamma)$ for every $\gamma\in Q$.
Moreover $Q=\bigsqcup_xu_-(x)B_k$, and $\sigma(u_-(x))=(u_-(x),1)$ by
Lemma~\ref{lemma:canonical}(c). The case of $s_2$ is identical, with
$Q$ replaced by $t\Gamma_kt^{-1}=\bigsqcup_xu_+(x)B_k$, which belongs
to $\cP$ because $k$ is deep.

($\overline{\mbox{C}}$)(ii). Let $s=s_1$. By
($\overline{\mbox{C}}$)(i) the left side is
$\fbar_{\tilde s\tilde u}*\fbar_{\tilde s}$. By the proof of
Proposition~\ref{prop:main}, it is supported over the single double
coset $B_k(sus)B_k$. By Lemma~\ref{lemma:support}(d),
\[
(\fbar_{\tilde s\tilde u}*\fbar_{\tilde s})(\tilde s\tilde u\tilde s)
=\sum_v\fbar_{\tilde s}\bigl((\tilde s\tilde u)^{-1}
\sigma(v)^{-1}(\tilde s\tilde u)\,\tilde s\bigr)
=\sum_v\fbar_{\tilde s}\bigl(\sigma(\gamma'_v)\,\tilde s\bigr),
\]
where $\gamma'_v=(su)^{-1}v^{-1}(su)\in\Gamma_k$, and we used
Lemma~\ref{lemma:canonical}(b). By the same lemma,
$\sigma(\gamma'_v)\tilde s=\tilde s\,\sigma(s^{-1}\gamma'_vs)$, with
$s^{-1}\gamma_v's\in\Gamma_k=(s^{-1}B_ks)B_k$. The argument for
($\overline{\mbox{A}}$)(i) shows that each term is $1$.

($\overline{\mbox{C}}$)(iii) is proved in the same way, with $s_1$,
$K$ and $\Gamma_k$ replaced by $s_2$, $tKt^{-1}$ and $t\Gamma_kt^{-1}$.
\end{proof}

\begin{remark}
In ($\overline{\mbox{C}}$)(ii) the element
$\tilde s_1\tilde u\tilde s_1$ is a lift of
$-u_+(-x^{-1})s_1h(-x)u_+(-x^{-1})$ by (\ref{eq:bruhat}). So the right
side is a root of unity times
$q\,\fbar_{\widetilde{-1}}*\fbar_{\tilde u'}*\fbar_{\tilde s_1}*
\fbar_{\tilde h(-x)}*\fbar_{\tilde u'}$, where $u'=u_+(-x^{-1})$. The
root of unity depends on the lifts and can be computed from $\beta$.
We do not need it.
\end{remark}

\begin{remark}[The twist]
\label{rem:twist}
Proposition~\ref{prop:main2} says that the genuine algebra differs from
$\Hlin$ only by roots of unity. These roots of unity can be organized
as follows. Fix a lift $\tilde g$ of each representative $g$ of a double
coset. For a pair $g_1,g_2$ whose volumes multiply,
Lemma~\ref{lemma:genuine} gives
\[
\fbar_{\tilde g_1}*\fbar_{\tilde g_2}
=\gamma(g_1,g_2)\,\fbar_{\widetilde{g_1g_2}},
\qquad\mbox{where}\quad
\tilde g_1\tilde g_2=\gamma(g_1,g_2)^{-1}\,\widetilde{g_1g_2}.
\]
For the lifts $(g,1)$ we get $\gamma=\beta^{-1}$, and another choice of
lifts changes $\gamma$ by a coboundary. So the genuine algebra is
$\Hlin$ twisted by the class of Kubota's cocycle, seen through the
double cosets of $B_k$. A map $f_g\mapsto c_g\fbar_{\tilde g}$ is
compatible with these products exactly when the constants $c_g$
trivialize $\gamma$ on the pairs that occur. For $n\geq 3$ this already
fails on the torus (Theorem~\ref{thrm:main}). For $n=2$ we do not know
whether it fails on the pairs whose volumes multiply. The obstruction
of Theorem~\ref{thm:unimodular} comes from the other products.
\end{remark}



%======================================================================
\section{The classical case \texorpdfstring{$n=2$}{n=2}}
\label{section:n2}

This section treats the classical metaplectic group, $n=2$, with $p$
odd. It first shows that everything supported on the maximal compact
subgroup $K=SL(2,\cO)$ transfers to the cover without change. It then
computes, from Kubota's cocycle, the signs that appear in the remaining
relations. Finally it shows that no support-preserving map with
constants of absolute value one is an algebra homomorphism. The first
and last results hold more generally, and we state them in their
natural generality.

\subsection{The subalgebra supported on \texorpdfstring{$K$}{K}}
This subsection shows that the part of the Hecke algebra supported on
$K$ behaves exactly as in the linear case. The reason is that the cover
splits over $K$ when $p\nmid n$.

\begin{proposition}
\label{prop:Ksub}
Assume $p\nmid n$ and $k\geq 1$, and let $s$ be a splitting of $K$
(Remark~\ref{vrem:kubota}). Then $s$ agrees with the canonical lift
$\sigma$ on $B_k$. For a function $f$ on $G$ supported on $K$, let
$\fbar$ be the genuine function supported on $\overline{K}$ with
$\fbar(\zeta s(x))=\zeta f(x)$ for $x\in K$ and $\zeta\in\mu_n$. Then
$f\mapsto\fbar$ is an isomorphism of algebras from the subalgebra of
$\Hlin$ of functions supported on $K$ onto the subalgebra of $\Hgen$ of
functions supported on $\overline{K}$. It sends $f_g$ to
$\fbar_{s(g)}$ for $g\in K$.
\end{proposition}

\begin{proof}
The restriction of $s$ to the pro-$p$ group $B_k$ is a splitting, so it
equals $\sigma$ by Lemma~\ref{lemma:canonical}(d). The group
$\overline{K}$ is the direct product of $\mu_n$ and $s(K)$. So a
genuine function supported on $\overline{K}$ is determined by its
composite with $s$, which is a function on $K$. It is
bi-$B_k^*$-invariant exactly when the composite is
bi-$B_k$-invariant, because $B_k^*=s(B_k)$. Let $\phi_1$ and $\phi_2$
be genuine functions supported on $\overline{K}$. Their convolution is
supported on $\overline{K}$. For $y\in K$ the integrand
$\phi_1(\tilde x)\phi_2(\tilde x^{-1}s(y))$ vanishes unless $x\in K$,
and for $x\in K$ we may take $\tilde x=s(x)$. Since
$s(x)^{-1}s(y)=s(x^{-1}y)$, we get
\[
(\phi_1*\phi_2)(s(y))=\int_K(\phi_1\circ s)(x)\,(\phi_2\circ
s)(x^{-1}y)\,dx=\bigl((\phi_1\circ s)*(\phi_2\circ s)\bigr)(y).
\]
So $\phi\mapsto\phi\circ s$ is an algebra isomorphism, and its inverse
is $f\mapsto\fbar$. The image of $f_g$ takes the value $1$ at $s(g)$,
so it equals $\fbar_{s(g)}$.
\end{proof}

\begin{example}
\label{ex:Ksub}
Proposition~\ref{prop:Ksub} applies to the running example for both
$n=2$ and $n=3$, because $7$ divides neither. All relations of
Proposition~\ref{prop:main} among elements of $K$ hold for the
$\fbar_{s(g)}$ without any sign. For instance relation A(i) becomes
$\fbar_{s(s_1)}*\fbar_{s(s_1)}=q\sum_x\fbar_{s(-u_-(x))}$, and relation
C(ii) becomes $\fbar_{s(s_1)}*\fbar_{s(u_+(x))}*\fbar_{s(s_1)}
=q\,\fbar_{s(-u_-(-x))}$ for $x\in\cO^\times$.
\end{example}

\subsection{Explicit signs for \texorpdfstring{$n=2$}{n=2}}
This subsection computes, for $n=2$, the signs that appear when the
relations involving $t_0$ are transferred. These relations leave $K$,
so Proposition~\ref{prop:Ksub} does not apply to them.

\begin{lemma}
\label{lemma:signs}
Let $n=2$ and $p$ odd. Let $\tilde t_0$ be a lift of $t_0$, and let
$\tilde t_0^{-1}$ denote its inverse in $\Gbar$.
\begin{enumerate}[leftmargin=2em,label=(\alph*)]
\item Every lift of $-1$ has square $(-1,-1)_2=1$.
\item $\tilde t_0$ commutes with every lift of $h(u)$, for
$u\in\cO^\times$.
\item For every lift $\tilde s_1$ of $s_1$,
$\tilde s_1\tilde t_0^{-1}\tilde s_1^{-1}=(-1,\varpi)_2\,\tilde t_0$,
and $(-1,\varpi)_2=(-1)^{(q-1)/2}$.
\end{enumerate}
None of these signs depends on the choice of lifts.
\end{lemma}

\begin{proof}
(a) We have $x(-1)=-1$ and $x(1)=1$, so
$(-1,\zeta)^2=(1,\beta(-1,-1)\zeta^2)=(1,(-1,-1)_2)$. Since $p$ is odd,
$-1$ is a unit, and $(-1,-1)_2=1$ by (H5).

(b) By Lemma~\ref{lemma:cover}(a), the commutator is
$(u,\varpi)_2^2=1$.

(c) Changing $\tilde s_1$ by an element of $\mu_2$ does not change the
conjugation. Changing $\tilde t_0$ to $\zeta\tilde t_0$ multiplies the
left side by $\zeta^{-1}$ and the right side by $\zeta$, and
$\zeta^{-2}=1$. So we may take $\tilde s_1=(s_1,1)$ and
$\tilde t_0=(t_0,1)$. We compute the values of $\beta$ that occur.
First, $\beta(s_1,s_1^{-1})=(1,-1)_2=1$, so $(s_1,1)^{-1}=(s_1^{-1},1)$.
Next, for $v=\varpi^{-1}$ we have $s_1h(v)=\mat 0{-v^{-1}}v0$, so
$\beta(s_1,h(v))=(v,v^2)_2=1$. Then $x(s_1h(v))=\varpi^{-1}$,
$x(s_1^{-1})=-1$ and $x(t_0)=\varpi^{-1}$, so
$\beta(s_1h(v),s_1^{-1})=(1,-\varpi^{-1})_2=1$. Hence
$(s_1,1)(t_0^{-1},1)(s_1,1)^{-1}=(t_0,1)$. Finally
$\beta(t_0,t_0^{-1})=(\varpi^{-1},\varpi)_2=(-1,\varpi)_2$, so
$(t_0,1)^{-1}=(t_0^{-1},(-1,\varpi)_2)$. Combining these,
$\tilde s_1\tilde t_0^{-1}\tilde s_1^{-1}=(-1,\varpi)_2\,(t_0,1)$. The
symbol $(-1,\varpi)_2$ is the Legendre symbol of $-1$ modulo $\varpi$,
which is $(-1)^{(q-1)/2}$.
\end{proof}

By Lemma~\ref{lemma:signs}(c), the constant $\varepsilon$ in relation
($\overline{\mbox{A}}$)(ii) of Proposition~\ref{prop:main2} equals
$(-1,\varpi)_2$ when the lift of $t_0^{-1}$ is the inverse of the lift
of $t_0$. Thus
\[
\fbar_{\tilde t_0}*\fbar_{\tilde s_1}
=(-1,\varpi)_2\ \fbar_{\tilde s_1}*\fbar_{\tilde t_0^{-1}} .
\]
This sign is not an obstruction by itself. Replacing the lift of
$t_0^{-1}$ by its negative removes it.

\begin{example}
\label{ex:signs}
In the running example $F=\mathbb{Q}_7$, $n=2$, we have $q=7\equiv 3$
modulo $4$, so $(-1,7)_2=-1$ and
$\fbar_{\tilde t_0}*\fbar_{\tilde s_1}
=-\fbar_{\tilde s_1}*\fbar_{\tilde t_0^{-1}}$. On the other hand
$\tilde t_0$ commutes with $\tilde h(3)$, in contrast with the case
$n=3$ (Example~\ref{ex:q7}). Over $F=\mathbb{Q}_5$ the sign is $+1$.
\end{example}

\subsection{No homomorphism with constants of absolute value one}
This subsection shows that, although no relation computed above is an
obstruction for $n=2$, the support-preserving maps of
Proposition~\ref{prop:admissible} are never algebra homomorphisms. The
argument is global. It rests on Lemma~\ref{lemma:nonsplit}, and it
works for every $n\geq 2$.

\begin{theorem}
\label{thm:unimodular}
Let $n\geq 2$. Let $K_0$ be an open compact subgroup of $G$ with a
splitting $s$, put $K_0^*=s(K_0)$, and assume that every element of $G$
is $K_0^*$-admissible. Then there is no algebra homomorphism
$\Psi:\cH(G/\!/K_0)\to\cH(\Gbar/\!/K_0^*)_{\gen}$ such that, for every
representative $g$ of a double coset $K_0gK_0$,
$\Psi(f_g)=c_g\fbar_{\tilde g}$ for some lift $\tilde g$ of $g$ and
some complex number $c_g$ with $|c_g|=1$.
\end{theorem}

\begin{proof}
The idea is to compare the structure constants of the two algebras.
Under $\Psi$ each structure constant of $\cH(G/\!/K_0)$, a positive
integer, becomes a sum of that many numbers of absolute value one. This
forces a splitting of the cover.

Suppose that $\Psi$ exists. Define $\psi$ on $\Gbar$ by
$\psi(\tilde y)=\Psi(f_y)(\tilde y)$. Then $\psi$ is genuine, $|\psi|=1$
everywhere, and $\Psi(f_g)=\psi\cdot(f_g\circ\pr)$ for every $g$. Let
$x,z\in G$ and $y=xz$. Let $S$ be the set of left cosets $x'K_0$ in
$K_0xK_0$ with $x'^{-1}y\in K_0zK_0$. Then $|S|=(f_x*f_z)(y)$, and
$S$ contains $xK_0$. Since $\vol(K_0)=1$,
\[
\bigl(\Psi(f_x)*\Psi(f_z)\bigr)(\tilde y)
=\sum_{x'K_0\in S}\psi(\tilde x')\,\psi(\tilde x'^{-1}\tilde y).
\]
On the other hand $f_x*f_z$ is a combination of the $f_h$, and its
coefficient at the double coset of $y$ is $|S|$, so
$\Psi(f_x*f_z)(\tilde y)=|S|\,\psi(\tilde y)$. Every term of the sum
has absolute value $1$, and the sum has absolute value $|S|$. So all
terms are equal, and each equals $\psi(\tilde y)$. The term for
$x'=x$ gives $\psi(\tilde x)\psi(\tilde z)=\psi(\tilde x\tilde z)$ for
all lifts $\tilde x$ and $\tilde z=\tilde x^{-1}\tilde y$. Hence $\psi$
is a homomorphism $\Gbar\to\C^\times$ with $\psi(\zeta)=\zeta$ for
$\zeta\in\mu_n$. Let $N$ be its kernel. Then $N\cap\mu_n=1$. The
quotient $\Gbar/\mu_nN$ is abelian and is a quotient of $G$, which is
perfect, so $\mu_nN=\Gbar$ and $\pr$ maps $N$ isomorphically onto $G$.
Its inverse is a splitting of $\Gbar\to G$, which contradicts
Lemma~\ref{lemma:nonsplit}.
\end{proof}

\begin{corollary}
\label{cor:unimodular}
Let $n\geq 2$ and let $k$ be deep. For no choice of representatives $g$ of
the double cosets $B_kgB_k$, of lifts $\tilde g$, and of constants $c_g$
of absolute value one is $f_g\mapsto c_g\fbar_{\tilde g}$ an algebra
homomorphism from $\Hlin$ to $\Hgen$.
In particular the map $\Phi$ of Proposition~\ref{prop:admissible} is
never an algebra homomorphism, for any choice of lifts.
\end{corollary}

\begin{proof}
By Proposition~\ref{prop:admissible} every element of $G$ is
$B_k^*$-admissible. Apply Theorem~\ref{thm:unimodular} with $K_0=B_k$.
\end{proof}

\begin{remark}
The constants that come from a choice of lifts lie in $\mu_n$, so they
have absolute value one. The proof uses this hypothesis only to force
the terms of a sum of roots of unity to be equal. For $n\geq 3$,
Corollary~\ref{corollary:main} below removes the hypothesis. For $n=2$
we do not know whether it can be removed (Question~\ref{q:n2}).
\end{remark}

\begin{remark}[Where the obstruction lives for $n=2$]
\label{rem:n2where}
For $n=2$ and $p$ odd, the constraints that we have computed from
products whose volumes multiply can all be met. On $SL(2,\cO)$ they are
met exactly (Proposition~\ref{prop:Ksub}). The torus gives no
constraint (Lemma~\ref{lemma:signs}(b)), and the sign in the relation
between $t_0$ and $s_1$ is absorbed by the lift of $t_0^{-1}$
(Lemma~\ref{lemma:signs}(c)). The proof of Theorem~\ref{thm:unimodular}
uses instead products whose volumes do not multiply. An example is the
product of $f_{w(\varpi)}$ and $f_{s_1}$. Here $w(\varpi)=-t_0s_1$, whose
double coset has volume $q^3$, while $w(\varpi)s_1=t_0$, whose double
coset has volume $q^2$. In such a product
$(f_x*f_z)(xz)$ is a sum of several terms. Constants of absolute value
one force these terms to agree, and the identity $w(a)w(-1)=h(a)$ of
Lemma~\ref{lemma:nonsplit} then carries the Hilbert symbol $(b,a)_2$
into the algebra. Constants of other absolute values might allow the
terms to cancel instead, which is why Question~\ref{q:n2} remains open.
\end{remark}

\begin{example}
\label{ex:unimodular}
In the running example $F=\mathbb{Q}_7$, $n=2$, the contradiction in
the proof of Theorem~\ref{thm:unimodular} comes from
Example~\ref{ex:nonsplit}. A homomorphism $\psi$ as in the proof would
give a splitting $s$ with $s(h(7))s(h(3))=-s(h(21))$.
\end{example}

\subsection{The structure constants \texorpdfstring{$N_k(x,z)$}{N} and
subalgebras on which \texorpdfstring{$\Phi$}{Phi} is multiplicative}
\label{subsec:Nk}
This subsection looks more closely at the structure constants used in
the proof of Theorem~\ref{thm:unimodular}. They do not shrink as the
level grows, so the obstruction is present at every deep level.
Products with structure constant one transfer individually. On a
monoid of such products, $\Phi$ is an algebra isomorphism exactly when
the lifts can be chosen multiplicatively. This holds for $n=2$ on the
monoid generated by $B$ and the positive translations, and it fails
for $n\geq3$ on the same monoid.

For deep $k$ and $x,z\in G$, let $N_k(x,z)$ be the coefficient of the
double coset of $xz$ in $f_x*f_z\in\Hlin$, that is, the number of cosets
$x'B_k$ in $B_kxB_k$ with $x'^{-1}xz\in B_kzB_k$. It depends on the
pair $(x,z)$ and on $k$.

\begin{remark}[Properties of $N_k(x,z)$]
\label{rem:Nk}
\begin{enumerate}[leftmargin=2em,label=(\alph*)]
\item Suppose that the volumes multiply, that is,
\[
\vol(B_kxB_k)\vol(B_kzB_k)=\vol(B_kxzB_k).
\]
Then $f_x*f_z=f_{xz}$ and $N_k(x,z)=1$, by
Lemma~\ref{lemma:general}(d).
These products transfer to $\Hgen$ by Lemma~\ref{lemma:genuine}.
\item If $z=x^{-1}$, every coset $x'B_k$ in $B_kxB_k$ is counted, since
$x'^{-1}\in B_kx^{-1}B_k$. So $N_k(x,x^{-1})=\vol(B_kxB_k)/\vol(B_k)$.
If $x\in BwB$ with $w\in W_a$, then $B_kxB_k$ is a translate of
$B_kwB_k$ on both sides by elements of $B$, which normalizes $B_k$, so
$N_k(x,x^{-1})=q^{\ell(w)}$ by Lemma~\ref{lemma:length}. This does not
depend on $k$. For example $N_k(s_1,s_1)=q$, since $s_1^{-1}=-s_1$ and
$-1$ is central, and $N_k(t_0,t_0^{-1})=q^2$, for every deep $k$.
\item Theorem~\ref{thm:unimodular} is a statement at each fixed level,
and its proof at level $k$ uses only pairs at level $k$. By (b), pairs
with $N_k\geq q$ exist at every level. So the obstruction does not
disappear as $k$ grows.
\item The condition $N_k=1$ is neither necessary nor sufficient for
multiplicativity. When $p\nmid n$, the part of $\Hlin$ supported on
$K$ transfers exactly (Proposition~\ref{prop:Ksub}), although it
contains $f_{s_1}*f_{s_1}$ with $N_k(s_1,s_1)=q$. The reason is that
the cover splits over $K$. Conversely, for $n\geq3$ the products
$f_{t_0}*f_{h(u)}$ and $f_{h(u)}*f_{t_0}$ both have structure constant
one, but no constants make them transfer (Theorem~\ref{thrm:main}).
Proposition~\ref{prop:monoid} makes this precise.
\end{enumerate}
\end{remark}

\begin{vremark}
\label{vrem:Nk}
\emph{Computer check.} For $q=3$ and $k=1,2,3$ we computed $N_k(x,z)$
for $x,z\in\pm W_a$ by counting cosets. The results are as follows.
\begin{center}
\begin{tabular}{@{}lll@{}}
\toprule
$x$, $z$ & $N_1,N_2,N_3$ & Iwahori-level constant \\
\midrule
$s_1$, $s_1$ & $3,3,3$ & $q$ \\
$s_2$, $s_2$ & $3,3,3$ & $q$ \\
$s_1$, $t_0$ & $3,3,3$ & $q$ \\
$t_0$, $t_0^{-1}$ & $9,9,9$ & $q^2$ \\
$t_0^2$, $t_0^{-1}$ & $9,9,9$ & $q^2$ \\
$t_0$, $s_1$ & $1,1,1$ & $1$ \\
$t_0$, $t_0$ & $1,1,1$ & $1$ \\
$s_1$, $s_2$ & $1,1,1$ & $1$ \\
\bottomrule
\end{tabular}
\end{center}
In these cases $N_k(x,z)$ does not depend on $k$, and it equals the
coefficient of $f_{xz}$ in $f_x*f_z\in\cH(G/\!/B)$. For instance
$f_{s_1}*f_{t_0}=(q-1)f_{s_1s_2}+qf_{s_2}$, and $s_1t_0=s_2$. Also
$N_k(x,z)=1$ exactly when $\ell(xz)=\ell(x)+\ell(z)$. We have not proved
these observations in general.
\end{vremark}

\begin{proposition}
\label{prop:monoid}
Let $k$ be deep, and let $M\subseteq G$ be a submonoid with
$B_kMB_k=M$ such that $\vol(B_kxB_k)\vol(B_kzB_k)=\vol(B_kxzB_k)$ for
all $x,z\in M$. Let $A_M\subseteq\Hlin$ and $\overline A_M\subseteq\Hgen$
be the spans of the $f_x$ and of the $\fbar_{\tilde x}$ with $x\in M$.
\begin{enumerate}[leftmargin=2em,label=(\alph*)]
\item $A_M$ and $\overline A_M$ are subalgebras, $f_x*f_z=f_{xz}$ and
$N_k(x,z)=1$ for $x,z\in M$, and $\Phi$ restricts to a linear
isomorphism from $A_M$ onto $\overline A_M$.
\item Suppose that $L:M\to\Gbar$ is a homomorphism of monoids that
lifts the inclusion and agrees with $\sigma$ on $B_k$. Then
$f_x\mapsto\fbar_{L(x)}$ is an isomorphism of algebras from $A_M$ onto
$\overline A_M$. It is $\Phi$ for the lifts $\tilde x=L(x)$.
\item Suppose that $x,z\in M$ commute and that some, hence every, pair
of lifts satisfies $\tilde x\tilde z\neq\tilde z\tilde x$. Then no map
$f_y\mapsto c_y\fbar_{\tilde y}$ with non-zero constants is an algebra
homomorphism on $A_M$.
\end{enumerate}
\end{proposition}

\begin{proof}
(a) The first two statements follow from Remark~\ref{rem:Nk}(a) and
Lemma~\ref{lemma:genuine}. The third follows from
Proposition~\ref{prop:admissible}. (b) If $x'=bxb'$ with $b,b'\in
B_k$, then $L(x')=\sigma(b)L(x)\sigma(b')$, so $\fbar_{L(x)}$ depends
only on $B_kxB_k$. By Lemma~\ref{lemma:genuine},
$\fbar_{L(x)}*\fbar_{L(z)}=\fbar_{L(x)L(z)}=\fbar_{L(xz)}$. (c) The
commutator $\zeta=\tilde z\tilde x\tilde z^{-1}\tilde x^{-1}$ lies in
$\mu_n$ and does not depend on the lifts. By Lemma~\ref{lemma:genuine},
$\fbar_{\tilde x}*\fbar_{\tilde z}=\fbar_{\tilde x\tilde z}$ and
$\fbar_{\tilde z}*\fbar_{\tilde x}=\fbar_{\zeta\tilde x\tilde
z}=\zeta^{-1}\fbar_{\tilde x\tilde z}$, which is non-zero. A
homomorphism would carry $f_x*f_z=f_z*f_x$ to
$c_xc_z\fbar_{\tilde x\tilde z}=c_xc_z\zeta^{-1}\fbar_{\tilde x\tilde
z}$, which is impossible.
\end{proof}

\begin{example}
\label{ex:monoid}
Let $M_+=\bigcup_{m\geq0}Bt_0^mB$. Since $\ell(t_0^m)=2m$ for $m\geq0$,
Proposition~\ref{prop:IM}(a) gives $Bt_0^mB\,Bt_0^{m'}B=Bt_0^{m+m'}B$
for $m,m'\geq0$, so $M_+$ is a monoid. By Lemma~\ref{lemma:length},
$\vol(B_kxB_k)=q^{2m}\vol(B_k)$ for $x\in Bt_0^mB$. So $M_+$ satisfies
the hypotheses of Proposition~\ref{prop:monoid}. Recall that
$t_0=h(\varpi)$ conjugates $u_+(y)$ to $u_+(\varpi^2y)$ and $u_-(y)$ to
$u_-(\varpi^{-2}y)$, and that $B=U^+(\cO)\,T(\cO)\,U^-_1$.
\begin{enumerate}[leftmargin=2em,label=(\roman*)]
\item Let $n=2$ and $p$ odd. Let $s$ be a splitting of $B$, fix a lift
$\tilde t_0$ of $t_0$, and put $L(bt_0^mb')=s(b)\tilde t_0^m s(b')$
for $b,b'\in B$ and $m\geq0$. Then $L$ is a well-defined homomorphism
of monoids, and it agrees with $\sigma$ on $B_k$ by
Proposition~\ref{prop:Ksub}. So by Proposition~\ref{prop:monoid}(b),
$\Phi$ with the lifts $L(x)$ is an algebra isomorphism from
$A_{M_+}$ onto $\overline A_{M_+}$.
\item Let $n\geq3$. The elements $t_0$ and $h(u)$ of $M_+$ commute,
but their lifts commute only up to $(u,\varpi)_n^2$
(Lemma~\ref{lemma:cover}(a)). This is not $1$ for a suitable unit $u$
(Lemma~\ref{lemma:hilbert}). By Proposition~\ref{prop:monoid}(c), no
constants make $\Phi$ multiplicative on $A_{M_+}$, although every
product in $M_+$ has structure constant one.
\item For every $n$, the smaller monoid
$M_k=\bigcup_{m\geq0}B_kt_0^mB_k$ gives the polynomial algebra
$A_{M_k}=\C[f_{t_0}]$, with $f_{t_0}^{*m}=f_{t_0^m}$. Here
$L(bt_0^mb')=\sigma(b)\tilde t_0^m\sigma(b')$, for $b,b'\in B_k$, is a
well-defined homomorphism of monoids. So $\Phi$ is an algebra
isomorphism from $A_{M_k}$ onto $\overline A_{M_k}$, and
$\fbar_{\tilde t_0}^{*m}=\fbar_{\tilde t_0^m}$.
\end{enumerate}
\end{example}

\begin{proof}[Proof of the claims in (i) and (iii)]
Consider (i). First, if $\gamma\in B$ and $t_0^j\gamma t_0^{-j}\in B$
for some integer $j$, then
\begin{equation}
\label{eq:conjs}
\tilde t_0^j\,s(\gamma)\,\tilde t_0^{-j}=s(t_0^j\gamma t_0^{-j}).
\end{equation}
Indeed, conjugation by $t_0^j$ preserves $U^\pm$ and $T$, so the
factors of $\gamma=\gamma_+\gamma_0\gamma_-$ in $U^+(\cO)$, $T(\cO)$
and $U^-_1$ have the same property as $\gamma$. For $\gamma_0$,
(\ref{eq:conjs}) holds because $\tilde t_0$ commutes with $s(\gamma_0)$
(Lemma~\ref{lemma:signs}(b)). For $\gamma_+$, note that $\beta$ is
trivial on $U^+\times U^+$ (Lemma~\ref{lemma:cover}(b)), so $u\mapsto
(u,1)$ is a splitting of $U^+(F)\cong F$. Any two splittings of a
subgroup isomorphic to $F$ or to $\cO$ differ by a homomorphism to
$\mu_n$, which is trivial because $F=nF$ and $\cO=n\cO$. So $s(u)=(u,1)$
on $U^+(\cO)$, and $u\mapsto\tilde t_0^j(t_0^{-j}ut_0^j,1)\tilde
t_0^{-j}$ is a splitting of $U^+(F)$, hence equals $u\mapsto(u,1)$.
This gives (\ref{eq:conjs}) for $\gamma_+$, and $\gamma_-$ is treated in
the same way.

If $bt_0^mb'=ct_0^mc'$, put $\gamma=c^{-1}b$. Then $\gamma\in B$ and
$t_0^{-m}\gamma t_0^m=c'b'^{-1}\in B$, so (\ref{eq:conjs}) with $j=-m$
gives $s(\gamma)\tilde t_0^m=\tilde t_0^ms(c'b'^{-1})$. Hence
$s(b)\tilde t_0^ms(b')=s(c)\tilde t_0^ms(c')$, and $L$ is well
defined. For multiplicativity, let $x=bt_0^mb'$ and $z=ct_0^{m'}c'$,
and write $b'c=\beta_+\beta_0\beta_-$ as above. Put
$\beta_+'=t_0^m\beta_+t_0^{-m}\in U^+(\cO)$ and
$\beta_-'=t_0^{-m'}\beta_-t_0^{m'}\in U^-_1$. Then
$xz=(b\beta_+'\beta_0)\,t_0^{m+m'}\,(\beta_-'c')$. By (\ref{eq:conjs}),
\[
L(x)L(z)=s(b)\,\tilde t_0^m s(\beta_+)\,s(\beta_0)\,s(\beta_-)\tilde
t_0^{m'}\,s(c') =s(b\beta_+'\beta_0)\,\tilde
t_0^{m+m'}\,s(\beta_-'c')=L(xz).
\]
For (iii), $M_k$ is a monoid by Lemma~\ref{lemma:general}(d), since
$\vol(B_kt_0^mB_k)=q^{2m}\vol(B_k)$. The same argument applies with
$B_k=U^+_kT_kU^-_{k+1}$ in place of $B$ and $\sigma$ in place of $s$.
Now (\ref{eq:conjs}) holds for every $n$ by
Lemma~\ref{lemma:canonical}(b).
\end{proof}

%======================================================================
\part{The case \texorpdfstring{$n\geq 3$}{n>=3}}

Part~II shows that for $n\geq 3$ the comparison fails for a local
reason, which no choice of constants can repair. It then relates the
algebras to representations and treats Iwahori level. The last section
covers every $n\geq 2$, including $n=2$.

\section{The torus obstruction}
\label{section:torus}

This section shows that for $n\geq 3$ the commutator of lifts of $t_0$
and of a diagonal unit is a non-trivial root of unity. This single
relation already rules out every support-preserving homomorphism, with
arbitrary non-zero constants. Throughout, $k$ is deep and the notation
is that of Section~\ref{section:metaplectic}.

By Proposition~\ref{prop:main2}, every relation of
Proposition~\ref{prop:main} holds in $\Hgen$ after the elements of $G$
are replaced by lifts adapted to that relation. For $n\geq 3$ the next
theorem shows that the lifts cannot be chosen once and for all, whatever
constants are allowed.

\begin{theorem}
\label{thrm:main}
Let $k$ be deep and $u\in\cO^\times$. Let $\tilde t_0$ and
$\tilde h(u)$ be any lifts of $t_0$ and $h(u)$. Then
\[
\fbar_{\tilde t_0}*\fbar_{\tilde h(u)}
=(u,\varpi)_n^{-2}\ \fbar_{\tilde h(u)}*\fbar_{\tilde t_0}\neq 0,
\qquad\mbox{whereas}\qquad
f_{t_0}*f_{h(u)}=f_{h(u)}*f_{t_0}.
\]
\end{theorem}

\begin{proof}
By C(i) of Proposition~\ref{prop:main}, both linear products equal
$f_{t_0h(u)}$. By ($\overline{\mbox{C}}$)(i) of
Proposition~\ref{prop:main2},
$\fbar_{\tilde t_0}*\fbar_{\tilde h(u)}=\fbar_{\tilde t_0\tilde h(u)}$
and
$\fbar_{\tilde h(u)}*\fbar_{\tilde t_0}=\fbar_{\tilde h(u)\tilde t_0}$.
By Lemma~\ref{lemma:cover}(a) with $a=\varpi$ and $b=u$, we have
$\tilde t_0\tilde h(u)=(u,\varpi)_n^2\,\tilde h(u)\tilde t_0$. The
claim follows from Lemma~\ref{lemma:support}(b).
\end{proof}

\begin{corollary}
\label{corollary:main}
Let $k$ be deep and $n\geq 3$. There is no algebra homomorphism
$\Psi:\Hlin\to\Hgen$ such that, for every $g\in G$, $\Psi(f_g)$ is a
non-zero function supported on $\overline{B_kgB_k}$. In particular, for
no choice of lifts is the linear isomorphism $\Phi$ of
Proposition~\ref{prop:admissible} an isomorphism of algebras.
\end{corollary}

\begin{proof}
By Lemma~\ref{lemma:support}(a), $\Psi(f_g)=c_g\fbar_{\tilde g}$ with
$c_g\neq 0$. By Lemma~\ref{lemma:hilbert} choose $u\in\cO^\times$ with
$\zeta=(u,\varpi)_n^2\neq 1$. Apply $\Psi$ to the relation
$f_{t_0}*f_{h(u)}=f_{h(u)}*f_{t_0}$ and use Theorem~\ref{thrm:main}.
This gives
\[
(\zeta^{-1}-1)\,c_{t_0}c_{h(u)}\,
\fbar_{\tilde h(u)}*\fbar_{\tilde t_0}=0,
\]
which is a contradiction.
\end{proof}

\begin{example}
\label{ex:q7}
In the running example $F=\mathbb{Q}_7$, $n=3$, take $u=3$. By
Example~\ref{ex:hilbert}, $\zeta=(3,7)_3^2$ is a primitive cube root of
unity. So $\fbar_{\tilde t_0}*\fbar_{\tilde h(3)}
=\zeta^{-1}\fbar_{\tilde h(3)}*\fbar_{\tilde t_0}$ in
$\cH(\Gbar/\!/B_k^*)_{\gen}$ for every $k\geq 1$.
\end{example}

\begin{remark}
Corollary~\ref{corollary:main} concerns support-preserving maps. It
does not by itself exclude an abstract isomorphism of $\C$-algebras
between $\Hlin$ and $\Hgen$, and we make no claim either way. An
abstract isomorphism that does not respect supports would not be
expected to respect traces of representations.
\end{remark}

%======================================================================
\section{Representation theory}
\label{section:reps}

This section explains how genuine Hecke algebras control genuine
representations. A genuine representation with a non-zero
$K_0^*$-fixed vector gives a module over the genuine Hecke algebra,
and every simple module arises in this way from exactly one irreducible
representation (Proposition~\ref{prop:borel}). We then show that for
$n\geq 3$ the traces of the genuine Hecke functions cannot match the
traces of the corresponding linear ones (Example~\ref{ex:trace}).
Throughout, $K_0$ is an open compact subgroup of $G$ with a splitting
$s$, and $K_0^*=s(K_0)$.

A smooth representation $(\pi,V)$ of $\Gbar$ on which $\mu_n$ acts by
the inverse of the fixed character $\mu_n\to\C^\times$ is called
\textbf{genuine}. With this convention a genuine function $\phi$ acts
on $V$ by $\pi(\phi)=\int_{\Gbar}\phi(\tilde x)\pi(\tilde x)\,d\tilde
x$, and the integrand is invariant under $\mu_n$. Replacing the fixed
character by its inverse exchanges this convention with the opposite
one. For a genuine smooth $(\pi,V)$, the space $V^{K_0^*}$ of
$K_0^*$-fixed vectors is a module over $\cH(\Gbar/\!/K_0^*)_{\gen}$.

For a module $(r,W)$ over $\cH=\cH(\Gbar/\!/K_0^*)_{\gen}$, let
\[
I_{\gen}(W)=C_c(\Gbar/K_0^*)_{\gen}\otimes_{\cH}W
\]
be the \textbf{induced} $\Gbar$-module of Borel \cite[\S2]{Bo76}. Here
$C_c(\Gbar/K_0^*)_{\gen}$ is the space of genuine, compactly supported,
right $K_0^*$-invariant functions on $\Gbar$. The group $\Gbar$ acts on
it by left translation, and $\cH$ acts by convolution on the right.

\begin{proposition}[{\cite[\S\S1--2]{Bo76}, \cite[\S4.3]{BH06}}]
\label{prop:borel}
Let $(r,W)$ be a module over $\cH=\cH(\Gbar/\!/K_0^*)_{\gen}$, and let
$(\pi,V)$ be a genuine smooth representation of $\Gbar$.
\begin{enumerate}[leftmargin=2em,label=(\alph*)]
\item There is a natural isomorphism of $\cH$-modules
$W\cong I_{\gen}(W)^{K_0^*}$.
\item $I_{\gen}(W)$ is generated, as a $\Gbar$-module, by
$I_{\gen}(W)^{K_0^*}$.
\item Restriction to $W$ is an isomorphism
$\Hom_{\Gbar}(I_{\gen}(W),V)\to\Hom_{\cH}(W,V^{K_0^*})$.
\item If $W$ is simple, then $I_{\gen}(W)$ has a unique irreducible
quotient $J_{\gen}(W)$, and $J_{\gen}(W)^{K_0^*}\cong W$. The maps
$V\mapsto V^{K_0^*}$ and $W\mapsto J_{\gen}(W)$ are mutually inverse
bijections between the set of equivalence classes of irreducible
genuine smooth representations $V$ of $\Gbar$ with $V^{K_0^*}\neq 0$
and the set of isomorphism classes of simple $\cH$-modules.
\end{enumerate}
\end{proposition}

\begin{proof}
The algebra $\cH(\Gbar/\!/K_0^*)$ of all bi-$K_0^*$-invariant compactly
supported functions is the direct sum, over the characters $\chi$ of
$\mu_n$, of the two-sided ideals of functions that transform by $\chi$
under $\mu_n$. The algebra $\cH$ is the summand for the fixed
character. A smooth representation of the kind considered is
annihilated by the other summands. The statements therefore follow
from the corresponding general facts for a locally profinite group and
an open compact subgroup. Parts (a)--(c) are in \cite[\S\S1--2]{Bo76}.
The bijection in (d) is in \cite[\S4.3]{BH06}. The description of the
inverse map follows from (c), since the sum of all
$\Gbar$-submodules of $I_{\gen}(W)$ with no non-zero $K_0^*$-fixed
vectors is the unique maximal proper submodule when $W$ is simple.
\end{proof}

\begin{vremark}
\label{vrem:borel}
\emph{Cited, not re-checked against the sources.} The section numbers
of \cite{Bo76} and \cite{BH06} in Proposition~\ref{prop:borel} are
quoted from memory and from the 1998 version.
\end{vremark}

\begin{remark}
In (d), $I_{\gen}(W)$ itself need not be irreducible for a general
open compact subgroup. Irreducibility for the Iwahori subgroup of a
linear group is a special feature established in \cite{Bo76}.
\end{remark}

Applied to $K_0=B_k$ with $k$ deep, Proposition~\ref{prop:borel}
identifies the irreducible genuine smooth representations of $\Gbar$
with non-zero $B_k^*$-fixed vectors with the simple modules of the
algebra described in Proposition~\ref{prop:main2}. In the same way, the
irreducible smooth representations of $G$ with non-zero $B_k$-fixed
vectors correspond to the simple modules of $\Hlin$.

\begin{example}
\label{ex:trace}
Let $k$ be deep and $n\geq 3$. Let $(\pi,V)$ be an admissible genuine
representation with $V^{B_k^*}\neq 0$, so that $V^{B_k^*}$ is
finite-dimensional. Choose a unit $u$ with
$\zeta=(u,\varpi)_n^2\neq1$. Put $A=\pi(\fbar_{\tilde t_0})$ and
$C=\pi(\fbar_{\tilde h(u)})$ on $V^{B_k^*}$. By
Theorem~\ref{thrm:main}, $AC=\zeta^{-1}CA$. The operator $C$ is
invertible, because $\fbar_{\tilde h(u)}*\fbar_{\tilde h(u)^{-1}}$ is
the identity by ($\overline{\mbox{B}}$)(ii). Hence
$A=\zeta^{-1}CAC^{-1}$ and $\tr A=\zeta^{-1}\tr A$, so $\tr A=0$. No
such vanishing holds for $G$. For the trivial representation of $G$,
$\tr\pi(f_{t_0})=\vol(B_kt_0B_k)=q^2$. Consequently there is no
bijection $\overline\pi\mapsto\pi$, from the irreducible admissible
genuine representations of $\Gbar$ with $B_k^*$-fixed vectors onto the
irreducible admissible representations of $G$ with $B_k$-fixed
vectors, such that
$\tr\overline\pi(\fbar_{\tilde g})=\tr\pi(f_g)$ for all $g$ and some
choice of lifts. In the running example this applies with $u=3$ and
every $k\geq 1$.
\end{example}

The argument of Example~\ref{ex:trace} applies to every translation that
fails to commute with the diagonal units.

\begin{proposition}
\label{prop:traceless}
Let $k$ be deep and let $(\pi,V)$ be an admissible genuine
representation of $\Gbar$. Let $m\in\Z$ with $n\nmid 2m$, and let
$\tilde t$ be a lift of $t_0^m$. Then $\tr\pi(\fbar_{\tilde t})=0$ on
$V^{B_k^*}$.
\end{proposition}

\begin{proof}
By Lemma~\ref{lemma:hilbert} there is a unit $u$ with $(u,\varpi)_n$ a
primitive $n$-th root of unity. Then $\zeta=(u,\varpi)_n^{2m}\neq 1$,
because $n\nmid 2m$. By Lemma~\ref{lemma:cover}(a) with $a=\varpi^m$ and
$b=u$, we have $\tilde t\,\tilde h(u)=\zeta\,\tilde h(u)\,\tilde t$.
Since $h(u)\in B$, relation ($\overline{\mbox{C}}$)(i) of
Proposition~\ref{prop:main2} and Lemma~\ref{lemma:support}(b) give
$\fbar_{\tilde t}*\fbar_{\tilde h(u)}
=\zeta^{-1}\fbar_{\tilde h(u)}*\fbar_{\tilde t}$. Put
$A=\pi(\fbar_{\tilde t})$ and $C=\pi(\fbar_{\tilde h(u)})$. Then $C$ is
invertible, $A=\zeta^{-1}CAC^{-1}$, and so $\tr A=\zeta^{-1}\tr A$.
Hence $\tr A=0$.
\end{proof}

\begin{remark}
For the linear group the corresponding traces need not vanish. For the
trivial representation of $G$, $\tr\pi(f_{t_0^m})=\vol(B_kt_0^mB_k)
=q^{2|m|}$ by Lemma~\ref{lemma:length}. For $n=2$ the condition
$n\nmid 2m$ never holds, so Proposition~\ref{prop:traceless} gives no
vanishing. This is consistent with $PGL(2)$, whose lattice of translations
contains that of $SL(2)$.
\end{remark}

\medskip\noindent\textbf{Spherical parts.} The algebra $\Hlin$ is not commutative (already the span of the $f_x$
with $x\in B$ is not, when $q>3$), so no Satake transform can be
injective on it. What can be computed is the spherical part of each
basis function, obtained by averaging over $K=SL(2,\cO)$ on both sides.
The next proposition computes it on the linear side and shows that on
the genuine side it vanishes for the same translations as in
Proposition~\ref{prop:traceless}.

Normalize Haar measure by $\vol(B_k)=1$, let $N=[K:B_k]$, and let
$e_K=N^{-1}\chi_K$, the identity of the subalgebra of $\Hlin$ of
bi-$K$-invariant functions. For the Satake transform we use the Haar
measure with $\vol(K)=1$, and for a bi-$K$-invariant function $\phi$ we
put
\[
\mathcal S(\phi)(X)=\sum_{j\in\Z}X^j\,q^{-j}\int_{U^+}\phi\bigl(h(\varpi^j)u\bigr)\,du,
\]
where $du$ gives $U^+\cap K$ volume one and $X$ is the value at $t_0$ of
an unramified character of the diagonal torus. With this normalization
$\mathcal S(\chi_K)=1$.

\begin{proposition}
\label{prop:spherical}
Let $k\geq 1$. Let $g\in G$, write $g=b_1wb_2$ with $b_1,b_2\in B$ and
$w\in\pm W_a$, and let $m\geq 0$ be the integer with $g\in Kt_0^mK$.
\begin{enumerate}[leftmargin=2em,label=(\alph*)]
\item $e_K*f_g*e_K=q^{\ell(w)}\,\vol(Kt_0^mK)^{-1}\,\chi_{Kt_0^mK}$, and
$\vol(Kt_0^mK)=N\,q^{2m}(1+q^{-1})$ for $m\geq 1$.
\item For $m\geq 1$,
\[
\mathcal S(\chi_{Kt_0^mK})=q^m\Bigl(\sum_{j=-m}^{m}X^j
-q^{-1}\sum_{j=-(m-1)}^{m-1}X^j\Bigr).
\]
Consequently the spherical part $N\,e_K*f_g*e_K$ of $f_g$, viewed with
the measure $\vol(K)=1$, has Satake transform
\[
\frac{q^{\ell(w)-m}}{1+q^{-1}}\Bigl(\sum_{j=-m}^{m}X^j
-q^{-1}\sum_{j=-(m-1)}^{m-1}X^j\Bigr)\quad(m\geq 1),
\qquad\mbox{and}\qquad q^{\ell(w)}\quad(m=0).
\]
\item Assume $p\nmid n$, let $s$ be a splitting of $K$, and let
$\overline e$ be $N^{-1}$ times the genuine function supported on
$\overline K$ with value $\zeta$ at $\zeta s(x)$. If $n\nmid 2m$, then
$\overline e*\fbar_{\tilde g}*\overline e=0$ for every lift $\tilde g$
of $g$.
\end{enumerate}
\end{proposition}

\begin{proof}
(a) The function $e_K*f_g*e_K$ is bi-$K$-invariant and supported on
$KgK$. Since $b_1,b_2\in K$, and $w$ is $\pm t_0^j$ or $\pm t_0^js_1$
with $|j|=m$, we have $KgK=Kt_0^mK$. A bi-$K$-invariant function
supported on one double coset is constant there. By
Lemma~\ref{lemma:general}(c) its integral is
$\int e_K\cdot\int f_g\cdot\int e_K=\vol(B_kgB_k)=q^{\ell(w)}$, by
Lemma~\ref{lemma:length}. This gives the constant. Next,
$\vol(Kt_0^mK)=N\,[K:K\cap t_0^{-m}Kt_0^m]$, and $K\cap t_0^{-m}Kt_0^m$
is the subgroup of $K$ of matrices whose lower left entry lies in
$\varpi^{2m}\cO$. Its index is the number of points of the projective
line over $\cO/\varpi^{2m}\cO$, which is $q^{2m}+q^{2m-1}$.

(b) An element of $G$ lies in $Kt_0^mK$ exactly when the least order of
its entries is $-m$. For $u=u_+(x)$ the entries of $h(\varpi^j)u$ are
$\varpi^j$, $\varpi^jx$, $0$ and $\varpi^{-j}$. Let $v$ be the order of
$x$. If $j=m$ the condition is $v\geq -2m$, a set of volume $q^{2m}$. If
$j=-m$ it is $v\geq 0$, of volume $1$. If $|j|<m$ it is $v=-m-j$, of
volume $q^{m+j}(1-q^{-1})$. For $|j|>m$ the set is empty. Summing
$X^jq^{-j}$ times these volumes gives
$q^m\bigl(X^m+X^{-m}+(1-q^{-1})\sum_{|j|<m}X^j\bigr)$, which is the
stated formula. The function $N\,e_K*f_g*e_K$ equals
$q^{\ell(w)}[Kt_0^mK:K]^{-1}\chi_{Kt_0^mK}$, where $[Kt_0^mK:K]$ is the
number of left cosets of $K$, which is $q^{2m}(1+q^{-1})$ by (a). This
gives the last formula.

(c) The function $\overline e*\fbar_{\tilde g}*\overline e$ is genuine,
bi-$s(K)$-invariant and supported on $\overline{Kt_0^mK}$. By
Lemma~\ref{lemma:support}(a), applied with $K_0=K$, such a function is
zero unless $t_0^m$ is $s(K)$-admissible. For a unit $v$, the element
$h(v)$ lies in $K\cap t_0^{-m}Kt_0^m$ and commutes with $t_0^m$, and by
Lemma~\ref{lemma:cover}(a) the commutator of the lifts is
$(v,\varpi)_n^{2m}$. By Lemma~\ref{lemma:hilbert} this is not $1$ for
some $v$ when $n\nmid 2m$. So $t_0^m$ is not admissible.
\end{proof}

\begin{remark}
\label{rem:spherical}
The formula in (b) is the case of $SL(2)$ of Macdonald's formula
\cite{Ma71}. At $X=q$ the transform in (b) equals $q^{\ell(w)}$, the
trace of $f_g$ on the trivial representation, as it should. By (c), the
spherical parts of the genuine basis functions vanish over exactly the
translations that Savin's theorem removes. For $n=2$ no vanishing is
forced. When $n\mid 2m$ the genuine spherical part is a multiple of the
genuine spherical function on $Kt_0^mK$. The multiple is a sum of roots
of unity that we have not computed. By Corollary~\ref{cor:spherical0}, Savin's isomorphism identifies the
genuine spherical Hecke algebra with the spherical Hecke algebra of
$G/Z_n$, by an explicit map with the rescaling of
Theorem~\ref{thm:savinsl2}(b). Composing with Macdonald's formula for
$G/Z_n$ gives a Satake transform for the genuine algebra. We have not
checked whether it agrees with a Satake isomorphism defined directly
through genuine unramified principal series.
\end{remark}

\begin{example}
\label{ex:spherical}
In the running example $F=\mathbb{Q}_7$ with $n=3$, take $g=t_0$, so
$\ell(w)=2$ and $m=1$. The linear spherical part has Satake transform
$\frac{49}{8}\bigl(X+1+X^{-1}-\frac17\bigr)$. The genuine spherical
part of $\fbar_{\tilde t_0}$ is zero, and so is that of
$\fbar_{\tilde t_0^2}$. With $n=2$ neither vanishing is forced.
\end{example}

%======================================================================
\section{The Iwahori level}
\label{section:savin}

This section treats level $k=0$ in the tame case $p\nmid n$, for every
$n\geq 2$. It determines which double cosets of the Iwahori subgroup
carry genuine functions, gives a presentation of the genuine
Iwahori--Hecke algebra by two generators and two quadratic relations,
and deduces Savin's theorem for $SL(2)$ with an explicit isomorphism.
Background and context are in Section~\ref{subsec:savinbg}. We use the
notation $n'=n/\gcd(n,2)$ and $r=t_0^{-n'}s_1$ of (\ref{eq:nprime}).

\subsection{Supports}
This subsection shows that the genuine Iwahori--Hecke algebra sees only
the translations $t_0^m$ and the reflections $t_0^ms_1$ with $n\mid
2m$. For $n\geq 3$ this already distinguishes it from $\cH(G/\!/B)$.

The cover splits over $B$. Indeed, the pro-$p$ radical $B^+$ of $B$
has the canonical splitting $\sigma$ of Lemma~\ref{lemma:canonical}. By
(H5) and Lemma~\ref{lemma:cover}(a), the map $u\mapsto(h(u),1)$ is a
splitting of $T(\cO)=\{h(u)\,|\,u\in\cO^\times\}$, and it agrees with
$\sigma$ on $T_1$ by uniqueness. Then
$s(h(u)b)=(h(u),1)\sigma(b)$, for $u\in\cO^\times$ and $b\in B^+$, is a
splitting of $B=T(\cO)B^+$ by Lemma~\ref{lemma:canonical}(b). Any other
splitting of $B$ differs from $s$ by a homomorphism $B\to\mu_n$. From
now on $B^*=s(B)$ for this $s$, except in
Proposition~\ref{prop:level0}, which holds for every splitting.

\begin{proposition}
\label{prop:level0}
Let $p\nmid n$, let $s$ be any splitting of $B$, and put $B^*=s(B)$.
The double coset $Bt_0^mB$ supports a non-zero genuine
bi-$B^*$-invariant function if and only if $n$ divides $2m$. In
particular, if $n\geq 3$, then $\cH(\Gbar/\!/B^*)_{\gen}$ contains no non-zero
function supported on $\overline{Bt_0B}$.
\end{proposition}

\begin{proof}
Let $g=t_0^m=h(\varpi^m)$ and $x=h(v)$ with $v\in\cO^\times$. Then
$x\in B\cap g^{-1}Bg$ and $gxg^{-1}=x$. By
Lemma~\ref{lemma:support}(a) the hypothesis implies
(\ref{eq:admissible}), so $s(x)$ commutes with $\tilde g$. By
Lemma~\ref{lemma:cover}(a), the commutator of lifts of $h(\varpi^m)$
and $h(v)$ is $(v,\varpi)_n^{2m}$. By Lemma~\ref{lemma:hilbert}, this
is trivial for all units $v$ only if $n$ divides $2m$.

Conversely, assume that $n$ divides $2m$. By
Lemma~\ref{lemma:support}(a) we must check (\ref{eq:admissible}) for
every $x\in B\cap g^{-1}Bg$. Write $x=h(v)b'$ with $v\in\cO^\times$
and $b'\in B^+$. Since $g$ is diagonal, $gxg^{-1}=h(v)\,gb'g^{-1}$, and
$gb'g^{-1}$ lies in $B$ and has the same diagonal entries as $b'$, so it
lies in $B^+$. The restriction of $s$ to the pro-$p$ group $B^+$ is
$\sigma$, by Lemma~\ref{lemma:canonical}(d). By
Lemma~\ref{lemma:canonical}(b),
$\tilde g\,\sigma(b')\,\tilde g^{-1}=\sigma(gb'g^{-1})$. By
Lemma~\ref{lemma:cover}(a),
$\tilde g\,s(h(v))\,\tilde g^{-1}=(v,\varpi)_n^{2m}s(h(v))=s(h(v))$.
Hence
\[
\tilde g\,s(x)\,\tilde g^{-1}=s(h(v))\,\sigma(gb'g^{-1})
=s(gxg^{-1}),
\]
which is (\ref{eq:admissible}).
\end{proof}

\begin{lemma}
\label{lemma:level0s1}
Let $p\nmid n$, and let $s$ be the splitting of $B$ constructed above.
The double coset $Bt_0^ms_1B$ supports a non-zero genuine
bi-$B^*$-invariant function if and only if $n$ divides $2m$.
\end{lemma}

\begin{proof}
For $v\in F^\times$ we have $s_1h(v)=h(v^{-1})s_1=w(-v^{-1})$. The
function $x(\cdot)$ takes the values $1$, $v^{-1}$, $v$ and $v$ on
$s_1$, $h(v)$, $h(v^{-1})$ and $w(-v^{-1})$. Hence
$\beta(s_1,h(v))=(v,v^2)_n=(-1,v)_n^2=1$ and
$\beta(h(v^{-1}),s_1)=(1,v)_n=1$, so
\[
(s_1,1)\,(h(v),1)=(h(v^{-1}),1)\,(s_1,1).
\]
Let $g=t_0^ms_1$, let $\tilde g=\tilde t_0^m(s_1,1)$ for a lift
$\tilde t_0^m$ of $t_0^m$, let $v$ be a unit, and let $x=h(v)$. Then
$x\in B\cap g^{-1}Bg$ and $gxg^{-1}=h(v^{-1})$. By the displayed
identity and Lemma~\ref{lemma:cover}(a),
\[
\tilde g\,s(x)\,\tilde g^{-1}
=\tilde t_0^m\,s(h(v^{-1}))\,\tilde t_0^{-m}
=(v,\varpi)_n^{-2m}\,s(gxg^{-1}).
\]
So (\ref{eq:admissible}) forces $n\mid 2m$, by
Lemma~\ref{lemma:hilbert}. Conversely, let $n\mid 2m$, and write
$x\in B\cap g^{-1}Bg$ as $h(v)b'$ with $b'\in B^+$. Then $gb'g^{-1}$
lies in $B$, and its diagonal entries are those of $b'$ in the other
order, so it lies in $B^+$. By Lemma~\ref{lemma:canonical}(b),
$\tilde g\sigma(b')\tilde g^{-1}=\sigma(gb'g^{-1})$, and
(\ref{eq:admissible}) follows as in the proof of
Proposition~\ref{prop:level0}.
\end{proof}

\begin{example}
\label{ex:level0}
For $n=2$ the condition $n\mid 2m$ holds for every $m$. In the running
example $n=3$, the double cosets $Bt_0^mB$ and $Bt_0^ms_1B$ can carry a
genuine function only when $3\mid m$. So $Bt_0B$ and $Bt_0^2B$ carry
none, while $Bt_0^3B$ does. The translations that survive form the
lattice $3\Z$, a rescaling of the lattice $\Z$ of the linear group.
\end{example}

\begin{remark}[The same lattice at every level]
\label{rem:samelattice}
Compare Proposition~\ref{prop:level0} with
Proposition~\ref{prop:traceless}. At Iwahori level the translations
$t_0^m$ that carry genuine functions are exactly those with $n\mid 2m$.
At deep level the others also carry genuine functions, but these act
with trace zero in every admissible genuine representation. So at
Iwahori level the genuine algebra has exactly the translations $t_0^m$
with $n\mid 2m$, and at deep level the translations that the genuine
theory can see through traces are among these. They are exactly the
translations whose lifts commute with the lifts of all diagonal units.
By Theorem~\ref{thm:savinsl2}, this is the lattice that Savin's
isomorphism matches, after rescaling, with the translations of $G/Z_n$.
This is the precise sense in which $SL(2)$ is the wrong partner and
$G/Z_n$ the right one, at least at Iwahori level and as far as traces
of translations are concerned.
\end{remark}

At Iwahori level, therefore, the linear and genuine Hecke algebras do
not even have the same supports when $n\geq 3$. This is in contrast
with Proposition~\ref{prop:admissible}.

\subsection{Generators and relations}
This subsection shows that the genuine Iwahori--Hecke algebra is
generated by $X=\fbar_{(s_1,1)}$ and $Y=\fbar_{(r,1)}$, and that two
quadratic relations, computed explicitly, are a defining set. The
computation uses the elements $w(y)=\mat 0y{-y^{-1}}0$ of
Section~\ref{subsec:cover} and two cocycle identities for them. Recall
that $s_1=w(-1)$ and $r=w(-\varpi^{-n'})$.

\begin{lemma}
\label{lemma:wformulas}
Let $y,a\in F^\times$, and write $\tilde u_\pm(z)=(u_\pm(z),1)$.
\begin{enumerate}[leftmargin=2em,label=(\alph*)]
\item $(w(y),1)=\tilde u_+(y)\,\tilde u_-(-y^{-1})\,\tilde u_+(y)
=\tilde u_-(-y^{-1})\,\tilde u_+(y)\,\tilde u_-(-y^{-1})$.
\item $(w(ay),1)=(-y^{-1},a^{-1})_n^{-1}\,(h(a),1)\,(w(y),1)$.
\item $(w(y),1)^{-1}=(w(-y),1)$.
\end{enumerate}
\end{lemma}

\begin{proof}
The first equality in (a) was shown in the proof of
Lemma~\ref{lemma:nonsplit}. For the second, put
$g=u_-(-y^{-1})u_+(y)=\mat 1y{-y^{-1}}0$. Then $gu_-(-y^{-1})=w(y)$.
The function $x(\cdot)$ takes the value $-y^{-1}$ on $g$, on
$u_-(-y^{-1})$ and on $w(y)$, and the value $1$ on $u_+(y)$. Hence
$\beta(u_-(-y^{-1}),u_+(y))=(1,-y^{-1})_n=1$ and
$\beta(g,u_-(-y^{-1}))=(1,1)_n=1$. For (b), $w(ay)=h(a)w(y)$, and
$x(\cdot)$ takes the values $a^{-1}$, $-y^{-1}$ and $-(ay)^{-1}$ on
$h(a)$, $w(y)$ and $w(ay)$. So $\beta(h(a),w(y))=(-y^{-1},a^{-1})_n$.
For (c), $w(y)w(-y)=1$, and $\beta(w(y),w(-y))=(-y,y)_n=1$ by (H3).
\end{proof}

\begin{lemma}
\label{lemma:words}
Let $W_n'$ be the set of the elements $\pm t_0^m$ and $\pm t_0^ms_1$ of
$G$ with $n'\mid m$.
\begin{enumerate}[leftmargin=2em,label=(\alph*)]
\item $W_n'$ is a group, and for $k\geq 0$
\[
(s_1r)^k=(-1)^kt_0^{kn'},\quad (rs_1)^k=(-1)^kt_0^{-kn'},\quad
(s_1r)^ks_1=(-1)^kt_0^{kn'}s_1,\quad
(rs_1)^kr=(-1)^kt_0^{-(k+1)n'}s_1 .
\]
\item Every element of $W_n'$ is, up to sign, a unique alternating
product $a_1a_2\cdots a_j$ of $s_1$ and $r$. Moreover
$\ell(a_1\cdots a_i)=\ell(a_1\cdots a_{i-1})+\ell(a_i)$ for $1\leq
i\leq j$. The product for $t_0^{n'k}$ or $t_0^{n'k}s_1$ has $|k|$
factors $r$.
\end{enumerate}
\end{lemma}

\begin{proof}
We have $rs_1=t_0^{-n'}s_1^2=-t_0^{-n'}$. Since
$s_1t_0^{-1}s_1^{-1}=t_0$, also $s_1r=s_1t_0^{-n'}s_1^{-1}s_1^2
=-t_0^{n'}$. The four formulas follow. They show that $W_n'$ is the
group generated by $s_1$, $r$ and $-1$, and that each of its elements
is, up to sign, exactly one of the listed products. The lengths of the
four products are $2kn'$, $2kn'$, $2kn'+1$ and $2(k+1)n'-1$. These are
the sums of the lengths of the factors, since $\ell(s_1)=1$ and
$\ell(r)=2n'-1$. Every initial segment of an alternating product is
again one, which gives the additivity at each step. The number of
factors $r$ can be read off from the formulas.
\end{proof}

\begin{proposition}
\label{prop:quadratic}
Let $p\nmid n$. Put $X=\fbar_{(s_1,1)}$ and $Y=\fbar_{(r,1)}$, and
write $1$ for the identity element $\fbar_{(1,1)}$ of
$\cH(\Gbar/\!/B^*)_{\gen}$. Then
\begin{enumerate}[leftmargin=2em,label=(\alph*)]
\item $X^2=(q-1)X+q$.
\item $Y^2=c_1Y+c_0$, where
\[
c_0=(-1,\varpi)_n^{n'}\,q^{2n'-1},\qquad
c_1=q^{n'-1}\sum_{u}(\varpi,-u)_n^{n'},
\]
and $u$ runs over representatives in $\cO^\times$ of $\Fq^\times$.
\end{enumerate}
Explicitly, $c_1=q^{n-1}(q-1)$ and $c_0=q^{2n-1}$ if $n$ is odd, and
$c_1=0$ and $c_0=(-1)^{(q-1)/2}q^{n-1}$ if $n$ is even.
\end{proposition}

\begin{proof}
We prove (b) first. Put $x_0=-\varpi^{-n'}$, so that $r=w(x_0)$, and
put $\tilde r=(r,1)$.

\emph{Support.} The function $Y*Y$ is supported over the support of
$f_r*f_r$, which lies in $BrBrB$. Since $r^2=-1\in B$, we have
$BrBrB=B(rBr^{-1})B$. By the Iwahori factorization
$B=U^+_0T(\cO)U^-_1$, and since conjugation by $r$ maps $U^+_0$ onto
$U^-_{2n'}$, $U^-_1$ onto $U^+_{1-2n'}$ and $T(\cO)$ onto itself, we
get $rBr^{-1}=U^-_{2n'}T(\cO)U^+_{1-2n'}$. Hence $BrBrB=BU^+_{1-2n'}B$.
Let $y\in\varpi^{1-2n'}\cO$. If $y\in\cO$, then $u_+(y)\in B$.
Otherwise $y=\varpi^iu$ with $1-2n'\leq i\leq -1$ and $u$ a unit, and
\begin{equation}
\label{eq:uplus}
u_+(y)=u_-(y^{-1})\,w(y)\,u_-(y^{-1}),\qquad w(y)=t_0^is_1h(-u^{-1}),
\end{equation}
by Lemma~\ref{lemma:wformulas}(a). Here $u_-(y^{-1})\in B$, so
$u_+(y)\in Bt_0^is_1B$. Hence $Y*Y$ is supported over $B$ and over the
double cosets $Bt_0^is_1B$ with $-2n'<i<0$. By
Lemma~\ref{lemma:level0s1}, among the latter only $i=-n'$ carries a
genuine function, and $BrB$ is that double coset. So $Y*Y=c_0+c_1Y$,
where $c_0=(Y*Y)(1,1)$ and $c_1=(Y*Y)(\tilde r)$.

\emph{The coefficient $c_0$.} In Lemma~\ref{lemma:support}(d) there
are $\vol(BrB)=q^{2n'-1}$ representatives $v$, and each gives the term
$Y(\tilde r^{-1}s(v)^{-1})=Y(\tilde r^{-1})$. Since $r^{-1}=w(-x_0)$
and $-x_0=x_0\cdot(-1)$, Lemma~\ref{lemma:wformulas}(b) and (c) give
\[
\tilde r^{-1}=(w(-x_0),1)=(\varpi^{n'},-1)_n^{-1}\,(h(-1),1)\,\tilde r .
\]
As $(h(-1),1)=s(h(-1))\in B^*$, we get
$Y(\tilde r^{-1})=(\varpi^{n'},-1)_n^{-1}=(-1,\varpi)_n^{n'}$ by (H1)
and (H2).

\emph{The coefficient $c_1$.} By the proof of
Lemma~\ref{lemma:length}, the elements $v=u_-(z)$, with $z$ running
over representatives of $\varpi\cO/\varpi^{2n'}\cO$, represent
$B/(B\cap rBr^{-1})$. By Lemma~\ref{lemma:canonical}(c),
$s(v)=\tilde u_-(z)$. Conjugation by $r^{-1}$ maps $u_-(-z)$ to
$u_+(z\varpi^{-2n'})$. By the uniqueness in
Lemma~\ref{lemma:cover}(b) the same holds for the lifts, since
$z\mapsto\tilde r^{-1}\tilde u_-(-z)\tilde r$ is a homomorphism. So
Lemma~\ref{lemma:support}(d) gives
\[
c_1=\sum_zY\bigl(\tilde r^{-1}\tilde u_-(-z)\tilde r\bigr)
=\sum_zY\bigl(\tilde u_+(z\varpi^{-2n'})\bigr).
\]
Put $y=z\varpi^{-2n'}$. By the first part of the proof, the term
vanishes unless $\ord(y)=-n'$. In that case $y=\varpi^{-n'}u$, where
$u$ runs over representatives of $(\cO/\varpi^{n'}\cO)^\times$, so
that each class in $\Fq^\times$ occurs $q^{n'-1}$ times. By
Lemma~\ref{lemma:wformulas}(a),
$\tilde u_+(y)=\tilde u_-(y^{-1})(w(y),1)\tilde u_-(y^{-1})$, and
$\tilde u_-(y^{-1})\in B^*$. Since $y=x_0\cdot(-u)$,
Lemma~\ref{lemma:wformulas}(b) gives
\[
Y(\tilde u_+(y))=Y\bigl((w(y),1)\bigr)=(\varpi^{n'},-u^{-1})_n^{-1}
=(\varpi,-u)_n^{n'} .
\]
This depends only on the residue of $u$, by (H5), and the formula for
$c_1$ follows.

\emph{The explicit values.} By Lemma~\ref{lemma:hilbert},
$u\mapsto(\varpi,u)_n$ maps $\cO^\times$ onto $\mu_n$. So
$u\mapsto(\varpi,u)_n^{n'}$ is trivial if $n'=n$, that is, if $n$ is
odd, and it is a character of order $2$ of $\Fq^\times$ if $n$ is even.
This gives $c_1$. For odd $n$, $(-1,\varpi)_n^{n}=1$. For even $n$,
$(-1,\varpi)_n^{n/2}=(-1)^{(q-1)/2}$ by Example~\ref{ex:hilbert}. This
gives $c_0$.

\emph{Part (a).} Replace $r$ by $s_1=w(-1)$. The support of $X*X$ lies
in $B\cup Bs_1B$. The elements $u_+(z)$, with $z$ running over
representatives of $\cO/\varpi\cO$, represent the cosets in
$B/(B\cap s_1Bs_1^{-1})$, and
\[
(s_1,1)^{-1}\,\tilde u_+(-z)\,(s_1,1)=\tilde u_-(z).
\]
For a unit $z$,
Lemma~\ref{lemma:wformulas}(a) gives $\tilde u_-(z)=\tilde
u_+(z^{-1})(w(-z^{-1}),1)\tilde u_+(z^{-1})$, and
Lemma~\ref{lemma:wformulas}(b) with $a=z^{-1}$ gives
$X((w(-z^{-1}),1))=(1,z)_n^{-1}=1$. So $(X*X)(s_1,1)=q-1$. Likewise
$(s_1,1)^{-1}=(w(1),1)=(1,-1)_n^{-1}(h(-1),1)(s_1,1)$, so
$(X*X)(1,1)=q$.
\end{proof}

\begin{theorem}
\label{thm:presentation}
Let $p\nmid n$, and let $X$, $Y$, $c_0$ and $c_1$ be as in
Proposition~\ref{prop:quadratic}.
\begin{enumerate}[leftmargin=2em,label=(\alph*)]
\item The functions $\fbar_{(w,1)}$, for $w=t_0^m$ and $w=t_0^ms_1$
with $n'\mid m$, form a basis of $\cH(\Gbar/\!/B^*)_{\gen}$.
\item If $w=a_1\cdots a_j$ is the alternating product of
Lemma~\ref{lemma:words}, then $\fbar_{(w,1)}$ is a root of unity times
the corresponding alternating product of $X$ and $Y$.
\item The algebra $\cH(\Gbar/\!/B^*)_{\gen}$ is generated by $X$ and
$Y$, and the relations $X^2=(q-1)X+q$ and $Y^2=c_1Y+c_0$ are defining
relations.
\end{enumerate}
\end{theorem}

\begin{proof}
Part (a) follows from Lemma~\ref{lemma:support}(c),
Proposition~\ref{prop:level0} and Lemma~\ref{lemma:level0s1}, since
$-1\in B$. For (b), Lemma~\ref{lemma:length} with $k=0$ and
Lemma~\ref{lemma:words}(b) show that the volumes multiply at each step
of the alternating product. By Lemma~\ref{lemma:genuine}, applied $j-1$
times, the alternating product of $X$ and $Y$ equals
$\fbar_{\tilde w}$, where $\tilde w$ is the product of the lifts
$(s_1,1)$ and $(r,1)$. This $\tilde w$ is a lift of $\pm w$, and
$(-1,1)=s(-1)\in B^*$, so $\fbar_{\tilde w}$ is a root of unity times
$\fbar_{(w,1)}$ by Lemma~\ref{lemma:support}(b).

For (c), let $A$ be the algebra with generators $\xi$ and $\eta$ and
relations $\xi^2=(q-1)\xi+q$ and $\eta^2=c_1\eta+c_0$. Using the
relations, every element of $A$ is a linear combination of alternating
products of $\xi$ and $\eta$, including the empty product. By
Proposition~\ref{prop:quadratic} there is a homomorphism from $A$ to
the genuine Iwahori--Hecke algebra with $\xi\mapsto X$ and
$\eta\mapsto Y$. By (a) and (b) it maps the alternating products
bijectively onto non-zero multiples of the basis elements. So the
alternating products are linearly independent in $A$, and the
homomorphism is an isomorphism.
\end{proof}

\begin{example}
\label{ex:presentation}
In the running example $F=\mathbb{Q}_7$ with $n=3$, the genuine
Iwahori--Hecke algebra has a basis indexed by the $t_0^{3j}$ and the
$t_0^{3j}s_1$. The function $Y$ is supported over $Bt_0^{-3}s_1B$, of
volume $7^5$, and Proposition~\ref{prop:quadratic} gives
$Y^2=294\,Y+16807$, since $q^2(q-1)=294$ and $q^5=16807$. Since
$t_0^{-3}=-rs_1$, the genuine function on $\overline{Bt_0^{-3}B}$ is a
root of unity times $Y*X$. With $n=2$ we have $Y=\fbar_{(s_2,1)}$ and
$Y^2=-7$, since $(-1)^{(7-1)/2}=-1$.
\end{example}

\begin{vremark}
\label{vrem:computer}
\emph{Computer check.} We checked Proposition~\ref{prop:quadratic}
numerically, independently of the proof. We computed $X*X$ and $Y*Y$
directly from Kubota's cocycle, with exact Hilbert symbols, for
$F=\mathbb{Q}_7$ with $n=2,3,6$ and for $F=\mathbb{Q}_5$ with $n=4$.
The splitting of $B$ used is the restriction of the splitting of $K$
given by $g\mapsto(c,d)_n$ when the lower left entry $c$ of $g$ lies in
$\varpi\cO\setminus\{0\}$, and by $g\mapsto 1$ otherwise. We checked on
random elements of $SL(2,\mathbb{Z}_p)$ that this is a splitting. In
each case the coefficients of $Y*Y$ on the double cosets that carry no
genuine function are vanishing sums of roots of unity, and the other
coefficients agree with Proposition~\ref{prop:quadratic}. For $n=3$ and
$q=7$ they are $294$ on $\overline{BrB}$ and $16807$ on
$\overline{B}$. The scripts are in the repository, next to the source
of this paper.
\end{vremark}

\subsection{Savin's isomorphism for \texorpdfstring{$SL(2)$}{SL(2)}}
This subsection compares the presentation of
Theorem~\ref{thm:presentation} with those of
Propositions~\ref{prop:IM} and \ref{prop:IMpgl}. The comparison proves
Savin's theorem for $SL(2)$ with an explicit isomorphism, and it shows
in which sense $SL(2)$ is the wrong partner when $n$ is even.

\begin{theorem}[Savin's theorem for $SL(2)$]
\label{thm:savinsl2}
Let $p\nmid n$. If $n$ is odd, let $G'=G$, $I'=B$, $f'_w=f_w$,
$\ell'=\ell$, $\mu=t_0$ and $r'=s_2$. If $n$ is even, let
$G'=PGL(2,F)$, let $I'$, $f'_w$ and $\ell'$ be as in
Section~\ref{subsec:savinbg}, and let $\mu=\lambda$ and $r'=[t]$. In
that case we write $s_1$ for $[s_1]$. Put $a_n=q^{n-1}$ if $n$ is odd
and $a_n=\epsilon^{1/2}q^{(n-1)/2}$ if $n$ is even, where
$\epsilon=(-1)^{(q-1)/2}$ and $\epsilon^{1/2}$ is $1$ if $\epsilon=1$
and $i$ if $\epsilon=-1$.
\begin{enumerate}[leftmargin=2em,label=(\alph*)]
\item There is a unique algebra isomorphism
$\Theta:\cH(\Gbar/\!/B^*)_{\gen}\to\cH(G'/\!/I')$ with
$\Theta(X)=f'_{s_1}$ and $\Theta(Y)=a_nf'_{r'}$.
\item Let $\Psi(t_0^{n'k})=\mu^k$ and $\Psi(t_0^{n'k}s_1)=\mu^ks_1$.
For $w=t_0^{n'k}$ and for $w=t_0^{n'k}s_1$,
\[
\Theta\bigl(\fbar_{(w,1)}\bigr)=c_w\,q^{e(w)}f'_{\Psi(w)},\qquad
e(w)=\frac{\ell(w)-\ell'(\Psi(w))}{2}=\frac{(n-1)|k|}{\gcd(n,2)},
\]
with $|c_w|=1$.
\item $\Theta$ preserves the trace $\phi\mapsto\phi(1)$, where $1$ is
the identity element of $\Gbar$, respectively of $G'$. It also
preserves the involution $\phi^*(x)=\overline{\phi(x^{-1})}$.
\item $\Theta$ maps $(1+X)/(1+q)$ to $(1+f'_{s_1})/(1+q)$.
\end{enumerate}
\end{theorem}

\begin{proof}
(a) By Propositions~\ref{prop:IM} and \ref{prop:IMpgl}, $f'_{s_1}$
satisfies the relation of $X$. For odd $n$,
\[
(q^{n-1}f_{s_2})^2=q^{n-1}(q-1)\,q^{n-1}f_{s_2}+q^{2n-1}
=c_1\,q^{n-1}f_{s_2}+c_0 .
\]
For even $n$, $(a_nf'_{[t]})^2=a_n^2=\epsilon q^{n-1}=c_0$, and
$c_1=0$. So by Theorem~\ref{thm:presentation}(c) there is a unique
homomorphism $\Theta$ with the stated values. By
Propositions~\ref{prop:IM}(c) and \ref{prop:IMpgl}(c), the map with
$f'_{s_1}\mapsto X$ and $f'_{r'}\mapsto a_n^{-1}Y$ is also a
well-defined homomorphism. It is inverse to $\Theta$ on generators, so
$\Theta$ is an isomorphism.

(b) We have $\Psi(s_1)=s_1$ and $\Psi(r)=\mu^{-1}s_1=r'$, since
$t_0^{-1}s_1=s_2$ and $\lambda^{-1}[s_1]=[t]$. Also $(s_1s_2)^k$ and
$([s_1][t])^k$ are $\pm t_0^k$ and $\lambda^k$. By
Lemma~\ref{lemma:words}, $w$ is, up to sign, an alternating product of
$s_1$ and $r$ with $|k|$ factors $r$, and $\Psi(w)$ is, up to sign, the
same alternating product of $s_1$ and $r'$. Lengths add along both
products. By Theorem~\ref{thm:presentation}(b), $\fbar_{(w,1)}$ is a
root of unity times the alternating product of $X$ and $Y$. By
Propositions~\ref{prop:IM}(a) and \ref{prop:IMpgl}(a), $f'_{\Psi(w)}$
is the alternating product of $f'_{s_1}$ and $f'_{r'}$, because $-1\in
B$ and $-1\in I'$. Hence $\Theta(\fbar_{(w,1)})$ is a root of unity
times $a_n^{|k|}f'_{\Psi(w)}$, and
$|a_n|^{|k|}=q^{(n-1)|k|/\gcd(n,2)}$. Finally,
$\ell(w)-\ell'(\Psi(w))=|k|\bigl(\ell(r)-\ell'(r')\bigr)$, which is
$|k|(2n-2)$ for odd $n$ and $|k|(n-1)$ for even $n$.

(c) By (b), $\Theta$ maps each $\fbar_{(w,1)}$ with $w\neq\pm 1$ to a
multiple of some $f'_{w'}$ with $w'\neq\pm1$, and both traces vanish
on these functions. Also $\Theta(1)=1$. For the involution, both sides
of $\Theta(\phi^*)=\Theta(\phi)^*$ are conjugate linear and reverse
products, so it suffices to check the identity on $X$ and $Y$. Since
$s_1^{-1}=-s_1$ and $r^{-1}=-r$, the functions $X^*$ and $Y^*$ are
multiples of $X$ and $Y$. The proof of Proposition~\ref{prop:quadratic}
gives $X((s_1,1)^{-1})=1$ and $Y(\tilde r^{-1})=(-1,\varpi)_n^{n'}$,
which is $\pm1$. So $X^*=X$ and $Y^*=(-1,\varpi)_n^{n'}Y$. In the same
way $f'^*_{s_1}=f'_{s_1}$ and $f'^*_{r'}=f'_{r'}$. It remains to check
that $(-1,\varpi)_n^{n'}a_n=\overline{a_n}$. For odd $n$ both sides are
$q^{n-1}$. For even $n$ the left side is $\epsilon\,\epsilon^{1/2}
q^{(n-1)/2}$, and $\epsilon\,\epsilon^{1/2}=\overline{\epsilon^{1/2}}$.

(d) This is clear from (a).
\end{proof}

\begin{corollary}
\label{cor:wronggroup}
Let $p\nmid n$.
\begin{enumerate}[leftmargin=2em,label=(\alph*)]
\item If $n$ is even, no algebra isomorphism from
$\cH(\Gbar/\!/B^*)_{\gen}$ onto $\cH(G/\!/B)$ maps $Y$ to a multiple of
$f_{s_1}$ or of $f_{s_2}$.
\item If $n$ is odd, the isomorphism $\Theta$ maps the genuine function
on $\overline{Bt_0^{nk}B}$ to a multiple of $f_{t_0^k}$, and the genuine
algebra contains no non-zero function supported over $Bt_0^mB$ when
$n\nmid m$.
\end{enumerate}
\end{corollary}

\begin{proof}
(a) An isomorphism preserves minimal polynomials. For even $n$ the
minimal polynomial of $Y$ is $z^2-c_0$ with $c_0\neq 0$, so its roots
have sum $0$. A non-zero multiple $cf_{s_i}$ has minimal polynomial
$z^2-c(q-1)z-c^2q$, whose roots $cq$ and $-c$ have sum $c(q-1)\neq 0$.
Part (b) restates Theorem~\ref{thm:savinsl2}(b) and
Proposition~\ref{prop:level0}.
\end{proof}

\begin{remark}[Abstract isomorphism]
\label{rem:abstract}
Corollary~\ref{cor:wronggroup}(a) concerns maps that respect a
generator. As abstract algebras, $\cH(\Gbar/\!/B^*)_{\gen}$,
$\cH(G/\!/B)$ and $\cH(PGL(2,F)/\!/I')$ are all isomorphic, for every
$n$. Indeed, in each presentation both generators have two distinct
eigenvalues: $q$ and $-1$ for $X$, $f_{s_i}$ and $f'_{[s_1]}$, then
$\pm1$ for $f'_{[t]}$, and for $Y$ the roots $q^n$ and $-q^{n-1}$ when
$n$ is odd and $\pm\sqrt{c_0}$ when $n$ is even. Replacing each
generator by the idempotent that projects onto one eigenvalue turns
each presentation into that of the algebra generated by two idempotents
with no further relation. That algebra is the group algebra of the
infinite dihedral group $\Z/2*\Z/2$. So at Iwahori level the abstract
algebra structure detects nothing. What distinguishes $SL(2)$ from
$PGL(2)$, and what Savin's isomorphism preserves, is the extra
structure in Theorem~\ref{thm:savinsl2}: the trace, the involution, the
finite part spanned by $1$ and $X$, and the lines through the basis
functions up to the rescaling $\Psi$. This should be kept in mind for
the abstract part of Question~\ref{q:n2}.
\end{remark}

\begin{corollary}
\label{cor:spherical0}
Let $p\nmid n$, and let $s_K$ be a splitting of $K=SL(2,\cO)$, which
exists by \cite{Ku67} (Remark~\ref{vrem:kubota}). Then $s_K$ restricts
to $s$ on $B$, and $s_K(s_1)=(s_1,1)$. Put $K^*=s_K(K)$. Then $\Theta$
restricts to an isomorphism from $\cH(\Gbar/\!/K^*)_{\gen}$ onto
$\cH(G'/\!/K')$, where $K'=K$ for odd $n$, and $K'$ is the image of
$GL(2,\cO)$ in $PGL(2,F)$ for even $n$. In particular the genuine
spherical Hecke algebra is commutative, and the irreducible genuine
representations of $\Gbar$ with non-zero $K^*$-fixed vectors correspond
bijectively to the irreducible unramified representations of $G'$.
\end{corollary}

\begin{proof}
On the pro-$p$ groups $U^\pm\cap K$, the splitting $s_K$ agrees with
$u\mapsto(u,1)$, as in the proof of Lemma~\ref{lemma:canonical}(d). So
$s_K(w(y))=(w(y),1)$ for every unit $y$, by
Lemma~\ref{lemma:wformulas}(a). In particular $s_K(s_1)=(s_1,1)$. For a
unit $u$ we have $h(u)=w(u)w(-1)$ and $\beta(w(u),w(-1))=(-1,u^{-1})_n
=1$ by (H5), so $s_K(h(u))=(h(u),1)=s(h(u))$. On $B^+$ both $s_K$ and
$s$ equal $\sigma$, by Lemma~\ref{lemma:canonical}(d). Hence
$s_K|_B=s$. Since $K=B\sqcup Bs_1B$ and $\vol(K)=1+q$, the identity
element of $\cH(\Gbar/\!/K^*)_{\gen}$ is $e=(1+X)/(1+q)$, and
$\cH(\Gbar/\!/K^*)_{\gen}=e*\cH(\Gbar/\!/B^*)_{\gen}*e$. Likewise
$K'=I'\sqcup I's_1I'$, and $e'=(1+f'_{s_1})/(1+q)$ is the identity
element of $\cH(G'/\!/K')=e'*\cH(G'/\!/I')*e'$. By
Theorem~\ref{thm:savinsl2}(d), $\Theta(e)=e'$. The spherical Hecke
algebra of $G'$ is commutative by the Satake isomorphism \cite{Ma71},
and the last statement follows from Proposition~\ref{prop:borel}(d) and
its linear analog.
\end{proof}

\subsection{Orbital integrals}
\label{subsec:orbital}
This subsection computes orbital integrals of basis functions
explicitly in some cases. At Iwahori level we compute the orbital
integrals at split regular elements of every basis function, on $G$,
on $PGL(2,F)$ and on $\Gbar$. Savin's isomorphism matches them, up to
an explicit constant of absolute value one
(Theorem~\ref{thm:orbtransfer}). This answers
Question~\ref{q:transfer} for split elements at Iwahori level. For the
basis of Proposition~\ref{prop:admissible} at deep level we obtain
vanishing results and a sum rule (Proposition~\ref{prop:orbdeep}).
Orbital integrals at elliptic elements are treated in
Section~\ref{subsec:elliptic}.

\medskip\noindent\textbf{Definitions.} Normalize Haar measure on $T$
by $\vol(T(\cO))=1$. For $a\in F^\times$ with $a^2\neq 1$, the element
$\gamma=h(a)$ is regular, with centralizer $T$, and
\[
O_\gamma(f)=\int_{T\backslash G}f(x^{-1}\gamma x)\,d\dot x,\qquad
F_f(\gamma)=|D(\gamma)|^{1/2}O_\gamma(f),
\]
where $D(h(a))=(1-a^2)(1-a^{-2})$ and $d\dot x$ is the quotient of the
Haar measure of $G$ by that of $T$. Thus
$|D(h(a))|^{1/2}=q^{|\ord(a)|}$ when $\ord(a)\neq 0$. A lift
$\tilde\gamma$ of $h(a)$ is called \textbf{good} if it commutes with
every lift of every element of $T$. By Lemma~\ref{lemma:cover}(a)
this holds if and only if $(b,a)_n^2=1$ for all $b\in F^\times$. For
example, $(t_0^m,1)$ is good if and only if $n\mid 2m$, by
Lemma~\ref{lemma:hilbert}. For a good $\tilde\gamma$ and a genuine
compactly supported locally constant $\phi$ on $\Gbar$, the function
$x\mapsto\phi(\tilde x^{-1}\tilde\gamma\tilde x)$ does not depend on
the lift $\tilde x$ of $x$ and is invariant under $T$. So
$O_{\tilde\gamma}(\phi)=\int_{T\backslash G}\phi(\tilde
x^{-1}\tilde\gamma\tilde x)\,d\dot x$ and
$F_\phi(\tilde\gamma)=|D(\gamma)|^{1/2}O_{\tilde\gamma}(\phi)$ are
defined. At other lifts the integrand is not invariant under $T$,
which is why genuine orbital integrals are taken only at good
elements. On $G'=PGL(2,F)$ we use the image $T'$ of the diagonal
torus, with $\vol(T'(\cO))=1$, the element
$\gamma'=[\mathrm{diag}(b,1)]$, and $D'(\gamma')=(1-b)(1-b^{-1})$, so
that $|D'(\gamma')|^{1/2}=q^{|\ord(b)|/2}$. At Iwahori level Haar
measure is normalized by $\vol(B)=1$, respectively $\vol(I')=1$.

\begin{lemma}
\label{lemma:orbtrace}
Let $\tilde\gamma$ be good, and let $\phi$, $\psi$ be genuine,
compactly supported and locally constant on $\Gbar$. Then
$O_{\tilde\gamma}(\phi*\psi)=O_{\tilde\gamma}(\psi*\phi)$. The same
holds on $G$ and on $G'$.
\end{lemma}

\begin{proof}
For $\tilde y\in\Gbar$ put $\psi_y(\tilde x)=\psi(\tilde y^{-1}\tilde
x)$. Then $\psi(\tilde x\tilde y^{-1})=\psi_y(\tilde y\tilde x\tilde
y^{-1})$. Hence
\[
(\phi*\psi)(\tilde x)=\int_G\phi(\tilde y)\psi_y(\tilde x)\,dy,\qquad
(\psi*\phi)(\tilde x)=\int_G\phi(\tilde y)\psi_y(\tilde y\tilde
x\tilde y^{-1})\,dy .
\]
The orbital integral is invariant under conjugation of its argument by
$\Gbar$, by the change of variables $x\mapsto yx$ in
$T\backslash G$. Integrating over $\tilde y$ gives the claim.
\end{proof}

\begin{lemma}
\label{lemma:orbred}
Let $f$ be a bi-$B$-invariant compactly supported function on $G$, and
let $\gamma=h(a)$ with $a^2\neq1$. Let $R$ consist of the elements
$u_-(y)$, for $y$ in a set of representatives of $\cO/\varpi\cO$,
together with $s_1$. Then
\[
O_\gamma(f)=|1-a^{-2}|^{-1}\sum_{\kappa\in R}\int_F
f\bigl(\kappa^{-1}\gamma u_+(z)\kappa\bigr)\,dz,
\]
where $\vol(\cO)=1$. The same formula holds for a genuine
bi-$B^*$-invariant $\phi$ and a good lift $\tilde\gamma$ of $\gamma$,
with $\phi(\tilde\kappa^{-1}\tilde\gamma\tilde u_+(z)\tilde\kappa)$
in the integrand, for any lifts $\tilde\kappa$.
\end{lemma}

\begin{proof}
By the Iwasawa decomposition $G=TU^+K$, and since
$\vol(K)=[K:B]=q+1$, we have $dg=(q+1)\,dt\,du\,dk$, where $du$ gives
$U^+(\cO)$ volume $1$ and $dk$ gives $K$ volume $1$. This can be
checked on the characteristic function of $K$. Hence
$O_\gamma(f)=(q+1)\int_{U^+}\int_Kf(k^{-1}u^{-1}\gamma uk)\,dk\,du$.
Since $f$ is invariant under conjugation by $B$, the integral over $K$
is $(q+1)^{-1}$ times the sum over a set of representatives of $K/B$,
and $R$ is such a set. Finally
$u_+(-x)\gamma u_+(x)=\gamma u_+((1-a^{-2})x)$. In the genuine case
the integrand does not depend on the lifts, it is invariant under
$B^*$-conjugation, and
$\tilde u_+(-x)\tilde\gamma\tilde u_+(x)=\tilde\gamma\tilde
u_+((1-a^{-2})x)$ by the uniqueness in Lemma~\ref{lemma:cover}(b).
\end{proof}

To locate the double coset of a matrix $g=\mat ABCD\in G$, put
$\nu(A)=\ord A$, $\nu(D)=\ord D$, $\nu(B)=\ord B+\frac12$ and
$\nu(C)=\ord C-\frac12$, with $\ord 0=\infty$. Row and column
operations by elements of $B$ show the following. If the minimum of
the four values is attained only at $A$, then $g\in Bt_0^{\ord A}B$.
If it is attained only at $D$, then $g\in Bt_0^{-\ord D}B$. If it is
attained only at $B$, then $g\in Bt_0^{\ord B}s_1B$. If it is attained
only at $C$, then $g\in Bt_0^{-\ord C}s_1B$. For example, if the
minimum is attained only at $C$, then $g=u_+(A/C)\,w(-C^{-1})\,u_+(D/C)$
with $A/C$ and $D/C$ in $\cO$, and $w(-C^{-1})\in Bt_0^{-\ord
C}s_1B$ by (\ref{eq:uplus}). For $GL(2,F)$ the same operations, by
elements of the Iwahori subgroup of $GL(2,\cO)$, reduce $g$ to a
diagonal or antidiagonal matrix. Its image in $PGL(2,F)$ lies in
$I'\lambda^jI'$, respectively in $I'\lambda^j[s_1]I'$, where $j$ is
the valuation of the upper entry minus that of the lower entry.

\begin{proposition}
\label{prop:orbsl2}
Let $\gamma=h(a)$ with $m=|\ord(a)|\geq1$. For $w\in W_a$,
\[
F_{f_w}(\gamma)=
\begin{cases}
q^m & \mbox{if } w=t_0^{\pm m},\\
(q-1)\,q^{(\ell(w)-1)/2} & \mbox{if } w=t_0^is_1 \mbox{ and }
\ell(w)\geq 2m+1,\\
0 & \mbox{otherwise.}
\end{cases}
\]
\end{proposition}

\begin{proof}
Since $s_1h(a)s_1^{-1}=h(a^{-1})$, we may assume $\ord(a)=m$. We apply
Lemma~\ref{lemma:orbred}. For $\kappa=1$, $\kappa=s_1$ and
$\kappa=u_-(y)$ with $y$ a unit, the matrix
$\kappa^{-1}\gamma u_+(z)\kappa$ is
\[
\mat a{az}0{a^{-1}},\qquad \mat{a^{-1}}0{-az}a,\qquad
\mat{a(1+v)}{av/y}{y(a^{-1}-a(1+v))}{a^{-1}-av},
\]
where $v=yz$. The rule above gives the following, where $j$ denotes
$\ord z$, respectively $\ord v$, and where the last column is the
measure of the set of $z$.
\begin{center}
\begin{tabular}{@{}llll@{}}
\toprule
$\kappa$ & condition & double coset & measure \\
\midrule
$1$ & $j\geq -2m$ & $Bt_0^mB$ & $q^{2m}$ \\
$1$ & $j\leq -2m-1$ & $Bt_0^{m+j}s_1B$ & $(1-q^{-1})q^{-j}$ \\
$s_1$ & $j\geq 1-2m$ & $Bt_0^{-m}B$ & $q^{2m-1}$ \\
$s_1$ & $j\leq -2m$ & $Bt_0^{-m-j}s_1B$ & $(1-q^{-1})q^{-j}$ \\
$u_-(y)$ & $j>-2m$ & $Bt_0^ms_1B$ & $q^{2m-1}$ \\
$u_-(y)$ & $v\in a^{-2}(1+\varpi\cO)$ & $Bt_0^{-m}B$ & $q^{2m-1}$ \\
$u_-(y)$ & $j=-2m$, $v\notin a^{-2}(1+\varpi\cO)$ & $Bt_0^ms_1B$ &
$q^{2m}-2q^{2m-1}$ \\
$u_-(y)$ & $j<-2m$ & $Bt_0^{-m-j}s_1B$ & $(1-q^{-1})q^{-j}$ \\
\bottomrule
\end{tabular}
\end{center}
For instance, in the second row the minimum of $\nu$ is attained only
at the upper right entry $az$, of valuation $m+j$. In the sixth row
write $v=a^{-2}(1+\epsilon)$ with $\ord\epsilon\geq1$. Then the upper
left entry $a^{-1}(1+\epsilon+a^2)$ has valuation $-m$, while the
other entries have $\nu>-m$. The other rows are similar. Summing,
with $q-1$ units $y$, the total measure is $q^{2m}$ for $t_0^m$ and
for $t_0^{-m}$. For $t_0^is_1$ it is $(q-1)q^{m+i}$ if $i\geq m$,
$(1-q^{-1})q^{m-i}$ if $i\leq-m-1$, and $0$ otherwise. Multiplying by
$|1-a^{-2}|^{-1}=q^{-2m}$ and by $|D(\gamma)|^{1/2}=q^m$, and using
$\ell(t_0^is_1)=|2i+1|$, gives the stated values.
\end{proof}

\begin{proposition}
\label{prop:orbpgl}
Let $\gamma'=[\mathrm{diag}(b,1)]$ with $k=|\ord(b)|\geq1$. For $w$ in
the extended affine Weyl group of $G'$,
\[
F'_{f'_w}(\gamma')=
\begin{cases}
q^{k/2} & \mbox{if } w=\lambda^{\pm k},\\
(q-1)\,q^{(\ell'(w)-1)/2} & \mbox{if } w=\lambda^i[s_1],\ \ell'(w)\geq
k+1\mbox{ and } i\equiv k\pmod 2,\\
0 & \mbox{otherwise.}
\end{cases}
\]
\end{proposition}

\begin{proof}
The proof is that of Proposition~\ref{prop:orbsl2}, with
$|1-b^{-1}|^{-1}=q^{-k}$ in place of $|1-a^{-2}|^{-1}$ and with
$\ord(b)=k$. The three matrices are $\mat b{bz}01$, $\mat 10{-bz}b$
and $\mat{b(1+v)}{bv/y}{y(1-b(1+v))}{1-bv}$, and the table becomes the
following.
\begin{center}
\begin{tabular}{@{}llll@{}}
\toprule
$\kappa$ & condition & double coset & measure \\
\midrule
$1$ & $j\geq -k$ & $I'\lambda^kI'$ & $q^{k}$ \\
$1$ & $j\leq -k-1$ & $I'\lambda^{k+2j}[s_1]I'$ & $(1-q^{-1})q^{-j}$ \\
$s_1$ & $j\geq 1-k$ & $I'\lambda^{-k}I'$ & $q^{k-1}$ \\
$s_1$ & $j\leq -k$ & $I'\lambda^{-k-2j}[s_1]I'$ & $(1-q^{-1})q^{-j}$ \\
$u_-(y)$ & $j>-k$ & $I'\lambda^k[s_1]I'$ & $q^{k-1}$ \\
$u_-(y)$ & $v\in b^{-1}(1+\varpi\cO)$ & $I'\lambda^{-k}I'$ & $q^{k-1}$ \\
$u_-(y)$ & $j=-k$, $v\notin b^{-1}(1+\varpi\cO)$ & $I'\lambda^k[s_1]I'$
& $q^{k}-2q^{k-1}$ \\
$u_-(y)$ & $j<-k$ & $I'\lambda^{-k-2j}[s_1]I'$ & $(1-q^{-1})q^{-j}$ \\
\bottomrule
\end{tabular}
\end{center}
Only reflections $\lambda^i[s_1]$ with $i\equiv k\pmod 2$ occur, as
they must, since $\det$ modulo squares is a conjugation invariant.
Summing as before gives the stated values.
\end{proof}

\begin{example}
\label{ex:orbsl2}
Proposition~\ref{prop:orbsl2} is consistent with
Proposition~\ref{prop:spherical}. The function $\chi_{Kt_0K}$ is the
sum of $f_w$ over $w=t_0^{\pm1}$, $t_0s_1$ and $t_0^{-1}s_1$, and at
$\gamma=t_0$ the proposition gives $q+q+(q-1)q+0=q^2+q$. With Haar
measure normalized by $\vol(K)=1$ instead, this becomes $q$, the
coefficient of $X$ in the Satake transform
$q(X+1+X^{-1}-q^{-1})$ of Proposition~\ref{prop:spherical}(b). In the
running example $q=7$, $m=1$, the value of $F_{f_{t_0s_1}}(t_0)$ is
$42$.
\end{example}

The genuine orbital integrals are determined by these, through Savin's
isomorphism. We need two facts about the generators.

\begin{lemma}
\label{lemma:cocenter}
\begin{enumerate}[leftmargin=2em,label=(\alph*)]
\item Let $A$ be an algebra generated by two elements $\xi$ and $\eta$
that satisfy quadratic relations. Then $A$ is spanned by $1$, $\xi$,
$\eta$, the elements $(\xi\eta)^j$ with $j\geq1$, and the commutators
$\alpha\beta-\beta\alpha$ with $\alpha,\beta\in A$.
\item Let $p\nmid n$ and $k\geq1$, and put $m=n'k$. Then
$(XY)^k=\zeta'_k\,\fbar_{(t_0^m,1)}$, where $\zeta'_k=1$ if $n$ is odd
and $\zeta'_k=\epsilon^{k(k-1)/2}$ if $n$ is even, with
$\epsilon=(-1)^{(q-1)/2}$.
\end{enumerate}
\end{lemma}

\begin{proof}
(a) Using the relations, $A$ is spanned by alternating products of
$\xi$ and $\eta$. Modulo commutators, a product may be rotated
cyclically. So $(\eta\xi)^j\equiv(\xi\eta)^j$. An alternating product
of odd length at least $3$, say $\xi(\eta\xi)^j$, is congruent to
$(\eta\xi)^j\xi=(\eta\xi)^{j-1}\eta\xi^2$, which is a combination of
shorter alternating products by the relation for $\xi$. The same holds
with $\xi$ and $\eta$ exchanged. Induction on the length gives (a).

(b) By Lemma~\ref{lemma:genuine}, $(XY)^k=\fbar_L$, where $L$ is the
$k$-th power of $(s_1,1)(r,1)$. We have $s_1r=h(c)$ with
$c=-\varpi^{n'}$, and $x(\cdot)$ takes the values $1$, $\varpi^{n'}$
and $c^{-1}$ on $s_1$, $r$ and $h(c)$. So
$(s_1,1)(r,1)=(h(c),(c^{-1},c^{-1}\varpi^{-n'})_n)$. By (H1)--(H3) and
(H5),
\[
(c^{-1},c^{-1}\varpi^{-n'})_n=(c^{-1},-c^{-2})_n=(-1,c)_n^{-1}
=(-1,\varpi)_n^{-n'} .
\] By Lemma~\ref{lemma:cover}(a),
$(h(c),1)^k=(h(c^k),(c,c)_n^{k(k-1)/2})$, and
$(c,c)_n=(-1,c)_n=(-1,\varpi)_n^{n'}$. Finally
$(h(c^k),1)=(\varpi^m,(-1)^k)_n^{-1}\,(h((-1)^k),1)\,(t_0^m,1)$, and
$(h((-1)^k),1)=s(h((-1)^k))\in B^*$. Hence $(XY)^k=\zeta\,
\fbar_{(t_0^m,1)}$ with
\[
\zeta=(-1,\varpi)_n^{\,n'k-n'k(k-1)/2-mk},
\]
by Lemma~\ref{lemma:support}(b). The value $(-1,\varpi)_n$ is $\pm1$.
It is $1$ if $n$ is odd, and $(-1,\varpi)_n^{n'}=\epsilon$ if $n$ is
even, by Example~\ref{ex:hilbert}. Since $m=n'k$, the exponent is
congruent to $n'k(k-1)/2$ modulo $2n'$, which gives $\zeta=\zeta'_k$.
\end{proof}

\begin{theorem}[Savin's isomorphism transfers split orbital integrals]
\label{thm:orbtransfer}
Let $p\nmid n$, $k\geq1$ and $m=n'k$, and let $\Theta$, $G'$, $\mu$
and $a_n$ be as in Theorem~\ref{thm:savinsl2}. Put
$\tilde\gamma=(t_0^m,1)$, a good element. For every $\phi$ in
$\cH(\Gbar/\!/B^*)_{\gen}$,
\[
F_\phi(\tilde\gamma)=\Delta_k\,F'_{\Theta(\phi)}(\mu^k),\qquad
\Delta_k=\zeta'_k\Bigl(\frac{|a_n|}{a_n}\Bigr)^{k}.
\]
Thus $\Delta_k=1$ for odd $n$, and
$\Delta_k=\epsilon^{k(k-1)/2}\,\overline{\epsilon^{1/2}}^{\,k}$ for
even $n$. Here $F'=F$ when $n$ is odd.
\end{theorem}

\begin{proof}
Consider the functionals $\tau_1(\phi)=F_\phi(\tilde\gamma)$ and
$\tau_2(\phi)=F'_{\Theta(\phi)}(\mu^k)$. They vanish on commutators, by
Lemma~\ref{lemma:orbtrace} and since $\Theta$ is an algebra
isomorphism. By Lemma~\ref{lemma:cocenter}(a) and
Theorem~\ref{thm:presentation}(c), it suffices to compare them on
$1$, $X$, $Y$ and $(XY)^j$ for $j\geq1$.

Since $|\fbar_{(w,1)}|=f_w\circ\pr$, we have $|O_{\tilde\gamma}(\fbar_{(w,1)})|\leq
O_{t_0^m}(f_w)$. By Proposition~\ref{prop:orbsl2} this vanishes for
$w=1$, for $w=s_1$, for $w=r$, whose length $2n'-1$ is less than
$2m+1$, and for $w=t_0^{n'j}$ with $j\neq k$. By
Lemma~\ref{lemma:words}(a), $(XY)^j$ is a multiple of
$\fbar_{(t_0^{n'j},1)}$. So $\tau_1$ vanishes on $1$, $X$, $Y$ and
$(XY)^j$ for $j\neq k$. By the proof of
Proposition~\ref{prop:orbsl2}, the double coset $Bt_0^mB$ is reached
only from $\kappa=1$ and $\ord z\geq-2m$. There
$\tilde\gamma\tilde u_+(z)=\tilde u_+(\varpi^{2m}z)\tilde\gamma$, and
$\tilde u_+(\varpi^{2m}z)\in B^*$, so the integrand equals $1$. Hence
$F_{\fbar_{(t_0^m,1)}}(\tilde\gamma)=q^m\cdot q^{-2m}\cdot q^{2m}=q^m$,
and $\tau_1((XY)^k)=\zeta'_kq^m$ by Lemma~\ref{lemma:cocenter}(b).

On the other side, $\Theta(1)=1$, $\Theta(X)=f'_{s_1}$ and
$\Theta(Y)=a_nf'_{r'}$, and Propositions~\ref{prop:orbsl2} and
\ref{prop:orbpgl} show that $F'$ vanishes at $\mu^k$ on $f'_1$,
$f'_{s_1}$ and $f'_{r'}$. Also $\Theta((XY)^j)=a_n^j(f'_{s_1}*f'_{r'})^j
=a_n^jf'_{\mu^j}$, because $s_1r'=\pm\mu$ and lengths add. So
$\tau_2((XY)^j)$ vanishes for $j\neq k$, and
$\tau_2((XY)^k)=a_n^k|D'(\mu^k)|^{1/2}$. Since
$|a_n|^k|D'(\mu^k)|^{1/2}=q^m$ in both cases, $\tau_1=\Delta_k\tau_2$.
\end{proof}

\begin{corollary}
\label{cor:orbgen}
In the situation of Theorem~\ref{thm:orbtransfer}, let $w=t_0^{n'j}$
or $w=t_0^{n'j}s_1$. Then $F_{\fbar_{(w,1)}}(\tilde\gamma)$ has
absolute value
\[
\begin{cases}
q^m & \mbox{if } w=t_0^{\pm m},\\
(q-1)\,q^{(\ell(w)-1)/2} & \mbox{if } w=t_0^{n'j}s_1,\ \ell(w)\geq
2m+1 \mbox{ and } n\mid m+n'j,\\
0 & \mbox{otherwise.}
\end{cases}
\]
In particular these are the values of Proposition~\ref{prop:orbsl2},
except that for even $n$ the reflections $t_0^{n'j}s_1$ with $j\not\equiv
k\pmod 2$ give $0$.
\end{corollary}

\begin{proof}
By Theorem~\ref{thm:savinsl2}(b),
$\Theta(\fbar_{(w,1)})=c_wq^{e(w)}f'_{\Psi(w)}$ with $|c_w|=1$. By
Theorem~\ref{thm:orbtransfer},
\[
|F_{\fbar_{(w,1)}}(\tilde\gamma)|=q^{e(w)}\,|F'_{f'_{\Psi(w)}}(\mu^k)| .
\]
Apply Propositions~\ref{prop:orbsl2} and \ref{prop:orbpgl}, and use
$e(w)=\frac12(\ell(w)-\ell'(\Psi(w)))$. For $w=t_0^{n'j}s_1$, the
condition on $\ell'(\Psi(w))$ is equivalent to $\ell(w)\geq 2m+1$,
and for even $n$ the parity condition $j\equiv k\pmod 2$ is equivalent
to $n\mid m+n'j$. For odd $n$ the divisibility always holds.
\end{proof}

\begin{vremark}
\label{vrem:orbphases}
\emph{Computer check, and phases.} We computed the genuine orbital
integrals of Corollary~\ref{cor:orbgen} directly from
Lemma~\ref{lemma:orbred}, with exact Hilbert symbols. For $n=2$ we
took $q=3$ with $k=1,2$ and $q=5$ with $k=1$, and for $n=4$ we took
$q=5$ with $k=1$. In these cases we computed the translations and between
four and six reflections $t_0^{n'j}s_1$ with $|j|\leq 4$, on both sides
of the vanishing range. For $n=3$ and $q=7$ we
computed only the two translations $t_0^{\pm3}$, since the reflections
need too many points. The results agree with
Corollary~\ref{cor:orbgen} and with Theorem~\ref{thm:orbtransfer},
including the constant $\Delta_k$, which is $-i$ for $q=3$ and $k=1$. In every computed case the non-zero values
are positive, so the integrand has constant phase $1$ on each
contributing region, and the vanishing for even $n$ comes from a
non-trivial character sum over the units. We have checked the phases by
hand along the lines of the proof of
Proposition~\ref{prop:quadratic}, but we have not written out that
computation here.
\end{vremark}

\begin{remark}
\label{rem:orbmeaning}
Theorem~\ref{thm:orbtransfer} is a transfer of orbital integrals at
split elements, with the bijection $t_0^{n'k}\mapsto\mu^k$ and the
transfer factor $\Delta_k$. It explains the factors $q^{e(w)}$ of
Theorem~\ref{thm:savinsl2}(b). The normalized orbital integrals of
$\fbar_{(t_0^m,1)}$ at $(t_0^m,1)$ and of $f'_{\mu^k}$ at $\mu^k$ are
$|D(t_0^m)|^{1/2}=q^m$ and $|D'(\mu^k)|^{1/2}$, and their ratio is the
factor $q^{e(w)}$. For $n=2$ and $q\equiv 3\pmod4$ the transfer factor
$\Delta_1=-i$ is not $1$. This is consistent with the expectation that
the transfer for $Mp(2)$ involves roots of unity \cite{Li11}. We have not
compared $\Delta_k$ with the transfer factors of \cite{Li11}.
\end{remark}

\medskip\noindent\textbf{Deep level.} For the basis of
Proposition~\ref{prop:admissible} we have the following.

\begin{proposition}
\label{prop:orbdeep}
Let $k\geq1$ be deep, normalize Haar measure by $\vol(B_k)=1$, and let
$\gamma=h(a)$ with $m=|\ord(a)|\geq1$.
\begin{enumerate}[leftmargin=2em,label=(\alph*)]
\item Let $B_kgB_k\subset BwB$ with $w\in W_a$. If $w$ is neither
$t_0^{\pm m}$ nor a reflection $t_0^is_1$ with $\ell(w)\geq 2m+1$, then
$O_\gamma(f_g)=0$, and $O_{\tilde\gamma}(\fbar_{\tilde g})=0$ for
every good lift $\tilde\gamma$ of $\gamma$.
\item For every $w\in W_a$,
\[
\sum_gF_{f_g}(\gamma)=(q-1)q^{3k-1}\,F_{f_w}(\gamma),
\]
where $g$ runs over representatives of the double cosets $B_kgB_k$
contained in $BwB$, and $F_{f_w}(\gamma)$ is the value in
Proposition~\ref{prop:orbsl2}.
\item If $n\nmid 2m'$ and $\tilde t$ is a lift of $t_0^{m'}$, then
$D(\fbar_{\tilde t})=0$ for every distribution $D$ on $\Gbar$ that is
invariant under conjugation. In particular every genuine orbital
integral of $\fbar_{\tilde t}$, at every good element, vanishes.
\end{enumerate}
\end{proposition}

\begin{proof}
(a) We have $0\leq f_g\leq\chi_{BwB}$ and
$|\fbar_{\tilde g}|=f_g\circ\pr$, and orbital integrals are integrals
of non-negative functions. So both orbital integrals are bounded by
the orbital integral of $\chi_{BwB}$, which vanishes by
Proposition~\ref{prop:orbsl2}. (b) The characteristic function of
$BwB$ is the sum of the $f_g$. With $\vol(B_k)=1$ the Haar measure is
$[B:B_k]=(q-1)q^{3k-1}$ times the measure with $\vol(B)=1$. Here
$[B:B_1]=(q-1)q^2$, since $B/B_1$ is described by the residues of the
diagonal entries, which are inverse units, and of the two off-diagonal
entries, and $[B_1:B_k]=q^{3k-3}$ by Lemma~\ref{lemma:subgroups}(b).
Orbital integrals scale in the same way. (c) Choose a unit $u$ with $\zeta=(u,\varpi)_n^{2m'}\neq1$, by
Lemma~\ref{lemma:hilbert}, and a lift $\tilde h$ of $h(u)$. Conjugation
by $\tilde h$ preserves $B_k^*$, by Lemma~\ref{lemma:canonical}(b),
and $\tilde h^{-1}\tilde t\tilde h=\zeta^{\pm1}\tilde t$ by
Lemma~\ref{lemma:cover}(a). So $\phi\mapsto\phi(\tilde
h\,\cdot\,\tilde h^{-1})$ maps $\fbar_{\tilde t}$ to
$\zeta^{\pm1}\fbar_{\tilde t}$, and $D(\fbar_{\tilde t})=\zeta^{\pm1}
D(\fbar_{\tilde t})$.
\end{proof}

Part (c) strengthens Proposition~\ref{prop:traceless}, since traces of
admissible representations are invariant distributions. The individual
orbital integrals $O_\gamma(f_g)$ at deep level we have not computed.

\begin{vremark}[What remains heuristic]
\label{vrem:savinaction}
\emph{Heuristic.} At deep level we expect
the comparison of orbital integrals to involve the same power map
$t_0^{n'k}\leftrightarrow\mu^k$ on split elements, as in the orbit map
of \cite{FK86}. The genuine orbital integrals live only on good
elements. For $\gamma=h(a)$ the commutator with $h(b)$ is $(b,a)_n^2$
by Lemma~\ref{lemma:cover}(a), so only those $a$ with
$a^2\in(F^\times)^n$ occur. For $n=2$ one expects the transfer factors
to be those of \cite{Li11}, which Remark~\ref{rem:orbmeaning} does not
check.
\end{vremark}

\subsection{Elliptic orbital integrals}
\label{subsec:elliptic}
This subsection treats regular elliptic elements at Iwahori level. On
the linear side the orbital integrals become counts of edges of the
Bruhat--Tits tree. We compute them for every basis function of
$\cH(G/\!/B)$ and of $\cH(G'/\!/I')$
(Propositions~\ref{prop:ellsl2} and \ref{prop:ellpgl}). They vanish on
translations, and they are determined by three numbers attached to the
subtree fixed by the element. On the genuine side the same reductions
hold (Proposition~\ref{prop:ellgen}), but the counts carry roots of
unity, and we have computed the remaining three values only
numerically. If the computed values hold in general, Savin's
isomorphism transfers elliptic orbital integrals as well
(Theorem~\ref{thm:ellcond}). For odd $n$ it matches each elliptic
class with itself. For even $n$ the statement has the shape of an
endoscopic transfer: sums over the two classes in a stable class match
$PGL(2)$, and differences match elements of $PGL(2)$ that flip an
edge.

\medskip\noindent\textbf{The tree.} Let $\mathcal T$ be the
Bruhat--Tits tree of $SL(2,F)$. Its vertices are the homothety classes
of $\cO$-lattices in $F^2$, and two vertices are adjacent if they have
representatives $L\supset L'\supset\varpi L$ with $L/L'$ of order $q$.
The group $G$ has two orbits on vertices. Vertices of type $0$ are
those in the orbit of $[\cO^2]$, with stabilizers conjugate to $K$, and
vertices of type $1$ are those in the orbit of $[\cO\oplus\varpi\cO]$.
The group $G$ acts transitively on edges, and the stabilizer of the
standard edge $C=\{[\cO^2],[\cO\oplus\varpi\cO]\}$ is $B$. So $x\mapsto
xC$ identifies $G/B$ with the set of edges. For edges $D=xC$ and
$D'=yC$, the \textbf{relative position} $\delta(D,D')$ is the $w\in
W_a$, up to sign, with $x^{-1}y\in BwB$. A reflection $t_0^is_1$ fixes
a vertex of type $0$ if $i$ is even and of type $1$ if $i$ is odd, and
its length is $2d+1$ with $d=i$ for $i\geq0$ and $d=-i-1$ for $i<0$.
Distances are measured in edges. On $G'=PGL(2,F)$ the set $G'/I'$ is
the set of directed edges, and the relative position $\delta'$ of two
directed edges lies in the extended affine Weyl group. Here $f'_{[s_1]}$
corresponds to directed edges with the same tail and different heads,
and $f'_{[t]}$ to an edge and its reverse.

An element $\gamma\in G$ is \textbf{elliptic regular} if its
characteristic polynomial is irreducible over $F$. Then $E=F[\gamma]$ is
a quadratic extension, the centralizer $T_\gamma\cong E^1$ is compact,
and we normalize $\vol(T_\gamma)=1$. The eigenvalues are units, so
$\gamma\in\cO_E$, and $\cO[\gamma]=\cO+\varpi^r\cO_E$ for a unique
$r\geq0$, the \textbf{conductor} of $\gamma$. Also
$|D(\gamma)|^{1/2}=|\gamma-\bar\gamma|$, which is $q^{-r}$ if $E/F$ is
unramified and $q^{-r-1/2}$ if $E/F$ is ramified. On $G'$ we use the
centralizer $E^\times/F^\times$, also with volume $1$.

\begin{lemma}
\label{lemma:ellcount}
Let $\gamma\in G$ be elliptic regular.
\begin{enumerate}[leftmargin=2em,label=(\alph*)]
\item For every bi-$B$-invariant $f$,
$O_\gamma(f)=\sum_Df(x_D^{-1}\gamma x_D)$, where $D=x_DC$ runs over the
edges of $\mathcal T$. In particular
$O_\gamma(f_w)=\#\{D\ |\ \delta(D,\gamma D)=w\}$.
\item If $p\nmid n$, every lift $\tilde\gamma$ of $\gamma$ is good, and
$O_{\tilde\gamma}(\phi)=\sum_D\phi(\tilde x_D^{-1}\tilde\gamma\tilde
x_D)$ for genuine bi-$B^*$-invariant $\phi$.
\end{enumerate}
The same holds on $G'$, with directed edges.
\end{lemma}

\begin{proof}
(a) The set $T_\gamma\backslash G/B$ is the set of $T_\gamma$-orbits
of edges. The set $T_\gamma\backslash T_\gamma xB$ has measure
\[
\vol(xB)/\vol(T_\gamma\cap xBx^{-1}),
\]
which is the size of the $T_\gamma$-orbit of $xC$, since $\vol(T_\gamma)=1$ and $xBx^{-1}$ is the
stabilizer of $xC$. Summing over orbits gives a sum over edges. It is
finite, by Lemma~\ref{lemma:fixedtree} and the proof of
Proposition~\ref{prop:ellsl2}. (b) The commutator of lifts is a
bimultiplicative pairing on $T_\gamma$ with values in $\mu_n$, and it is
trivial on the diagonal. The group $E^1$ is the product of a finite
cyclic group, of order $q+1$ or $2$, and a pro-$p$ group. A
bimultiplicative pairing into $\mu_n$ is trivial on a pro-$p$ factor,
since $p\nmid n$, and an alternating pairing on a cyclic group is
trivial. So every lift is good, and the proof of (a) applies.
\end{proof}

\begin{lemma}
\label{lemma:fixedtree}
Let $\gamma\in G$ be elliptic regular with conductor $r$. The vertices
fixed by $\gamma$ are the vertices at distance at most $r$ from a
vertex $c$ if $E/F$ is unramified, and from an edge $c$ if $E/F$ is
ramified.
\end{lemma}

\begin{proof}
Since $\det\gamma=1$, a vertex $[L]$ is fixed if and only if $\gamma
L=L$, that is, if $L$ is a module over $R_r=\cO[\gamma]$. Choose
$v\neq0$ in $F^2$ and identify $F^2$ with $E$ by $\xi\mapsto\xi v$, with
$\gamma$ acting by multiplication. For a lattice $L\subset E$, the ring
$\{\xi\in E\ |\ \xi L\subset L\}$ is an order, hence equal to
$R_c=\cO+\varpi^c\cO_E$ for some $c\geq0$. Since $R_c=\cO[\varpi^c\delta]$
for a generator $\delta$ of $\cO_E$ over $\cO$, it is a Gorenstein local
ring, so $L$ is an invertible, hence principal, $R_c$-module:
$L=yR_c$ with $y\in E^\times$. So the fixed vertices are the classes
$[yR_c]$ with $c\leq r$. Let $c_0$ be the vertex $[\cO_E]$ if $E/F$ is
unramified, and the edge $\{[\cO_E],[\varpi_E\cO_E]\}$ if it is
ramified. The vertex $[R_c]$ lies at distance $c$ from $c_0$, since
$\cO_E/R_c\cong\cO/\varpi^c\cO$. The group $E^\times$ fixes $c_0$, so
all $[yR_c]$ lie at distance $c$. Their number is
$[E^\times:F^\times R_c^\times]$, which is $(q+1)q^{c-1}$ for $c\geq1$
in the unramified case and $2q^c$ in the ramified case. These are the
numbers of vertices at distance exactly $c$ from $c_0$. Hence the fixed
vertices are exactly the vertices at distance at most $r$ from $c_0$.
\end{proof}

Let $\mathcal F_\gamma$ be this fixed subtree. Put $E_\gamma$ for its
number of edges, and $S_\tau(\gamma)$ for the number of edges of
$\mathcal T$ that have one endpoint in $\mathcal F_\gamma$, of type
$\tau$, and the other endpoint outside. By
Lemma~\ref{lemma:fixedtree}, if $E/F$ is unramified and $c$ has type
$\tau_c$, then
\[
E_\gamma=\frac{(q+1)(q^r-1)}{q-1},\qquad S_{\tau_c+r}(\gamma)=(q+1)q^r,
\qquad S_{\tau_c+r+1}(\gamma)=0,
\]
with types taken modulo $2$. If $E/F$ is ramified, then
$E_\gamma=1+2(q+q^2+\cdots+q^r)$ and $S_0(\gamma)=S_1(\gamma)=q^{r+1}$.

\begin{proposition}
\label{prop:ellsl2}
Let $\gamma\in G$ be elliptic regular, and $w\in W_a$. Then
$O_\gamma(f_w)$ is $E_\gamma$ if $w=1$, and $0$ if $w=t_0^j$ with
$j\neq0$. If $w$ is a reflection of length $2d+1$ fixing a vertex of
type $\tau$, then $O_\gamma(f_w)=q^dS_\tau(\gamma)$.
\end{proposition}

\begin{proof}
We use Lemma~\ref{lemma:ellcount}(a). An edge $D$ with
$\delta(D,\gamma D)=1$ is an edge fixed by $\gamma$, that is, an edge
of $\mathcal F_\gamma$. Let $D$ be an edge not in $\mathcal
F_\gamma$. Let $v$ be the vertex of $\mathcal F_\gamma$ nearest to
$D$, and let $d$ be the distance from $v$ to the nearer endpoint of
$D$. The first edge of the path from $v$ to $D$ is not fixed, so the
geodesic from $D$ to $\gamma D$ passes through $v$. In an apartment
containing this geodesic, $\gamma D$ is the mirror image of $D$ in $v$,
because $\gamma$ preserves types and fixes $v$. Hence $\delta(D,\gamma
D)$ is the reflection of length $2d+1$ fixing a vertex of the type of
$v$. In particular it is never a translation. The vertex $v$ has
$q+1$ neighbors, and the edges from $v$ to vertices outside $\mathcal
F_\gamma$ are counted in $S_\tau(\gamma)$. Each such edge starts a
branch with $q^d$ edges whose nearer endpoint is at distance $d$ from
$v$. This gives $q^dS_\tau(\gamma)$.
\end{proof}

\begin{proposition}
\label{prop:ellpgl}
Let $\gamma'\in G'$ be elliptic regular.
\begin{enumerate}[leftmargin=2em,label=(\alph*)]
\item If $\gamma'$ fixes a vertex, its fixed vertices again form a
subtree $\mathcal F_{\gamma'}$, and $O'_{\gamma'}(f'_1)=2E_{\gamma'}$,
$O'_{\gamma'}(f'_{[s_1]})=S_0(\gamma')+S_1(\gamma')$,
$O'_{\gamma'}(f'_{[t]})=0$. For $\gamma'$ the image of an elliptic
regular $\gamma\in G$, the subtree $\mathcal F_{\gamma'}$ equals
$\mathcal F_\gamma$.
\item If $\gamma'$ interchanges the endpoints of exactly one edge, then
$O'_{\gamma'}(f'_1)=O'_{\gamma'}(f'_{[s_1]})=0$ and
$O'_{\gamma'}(f'_{[t]})=2$.
\item In both cases $O'_{\gamma'}(f'_{\lambda^k})=0$ for $k\neq0$.
\end{enumerate}
\end{proposition}

\begin{proof}
An elliptic $\gamma'$ fixes a point of the geometric realization of
$\mathcal T$, either a vertex or the midpoint of an edge. As in the
proof of Proposition~\ref{prop:ellsl2}, a directed edge $D$ and
$\gamma'D$ are mirror images in the fixed point nearest to $D$, so
$\delta'(D,\gamma'D)$ is never a translation. This gives (c). In (a),
a directed edge is fixed if and only if its edge lies in $\mathcal
F_{\gamma'}$, which gives $2E_{\gamma'}$. The relative position is
$[s_1]$ exactly when the tail is fixed and the head is not, which gives
the edges counted in $S_0+S_1$. No edge is reversed, since a reversal
does not fix a vertex. If $\gamma'$ is the image of $\gamma$, it fixes
the same vertices. In (b), reversing an edge forces $\ord\det\gamma'$
to be odd, so $\gamma'$ moves every vertex to a vertex at odd distance
and fixes no vertex. A directed edge $D$ has
$\delta'(D,\gamma'D)=[t]$ exactly when $\gamma'D$ is $D$ reversed, which
happens for the two orientations of the one reversed edge.
\end{proof}

All these functionals are traces on algebras with two quadratic
generators that vanish on the translations. Such traces are determined
by three values.

\begin{lemma}
\label{lemma:ellrec}
Let $A$ be generated by $\xi$ and $\eta$ with $\xi^2=\alpha\xi+\beta$
and $\eta^2=\gamma_0\eta+\delta_0$, and let $\tau$ be a linear form on
$A$ that vanishes on commutators and on $(\xi\eta)^j$ for $j\geq1$. Put
$A_j=\tau(\xi(\eta\xi)^j)$ and $B_j=\tau(\eta(\xi\eta)^j)$. Then
$A_j=\beta B_{j-1}$ and $B_j=\delta_0A_{j-1}$ for $j\geq1$. So $\tau$ is
determined by $\tau(1)$, $\tau(\xi)$ and $\tau(\eta)$.
\end{lemma}

\begin{proof}
Rotating, $\tau(\xi(\eta\xi)^j)=\tau((\eta\xi)^{j-1}\eta\xi^2)
=\alpha\tau((\eta\xi)^j)+\beta\tau(\eta(\xi\eta)^{j-1})$, and
$\tau((\eta\xi)^j)=\tau((\xi\eta)^j)=0$. The other formula is the same
with $\xi$ and $\eta$ exchanged. The last statement follows from
Lemma~\ref{lemma:cocenter}(a).
\end{proof}

For $\cH(G/\!/B)$, with $\beta=\delta_0=q$, this recovers the factor
$q^d$ of Proposition~\ref{prop:ellsl2}. For the genuine algebra
$\beta=q$ and $\delta_0=c_0$, by Proposition~\ref{prop:quadratic}.

\begin{proposition}
\label{prop:ellgen}
Let $p\nmid n$, let $\gamma\in G$ be elliptic regular, and put
$\tilde\gamma=(\gamma,1)$ and $\bar\tau(\phi)=O_{\tilde\gamma}(\phi)$.
Let $\tau_r\in\{0,1\}$ be the type of the vertex fixed by $r$, so that
$\tau_r\equiv n'\pmod2$.
\begin{enumerate}[leftmargin=2em,label=(\alph*)]
\item $\bar\tau(\fbar_{(w,1)})=0$ for every translation $w\neq\pm1$, and
$\bar\tau$ is determined by $\bar\tau(1)$, $\bar\tau(X)$ and
$\bar\tau(Y)$ through Lemma~\ref{lemma:ellrec}.
\item $|\bar\tau(1)|\leq E_\gamma$, $|\bar\tau(X)|\leq S_0(\gamma)$ and
$|\bar\tau(Y)|\leq q^{n'-1}S_{\tau_r}(\gamma)$.
\end{enumerate}
\end{proposition}

\begin{proof}
Since $|\fbar_{(w,1)}|=f_w\circ\pr$, we have
$|\bar\tau(\fbar_{(w,1)})|\leq O_\gamma(f_w)$. Apply
Proposition~\ref{prop:ellsl2}, with $\ell(r)=2(n'-1)+1$. The trace
property is Lemma~\ref{lemma:orbtrace}, and Lemma~\ref{lemma:ellrec}
applies by Theorem~\ref{thm:presentation}(c).
\end{proof}

\begin{vremark}
\label{vrem:ellcomputed}
\emph{Computed, not proved.} Let $h\in GL(2,F)$ with $\det h\notin
N_{E/F}(E^\times)$, so that $\gamma$ and $\gamma^h=h^{-1}\gamma h$ are
the two $SL(2,F)$-classes in the stable class of $\gamma$. We computed
$\bar\tau(1)$, $\bar\tau(X)$ and $\bar\tau(Y)$ for both classes,
directly from Kubota's cocycle, for $\gamma=\mat 0{-1}1t$ with
conductor $0$ and $1$, in the ramified and unramified cases. We took
$q=3,5,7$ with $n=2$, $q=7$ with $n=3$, and $q=5$ with $n=4$. In every
case:
\begin{enumerate}[leftmargin=2em,label=(\roman*)]
\item $\bar\tau(1)=E_\gamma$ and $\bar\tau(X)=S_0(\gamma)$.
\item for odd $n$, $\bar\tau(Y)=q^{n-1}S_1(\gamma)$, so that all phases
are $1$.
\item for even $n$ and $E/F$ unramified, $\bar\tau(Y)=0$, although the
corresponding count $q^{n'-1}S_{\tau_r}(\gamma)$ can be non-zero.
\item for even $n$ and $E/F$ ramified, $\bar\tau(Y)=\pm
q^{n'-1}S_{\tau_r}(\gamma)$, with opposite signs for $\gamma$ and
$\gamma^h$.
\end{enumerate}
For example, for $q=3$, $n=2$ and $t=5$ (ramified, conductor $0$), the
three values are $(1,3,-3)$ for $\gamma$ and $(1,3,3)$ for
$\gamma^h$. For $q=7$, $n=3$ and $t=296$ (unramified, conductor $1$),
they are $(8,56,0)$ and $(8,0,2744)$.
\end{vremark}

\begin{theorem}[Elliptic transfer, conditional on the computed values]
\label{thm:ellcond}
Let $p\nmid n$, let $\gamma\in G$ be elliptic regular, and suppose that
(i)--(iv) of Remark~\ref{vrem:ellcomputed} hold for $\gamma$ and
$\gamma^h$. Let $\Theta$ and $a_n$ be as in
Theorem~\ref{thm:savinsl2}. Then for every $\phi$ in the genuine
Iwahori--Hecke algebra:
\begin{enumerate}[leftmargin=2em,label=(\alph*)]
\item If $n$ is odd, $O_{(\gamma,1)}(\phi)=O_\gamma(\Theta(\phi))$.
\item If $n$ is even,
$O_{(\gamma,1)}(\phi)+O_{(\gamma^h,1)}(\phi)=O'_{[\gamma]}(\Theta(\phi))$,
where $[\gamma]$ is the image of $\gamma$ in $PGL(2,F)$.
\item If $n$ is even and $E/F$ is ramified, and if $\gamma''\in G'$ is
elliptic regular and interchanges the endpoints of exactly one edge,
then
\[
F_{(\gamma,1)}(\phi)-F_{(\gamma^h,1)}(\phi)=
\pm\,\overline{\epsilon^{1/2}}\;F'_{\Theta(\phi)}(\gamma''),
\]
with $F=|D|^{1/2}O$ and $\epsilon^{1/2}$ as in
Theorem~\ref{thm:savinsl2}. The sign depends on which class is called
$\gamma$.
\end{enumerate}
\end{theorem}

\begin{proof}
Both sides of each identity are traces that vanish on translations, by
Lemma~\ref{lemma:orbtrace} and Propositions~\ref{prop:ellsl2},
\ref{prop:ellpgl} and \ref{prop:ellgen}. By Lemma~\ref{lemma:ellrec} it
suffices to compare them on $1$, $X$ and $Y$, where
$\Theta(1)=1$, $\Theta(X)=f'_{s_1}$ and $\Theta(Y)=a_nf'_{r'}$.

(a) By Proposition~\ref{prop:ellsl2}, $O_\gamma$ takes the values
$E_\gamma$, $S_0(\gamma)$ and $q^{n-1}S_1(\gamma)$ on $1$, $f_{s_1}$ and
$q^{n-1}f_{s_2}$, which are the values in (i) and (ii).

(b) By (i), (iii) and (iv), the left side takes the values
$2E_\gamma$, $S_0(\gamma)+S_0(\gamma^h)$ and $0$ on $1$, $X$ and $Y$.
Conjugation by $h$ maps $\mathcal F_\gamma$ onto $\mathcal
F_{\gamma^h}$. If $E/F$ is unramified, $\ord\det h$ is odd, so $h$
exchanges the two types, and $S_0(\gamma^h)=S_1(\gamma)$. If $E/F$ is
ramified, $S_0=S_1$ for both. In both cases the values are $2E_\gamma$,
$S_0(\gamma)+S_1(\gamma)$ and $0$. These are the values of
$O'_{[\gamma]}\circ\Theta$, by Proposition~\ref{prop:ellpgl}(a), since
$\mathcal F_{[\gamma]}=\mathcal F_\gamma$ and $\Theta(Y)$ is a multiple
of $f'_{[t]}$.

(c) By (i) and (iv), the left side vanishes on $1$ and $X$, and on $Y$
it equals $\pm2|D(\gamma)|^{1/2}q^{n'-1}q^{r+1}=\pm2q^{n'-1/2}$. By
Proposition~\ref{prop:ellpgl}(b), the right side vanishes on $1$ and
$X$. On $Y$ it is
\[
\pm\overline{\epsilon^{1/2}}\,|D'(\gamma'')|^{1/2}a_n\cdot2 .
\]
Here $|D'(\gamma'')|=|(1-\alpha)(1-\alpha^{-1})|=1$, where
$\alpha$ is the ratio of the eigenvalues, because $\alpha$ is close to
$-1$ and $p$ is odd. Since $\overline{\epsilon^{1/2}}a_n=q^{n'-1/2}$,
the two sides agree.
\end{proof}

\begin{remark}
\label{rem:ellmeaning}
For odd $n$, Theorem~\ref{thm:ellcond} matches each elliptic class with
itself, and together with Theorem~\ref{thm:orbtransfer} it would make
Savin's isomorphism a transfer of all regular orbital integrals at
Iwahori level. For even $n$ it matches the stable orbital integral, the
sum over the two classes, with $PGL(2)$. For ramified tori it matches
the difference with elements that reverse an edge. For unramified tori
the difference is not seen by $PGL(2)$, which has no such elements
there, and it should correspond to an endoscopic group. This is the
shape of the transfer for $Mp(2)$ in \cite{Li11}. The constant
$\pm\overline{\epsilon^{1/2}}$ is, up to sign, the constant $\Delta_1$
of Theorem~\ref{thm:orbtransfer}. We have not proved (i)--(iv) in
general, and we have not compared the constants with \cite{Li11}.
\end{remark}

%======================================================================
\section{Summary of results and open questions}
\label{section:summary}
\addtocontents{toc}{\protect\setcounter{tocdepth}{1}}

This section collects the main results and the questions they leave
open.

\subsection*{Results}
Let $G=SL(2,F)$ and let $\Gbar$ be its $n$-fold metaplectic cover.
\begin{enumerate}[leftmargin=2em]
\item The algebra $\Hlin$ is generated by the $f_g$ with $g$ in
$\{s_1,s_2\}\cup B$, subject to the relations of
Proposition~\ref{prop:main}. The element $f_{t_0}$ is not invertible
up to a scalar (Example~\ref{ex:t0}).
\item At deep level, $B_k$ has a canonical lift, and every double coset
$B_kgB_k$ carries a genuine function $\fbar_{\tilde g}$, unique up to a
root of unity. Hence $\Phi:f_g\mapsto\fbar_{\tilde g}$ is an
isomorphism of vector spaces from $\Hlin$ onto $\Hgen$, natural up to
these roots of unity (Proposition~\ref{prop:admissible}). It is
support-preserving (Definition~\ref{def:supportpreserving}), and every
support-preserving map is $\Phi$ followed by a rescaling of the basis.
\item Every relation of Proposition~\ref{prop:main} transfers to
$\Hgen$ for suitable lifts (Proposition~\ref{prop:main2}). When
$p\nmid n$, the part supported on $SL(2,\cO)$ transfers exactly
(Proposition~\ref{prop:Ksub}).
\item For $n=2$ and $p$ odd the torus gives no obstruction, and the
only sign in the relations involving $t_0$ and $s_1$ is $(-1,\varpi)_2$
(Lemma~\ref{lemma:signs}). In the running example it is $-1$.
\item For every $n\geq 2$, no support-preserving homomorphism with
constants of absolute value one exists, because the cover does not
split (Theorem~\ref{thm:unimodular} and Lemma~\ref{lemma:nonsplit}).
\item For $n\geq 3$, no support-preserving homomorphism exists for any
non-zero constants, and the trace identities fail
(Theorem~\ref{thrm:main}, Corollary~\ref{corollary:main} and
Example~\ref{ex:trace}). In the running example with $n=3$ this
happens at every level $k\geq 1$.
\item At Iwahori level with $p\nmid n$, the double cosets $Bt_0^mB$
and $Bt_0^ms_1B$ carry genuine functions if and only if $n\mid 2m$
(Proposition~\ref{prop:level0} and Lemma~\ref{lemma:level0s1}). The
genuine Iwahori--Hecke algebra has an explicit presentation
(Theorem~\ref{thm:presentation}), which gives a direct proof of
Savin's theorem for $SL(2)$ (Theorem~\ref{thm:savinsl2}). For odd $n$
the target is the Iwahori--Hecke algebra of $SL(2)$, with the
translations rescaled by $n$. For even $n$ it is that of $PGL(2)$, and
no isomorphism with that of $SL(2)$ respects the affine generator
(Corollary~\ref{cor:wronggroup}). The orbital integrals of all basis
functions at split elements are explicit, and Savin's isomorphism
transfers them (Theorem~\ref{thm:orbtransfer}). The linear elliptic
orbital integrals are explicit too, and the genuine ones are reduced to
three values, whose computed values imply an elliptic transfer
(Theorem~\ref{thm:ellcond}). At deep
level the translations $t_0^m$ with $n\nmid 2m$ carry genuine
functions, but these have trace zero in every admissible genuine representation
(Proposition~\ref{prop:traceless}), and their spherical parts vanish
(Proposition~\ref{prop:spherical}). The spherical parts of the linear
basis functions have explicit Satake transforms.
\end{enumerate}

\noindent
In short, the comparison with $SL(2)$ is correct for supports, for
volumes, for each relation separately and for everything supported on
$SL(2,\cO)$. It fails when the relations are combined. For $n=2$ the
failure is global and is the non-splitting of the cover. For $n\geq 3$
it is local, and it is concentrated on the translations $t_0^m$ with
$n\nmid 2m$, which are the ones that Savin's theorem removes
(Table~\ref{tab:howfar}).

\subsection*{Open questions}
Question~\ref{q:defining} asks whether the relations of
Proposition~\ref{prop:main} are a defining set. The following questions
concern the genuine algebra.

\begin{question}
\label{q:n2}
Let $n=2$, $p$ odd and $k$ deep. Is there an algebra homomorphism
$\Psi:\Hlin\to\Hgen$ with $\Psi(f_g)=c_g\fbar_{\tilde g}$ for every
representative $g$ of a double coset $B_kgB_k$, with non-zero constants
$c_g$? By Theorem~\ref{thm:unimodular}
such constants could not all have absolute value one. Is $\Hlin$
isomorphic to $\Hgen$ as an abstract algebra? The situation at Iwahori
level, where the genuine algebra is of $PGL(2)$ type, suggests that the
answer to the first question is negative. We have not proved this. For
the second question, note that at Iwahori level the corresponding
abstract isomorphism exists for trivial reasons
(Remark~\ref{rem:abstract}).
\end{question}

\begin{question}
\label{q:savin}
For deep $k$, is $\Hgen$ isomorphic to a generalized Hecke algebra of
the linear group $G/Z_n$, or to a twisted form of one, compatibly with
supports up to a rescaling of the lattice of translations, in the
spirit of Savin's theorem at level zero?
\end{question}

\begin{question}
\label{q:reps}
For deep $k$, is there a natural correspondence between the
irreducible genuine representations of $\Gbar$ with $B_k^*$-fixed
vectors and irreducible representations of a linear group? Does it
preserve supercuspidality and square-integrability?
\end{question}

\begin{question}[Orbital integrals]
\label{q:transfer}
Theorem~\ref{thm:orbtransfer} shows that Savin's isomorphism of
Theorem~\ref{thm:savinsl2} transfers orbital integrals at the split
elements $(t_0^{n'k},1)$. Do the values (i)--(iv) of
Remark~\ref{vrem:ellcomputed} hold for all elliptic regular elements, so
that the elliptic transfer of Theorem~\ref{thm:ellcond} holds
unconditionally? How are the constants $\Delta_k$ and
$\pm\overline{\epsilon^{1/2}}$ related to the transfer factors of
\cite{Li11} for $n=2$? For
$n=2$ this would be an Iwahori-level analog of the spherical
fundamental lemma for $Mp(2)$ \cite{Li11,Lu18}. For $n\geq 3$ we do not
know of a spherical fundamental lemma for these covers. Is there an
analog at deep level?
\end{question}

%======================================================================
\section*{Glossary}

\begin{description}[leftmargin=1.5em,style=nextline]
\item[Admissible element] An element $g$ of $G$ is admissible, relative
to a lift $K_0^*$ of an open compact subgroup $K_0$, if conjugation by a
lift of $g$ carries the lift of $K_0\cap g^{-1}K_0g$ onto the lift of
$gK_0g^{-1}\cap K_0$. The condition depends only on the double coset
$K_0gK_0$, and it holds exactly when that double coset carries a
non-zero genuine function that is invariant on both sides under
$K_0^*$. At deep level every element of $G$ is admissible for the
canonical lift of $B_k$, whereas at Iwahori level $t_0$ is not
admissible when $n\geq 3$.
\item[Affine Weyl group] For $SL(2,F)$ this is the infinite dihedral
group generated by two reflections, represented here by $s_1$ and $s_2$
up to the sign $\pm 1$. It is the semidirect product of the finite Weyl
group of order $2$ with the group of translations generated by
$t_0=h(\varpi)$. Its elements index the double cosets of the Iwahori
subgroup in $SL(2,F)$.
\item[Bernstein presentation] A description of an Iwahori--Hecke algebra
as generated by a finite Hecke algebra and a commutative algebra of
normalized translations, with an explicit commutation rule between
them. The translation elements are characteristic functions of double
cosets scaled by $q^{-\ell/2}$, where $\ell$ is the length. Savin's
isomorphism is defined on these generators. For $SL(2)$ the relation is
given in Remark~\ref{rem:bernstein}.
\item[Canonical lift] A distinguished splitting of a sufficiently small
open subgroup of $G$ into the cover, compatible with restriction to
smaller subgroups and with conjugation by every element of the cover.
When $p\nmid n$ its image is the unique Sylow pro-$p$ subgroup of the
inverse image. When $p\mid n$ it is built on deep congruence subgroups
by taking $n$-th roots of matrices with the logarithm and exponential
series, lifting them, and raising the lifts to the $n$-th power.
\item[Deep level] A level $k\geq 1$ at which the filtration subgroup
$B_k$, the principal congruence subgroup $\Gamma_k$ and its conjugate by
$t$ all have canonical lifts. Every $k\geq 1$ is deep when $p\nmid n$.
When $p\mid n$ the deep levels are those beyond an explicit bound that
depends on the cover and on the power of $p$ dividing $n$. At deep level
every double coset of $B_k$ carries a genuine Hecke function.
\item[Filtration subgroup] For $k\geq 1$, the subgroup $B_k$ of the
Iwahori subgroup consisting of matrices congruent to the identity modulo
$\varpi^k$ on and above the diagonal and modulo $\varpi^{k+1}$ below it.
These subgroups are normal in the Iwahori subgroup, they shrink to the
identity as $k$ grows, and their successive quotients are isomorphic to
$\Fq^3$. As $k$ grows, representations with $B_k$-fixed vectors include
representations of greater depth.
\item[Generalized Hecke algebra] For an open compact subgroup $K_0$ of
a locally profinite unimodular group $H$, the algebra of compactly
supported functions on $H$ that are invariant on both sides under $K_0$,
with convolution as product. The characteristic functions of the double
cosets of $K_0$ form a basis, and the characteristic function of $K_0$
is the identity when $K_0$ has volume one. Its simple modules correspond
to the irreducible smooth representations of $H$ with non-zero
$K_0$-fixed vectors.
\item[Genuine function] A function on the cover that transforms under
the central subgroup $\mu_n$ by the fixed embedding of $\mu_n$ into the
complex numbers. It is determined by its values on one lift of each
element of $G$, and it does not descend to a function on $G$ unless it
is zero. The genuine Hecke algebra consists of the compactly supported
genuine functions that are invariant on both sides under a lifted open
compact subgroup.
\item[Genuine representation] A smooth representation of the cover on
which the central subgroup $\mu_n$ acts by a fixed faithful character.
Since the character is faithful, such a representation does not factor
through $G$ or through any smaller cover. Genuine representations are
the local counterparts of modular forms of fractional weight, and
genuine functions act on them by integration.
\item[Hilbert symbol] The pairing $F^\times\times F^\times\to\mu_n$ of
local class field theory, for a $p$-adic field $F$ containing $\mu_n$. It
is bimultiplicative and skew-symmetric, it is trivial on the pairs
$(a,-a)$ and $(a,1-a)$, and it is non-degenerate modulo $n$-th powers.
Kubota's cocycle is built from it, and the non-triviality of
$(u,\varpi)_n^2$ for $n\geq 3$ is what makes the inverse image of the
diagonal torus non-abelian.
\item[Iwahori subgroup] The open compact subgroup of $SL(2,\cO)$ of
matrices whose lower left entry lies in $\varpi\cO$, so that they are
upper triangular modulo $\varpi$. It is the stabilizer of an edge of the
Bruhat--Tits tree of $SL(2,F)$, and it plays the role of a Borel
subgroup at the integral level. When $p\nmid n$ the metaplectic cover
splits over it.
\item[Iwahori--Hecke algebra] The Hecke algebra of an Iwahori
subgroup. For $SL(2)$ and $PGL(2)$ it is generated by two elements, each
satisfying a quadratic relation, with no further relation, because the
affine Weyl group is infinite dihedral. The two presentations differ
only in the eigenvalues of the second generator: $q$ and $-1$ for
$SL(2)$, and $\pm1$ for $PGL(2)$.
\item[Iwahori--Bruhat decomposition] The decomposition of $SL(2,F)$ into
the disjoint double cosets $BwB$ of the Iwahori subgroup $B$, where $w$
runs over representatives of the affine Weyl group. It shows that every
element of $G$ can be written as $b_1wb_2$ with $b_1,b_2\in B$. This is
how the Hecke algebras of the filtration subgroups are generated by
functions attached to $B$ and to the two reflections.
\item[Kubota's cocycle] The explicit $2$-cocycle on $SL(2,F)$ with
values in $\mu_n$, built from the Hilbert symbol and from the lower left
entry of a matrix, or the lower right entry when the lower left entry is
zero. It defines the multiplication of the metaplectic cover on the set
of pairs $(g,\zeta)$. It is a coboundary on sufficiently small open
subgroups, so the cover splits locally.
\item[Length] For a representative $w$ of the affine Weyl group, the
least number of reflections $s_1$, $s_2$ needed to write $w$ up to sign.
The translation $t_0^m$ has length $2|m|$, and $t_0^ms_1$ has length
$|2m+1|$. The volume of the double coset $B_kwB_k$ is $q$ to the power
of the length, so characteristic functions of such double cosets
multiply exactly when lengths add.
\item[Lift] An element of the cover that projects to a given element of
$G$. Two lifts of the same element differ by a unique element of the
central subgroup $\mu_n$. Working with lifts rather than with cocycle
coordinates makes the genuine Hecke basis independent of the choice of
cocycle within its cohomology class.
\item[Metaplectic cover] The central extension of $SL(2,F)$ by $\mu_n$
defined by Kubota's cocycle. As a set it is the product of $SL(2,F)$ and
$\mu_n$, but the multiplication is twisted by the cocycle. It splits over
sufficiently small open subgroups but not over the whole group, and for
$n\geq 3$ the inverse image of the diagonal torus is non-abelian.
\item[Orbital integral] The integral of a function over a conjugacy
class, taken with respect to an invariant measure on the class. For a
genuine function on the cover it is defined only at elements whose
lift commutes with the lifts of its centralizer. Matching of orbital
integrals between a cover and a linear group, up to transfer factors,
is the geometric form of a correspondence, and the fundamental lemma
asserts it for spherical functions. For the Iwahori--Hecke algebras of
this paper the orbital integrals at split elements are computed in
Sections~\ref{subsec:orbital} and~\ref{subsec:elliptic}.
\item[Principal congruence subgroup] For $j\geq 1$, the normal subgroup
$\Gamma_j$ of $SL(2,\cO)$ of matrices congruent to the identity modulo
$\varpi^j$ in every entry. It factors as a product of lower unipotent,
diagonal and upper unipotent pieces, and it contains $B_j$ as a normal
subgroup of index $q$. The square of the Hecke function attached to
$s_1$ is supported on $-\Gamma_k$.
\item[Savin's isomorphism] An isomorphism between the genuine
Iwahori--Hecke algebra of an $n$-fold cover and the Iwahori--Hecke
algebra of a linear group, which for $SL(2)$ is $SL(2)$ itself when $n$
is odd and $PGL(2)$ when $n$ is even. It keeps the finite part, sends
the genuine translations, which form a sublattice, to all the
translations of the linear group, and multiplies each basis function by
the square root of the ratio of the volumes of the two double cosets.
Savin calls the resulting bijection of representations a local Shimura
correspondence.
\item[Support-preserving map] A linear map between Hecke algebras of
$G$ and of $\Gbar$ that sends a function supported on a double coset to
a function supported over the same double coset
(Definition~\ref{def:supportpreserving}). At deep level these are the
maps $f_g\mapsto c_g\fbar_{\tilde g}$.
\item[Splitting] A homomorphism from a subgroup $H$ of $G$ into the
cover that is a right inverse of the projection. It exists exactly when
Kubota's cocycle restricted to $H$ is a coboundary, and two splittings
differ by a homomorphism from $H$ to $\mu_n$. The unipotent subgroups
split uniquely, and a genuine Hecke algebra requires a choice of
splitting of the open compact subgroup.
\item[Type] In the sense of Bushnell and Kutzko, a pair consisting of an
open compact subgroup and an irreducible representation of it whose
Hecke algebra describes one block, or finitely many blocks, of the
category of smooth representations. For the Iwahori subgroup of a split
linear group, the trivial representation is a type for the unramified
principal series block. In that case the representation induced from a
simple module of the Hecke algebra is irreducible.
\end{description}

%======================================================================
\section*{Acknowledgments}

D.~J.\ thanks Bill Wardlaw for helpful conversations on the proof of
the relations of Proposition~\ref{prop:main} in its original form, and
the anonymous referee of \cite{Jo98} for many helpful corrections and
suggestions.

AI systems contributed substantially to this version. Gemini (Google),
referred to as ``Carbon'', acted as referee of \cite{Jo98}, identified
that its main theorem could not hold as stated, refereed Draft~2 of this
revision, provided a summary of \cite{Sa88}, which we used only for
orientation, and commented on the relevance of \cite{Wi22}. We checked
those comments against \cite{Wi22} before using them. Claude
(Anthropic) is a co-author. Claude checked the referee report point by
point, drafted the corrected statements and the new proofs, revised the
paper to the authors' house style, checked bibliographic data, and
prepared the revision log. Two Claude models were used, Claude Fable
5.2 for Draft~1 and Claude Opus 5.5 for Drafts~2 to~11. AI review is not an
independent certification of the results, and the new arguments have
not yet been checked by an independent human reader. Further details on
the AI usage can be found in the revision log.

%======================================================================
\begin{thebibliography}{HM85}

\bibitem[Bo76]{Bo76} A.~Borel, Admissible representations of a
semi-simple group over a local field with vectors fixed under an
Iwahori subgroup, \emph{Invent. Math.} \textbf{35} (1976), 233--259.

\bibitem[BH06]{BH06} C.~J.~Bushnell and G.~Henniart, \emph{The local
Langlands conjecture for $GL(2)$}, Grundlehren der mathematischen
Wissenschaften 335, Springer, Berlin, 2006.

\bibitem[BK98]{BK98} C.~J.~Bushnell and P.~C.~Kutzko, Smooth
representations of reductive $p$-adic groups: structure theory via
types, \emph{Proc. London Math. Soc. (3)} \textbf{77} (1998), 582--634.

\bibitem[Fl80]{Fl80} Y.~Flicker, Automorphic forms on covering groups
of $GL(2)$, \emph{Invent. Math.} \textbf{57} (1980), 119--182.

\bibitem[FK86]{FK86} Y.~Flicker and D.~Kazhdan, Metaplectic
correspondence, \emph{Publ. Math. IH\'ES} \textbf{64} (1986), 53--110.

\bibitem[Ge76]{Ge76} S.~Gelbart, \emph{Weil's representation and the
spectrum of the metaplectic group}, Lecture Notes in Mathematics 530,
Springer, Berlin, 1976.

\bibitem[GS12]{GS12} W.~T.~Gan and G.~Savin, Representations of
metaplectic groups II: Hecke algebra correspondences, \emph{Represent.
Theory} \textbf{16} (2012), 513--539.

\bibitem[HM85]{HM85} R.~Howe (with the collaboration of A.~Moy),
\emph{Harish-Chandra homomorphisms for $p$-adic groups}, CBMS Regional
Conference Series in Mathematics 59, Amer. Math. Soc., Providence,
1985.

\bibitem[IM65]{IM65} N.~Iwahori and H.~Matsumoto, On some Bruhat
decomposition and the structure of the Hecke rings of $p$-adic
Chevalley groups, \emph{Publ. Math. IH\'ES} \textbf{25} (1965),
5--48.

\bibitem[Jo98]{Jo98} D.~Joyner, A correspondence for the generalized
Hecke algebra of the metaplectic cover $\overline{SL(2,F)}$, $F$
$p$-adic, \emph{New York J. Math.} \textbf{4} (1998), 223--235.

\bibitem[Ka21]{Ka21} E.~Karasiewicz, A Hecke algebra on the double
cover of a Chevalley group over $\mathbb{Q}_2$, \emph{Algebra Number
Theory} \textbf{15} (2021), 1729--1753.

\bibitem[KP84]{KP84} D.~A.~Kazhdan and S.~J.~Patterson, Metaplectic
forms, \emph{Publ. Math. IH\'ES} \textbf{59} (1984), 35--142.

\bibitem[Ku67]{Ku67} T.~Kubota, Topological covering of $SL(2)$ over a
local field, \emph{J. Math. Soc. Japan} \textbf{19} (1967), 114--121.

\bibitem[Li11]{Li11} W.-W.~Li, Transfert d'int\'egrales orbitales pour le
groupe m\'etaplectique, \emph{Compos. Math.} \textbf{147} (2011),
524--590.

\bibitem[LS10]{LS10} H.~Y.~Loke and G.~Savin, Representations of the
two-fold central extension of $SL_2(\mathbb{Q}_2)$, \emph{Pacific J.
Math.} \textbf{247} (2010), 435--454.

\bibitem[Lu18]{Lu18} C.~Luo, Spherical fundamental lemma for metaplectic
groups, \emph{Canad. J. Math.} \textbf{70} (2018), 898--924.

\bibitem[Ma71]{Ma71} I.~G.~Macdonald, \emph{Spherical functions on a group
of $p$-adic type}, Ramanujan Institute, University of Madras, 1971.

\bibitem[Pe14]{Pe14} L.~Peskin, A classification of the irreducible
admissible genuine mod $p$ representations of the metaplectic cover of
$p$-adic $SL(2)$, preprint, 2014, as cited in \cite{Wi22}.

\bibitem[RS77]{RS77} S.~Rallis and G.~Schiffmann, Repr\'esentations
supercuspidales du groupe m\'etaplectique, \emph{J. Math. Kyoto Univ.}
\textbf{17} (1977), 567--603.

\bibitem[Sa88]{Sa88} G.~Savin, Local Shimura correspondence,
\emph{Math. Ann.} \textbf{280} (1988), 185--190.

\bibitem[Sa04]{Sa04} G.~Savin, On unramified representations of
covering groups, \emph{J. Reine Angew. Math.} \textbf{566} (2004),
111--134.

\bibitem[Sc98]{Sc98} J.~Schultz, \emph{Lifting of characters of
$\overline{SL_2(F)}$ and $SO_{2,1}(F)$, for $F$ nonarchimedean}, Ph.D.
thesis (supervised by J.~Adams), University of Maryland, College Park,
1998.

\bibitem[TW18]{TW18} S.~Takeda and A.~Wood, Hecke algebra
correspondences for the metaplectic group, \emph{Trans. Amer. Math.
Soc.} \textbf{370} (2018), 1101--1121.

\bibitem[Wa80]{Wa80} J.-L.~Waldspurger, Correspondance de Shimura,
\emph{J. Math. Pures Appl.} \textbf{59} (1980), 1--133.

\bibitem[Wa24]{Wa24} R.~Wang, Hecke algebras for tame genuine principal
series and local Shimura correspondence, preprint, arXiv:2405.06254,
2024.

\bibitem[We18]{We18} M.~H.~Weissman, L-groups and parameters for
covering groups, \emph{Ast\'erisque} \textbf{398} (2018), 33--186.

\bibitem[Wi22]{Wi22} R.~Witthaus, The mod-$p$ representation theory of
the metaplectic cover of $GL_2(\mathbb{Q}_p)$, arXiv preprint
2208.12478, 2022.

\end{thebibliography}

\end{document}
